\quellseite{175}
er hätte die Faust auf dem Kopfe, so ginge es besser zu; sonst wird es kein nütze. Wenn sie Weiber nehmen, sind es ungezwogene Leute, wilde und wüste, deren niemand gebrauchen, noch mit ihnen umgehen kann.

32. Aber es ist vergebens, daß wir es sagen; allein, daß wir wissen, daß dies sind fromme, heilige Leute gewesen, haben kein Regiment gehabt, auch unter den Heiden; jetzt ist es gar nichts. Ein Knecht galt dazumal einen Gulden oder achte, eine Magd einen Gulden oder sechse, und mußte thun, was die Frau mit ihr machte. Und sollte die Welt lange stehen, könnte man es nicht wohl wieder halten im Schwange, man müßte es wieder aufrichten. Die Juden hatten einen Vortheil, daß sie sich sechs Jahre verkauften, im siebenten gingen sie frei davon; wie alles auf das feinste geordnet war von Gott.

33. Zuletzt, daß Abimelech sagt: „Siehe, der soll dir ein Deckel deiner Augen sein, und ein Verantworter“, ist ein finsterer Text, [ich] weiß nicht, ob wir es treffen werden. Das ist, dein Mann und Bruder, oder das\footnote{„das“ fehlt in der Erlanger.} Geld, das ich dir da schenke (denn es ist nicht gewiß, ob es von Abraham, oder vom Geschenke gesagt ist), soll dir ein Deckel sein, damit du die Augen bedeckt, und dich beschönet. Daß dies die Meinung sei: Ich habe dich genommen, und gemeint, du wärest des Mannes Schwester, nicht sein Eheweib, und bist nun in meinem Hause gewesen; auf daß du nun ein Zeugniß mitbringest, daß du deiner Ehren unverrückt bist, so hast du das Geschenk, und den Mann dabei, daß er es sieht und zuhört, auf daß er dadurch deine Ehre verantworte. So will der König, daß er es gethan habe, daß er die Sünde redlich büße, doch unwissend gethan; und will damit aufheben, daß er ihm nichts könne aufrücken, und sie die Augen zudecke, das ist, sich schmücke und beschöne vor jedermann, daß ihre Ehre ganz geblieben ist. Ist das nicht die Meinung, so weiß ich nichts anders; doch liegt hier nicht große Macht an.

34. So hielt nun Gott die seltsame Strafe und Plage der Weiber Abimelechs, welcher zuvor geschlagen war, wie oben gehört, mit einem großen, greulichen Schrecken, daß er den Tod und Hölle fühlte. Sein Weib und Mägde, oder eigene Leute mit ihren Kindern, die auch des Herrn waren, die hatte Gott alle geplagt, daß sie nicht konnten der Frucht los werden; also, daß Frucht und Zeit, da sie gelegen sollten (aber nicht konnten), in Fahr standen, und den Tod vor Augen hatten.

\section{Das einundzwanzigste Capitel}

V. 1—5. Und der Herr suchte heim Sara, wie er geredet hatte, und that mit ihr, wie er geredet hatte. Und Sara ward schwanger, und gebar Abraham einen Sohn in seinem Alter, um die Zeit, die ihm Gott geredet hatte. Und Abraham hieß seinen Sohn, der ihm geboren war, Isaak, den ihm Sara gebar. Und beschnitt ihn am achten Tage, wie ihm Gott geboten hatte. Hundert Jahr war Abraham alt, da ihm sein Sohn Isaak geboren ward.

1. Moses steht lange auf dem guten Vater Abraham, wird noch viel von ihm sagen. Ist eine Historia, die Gott lieb hat, und ihm sonderlich gefallen läßt, um des Glaubens willen.

2. In diesem Capitel sind zwei Stücke; das erste, wie Abraham ein Sohn geboren ist, und wie es mit dem Stiefbruder ist gegangen. Das andere, was er abermal schafft mit dem Könige Abimelech. St. Paulus hat [Gal. 4, 23. Hebr. 11, 11.] dies Capitel auch gerühret und angezogen, daß auch viel daran gelegen ist. Es sieht, als sei es eine neidische, häßige und feindselige Historia, die unfreundlich lautet; [wir] wollen aber hören, was sie in sich hat.

3. Mit vielen Worten hebt er an und sagt, wie Gott die Sara habe heimgesucht, daß sie schwanger worden ist, und einen Sohn geboren hat etc. Diese Worte alle wollen das, daß der Sohn, wie droben [Cap. 16, § 9] gesagt, wiewohl er Fleisch und Blut war von Abraham und Sara, doch nicht durch Menschenkräfte, noch
\quellseite{176}
nach natürlichem Laufe geboren ward. Denn sie waren beide zu alt, dazu war sie unfruchtbar, daß es gar schwer, ja, unmöglich war, daß hieraus eine Frucht sollte kommen. Darum ist es das Wunder, das Gott mit ihm thut, daß er ihm den Sohn gibt von seiner rechten Hausfrau. Den Sohn hat er auch lieber denn jenen von Hagar; wiewohl das Vaterherz auch da war; aber hier noch viel mehr. Das ist es, das Mosis hier erstlich will anzeigen.

4. Wir haben aber [Cap. 17, § 19 ff.] gehört, wie Gott geboten hat die Beschneidung, und wie die Werktheiligen sich stoßen an die Schrift, meinen, es sei ein schlechter Tand, sei anders nichts, denn daß die Alten haben Weiber genommen, und Kinder gezeugt, wie der gemeine Lauf ist; so spotten sie und lachen Gottes. [Es] muß also gehen, daß man ihn für einen Narren hält. Er geht aber nichtsdestoweniger immer fort; hätte hie auch wohl ander Wunder können thun, daß für groß angesehen wäre, da die ganze Welt das Maul aufsperrete, und davon sänge und sagte, so nimmt er, daß die ganze Welt lacht und spottet, und sie, die Sara selbst, lachen muß, und schimpft,\footnote{schimpft = scherzt.} daß man soll sagen, wie das alte Weib habe beim Mann geschlafen. Wie das auch war, daß er den alten Mann heißt beschneiden; er hätte sich doch mögen schämen seines Alters halben; warum thut er ihm die Schande auf? Er hat Lust dazu, daß er solch verächtlich, spöttisch Ding anrichtet, läßt die Welt wohl hohnlachen und spotten, sieht ihr zu, wie lange sie lachen kann. Darum sollen wir nicht gering achten die Hand Gottes, die darinne ist.

5. Es ist ein theuer Ding um das\footnote{„das“ fehlt in der Erlanger.} Leben. Wo kommt es aber her? Von Fleisch und Blute, dem schlammichten und lauteren Dreckjache, daß man sich es schämt zu sagen. Das sind Gottes Werke und edle Schätze, die bringt er heraus, da die Welt ihre Lächerei austreibt. Auch hat Gott zwar allzumal gesetzt Mann und Weib, daß er der Welt die Augen zuthue, daß sie selbst sieht, wie sie es nicht umgehen kann; es muß zuhause, und Kinder haben. Was hat es sie gekostet, die Eltern, daß sie uns das Leben geben? Warum sollten wir es denn verachten, und nicht wieder so mit Kindern umgehen, wie sie mit uns?

Warum verlacht man es denn, weil es Gott so angegriffen hat, sonderlich im alten Testamente, da er es so hoch hebt und ehrt; will es nicht belacht haben, sondern Ernst lassen sein, und ein sonderlich Gebot darauf gelegt, daß man soll Vater und Mutter ehren [2 Mos. 20, 12.]. Aber die Welt muß also narren, daß sie Gottes Wunder mit Frauen und Kindern lacht. Darum spricht sie nun:

B. 6. Gott hat mir ein Gelächter zugerichtet, denn wer es hören wird, der wird mein lachen.

6. Die Schrift zeucht es so an, daß sie züchtig und fromm ist gewesen, wie sie droben [Cap. 18, 12.] gesagt hat: „Soll ich erst mit Wohllust zu schaffen haben, weil ich nun von neunzig Jahren bin?“ Das ist das Gelächter. So hat auch Lucas [Cap. 1, 36. 57.] von Elisabeth, Johannes des Täufers Mutter, geschrieben; die ging auch schwanger daran, daß sie die Schande sollte freien. So züchtig und fromm sind die heiligen Weiber auch vor der Welt gewesen, daß sie sich auch des Gerüchtes schämten, da eitel Ehre und Gottes Werk war.

7. So muß Gottes Werk allenthalben lächerlich und spöttisch gehen, daß es scheint, als zeucht er eine Narrenkappe an, läßt es wohl lachen, wird aber uns die Schellen wieder aufknüpfen; darum will er haben, daß ein jeglicher die Augen zuthue, sehe nur, wo das Wort ist. Da hält er groß von; ob es gleich jedermann recht schändlich dünkt, halte er stille, es muß doch zuletzt zu allen Ehren werden.

8. So muß man auch nicht sehen, daß Sara, das alte Weib, beim Manne schläft, sondern dahin sehen, daß Gott sagt: Sara muß zu Schanden werden; aber ich will sie zu Ehren machen, daß sie gepreiset und gepredigt soll werden vor allen Weibern. Wie viel Königinnen und große Frauen sind, die zu der Ehre nicht gekommen sind, daß also von ihrem Leibe und Frucht gepredigt würde, und hervor gezogen in aller Welt? Es ist eine ewige Ehre, und ein Exempel, damit man alle Welt regieren kann. Also ehrt Gott die Seinen mit seinen Werken, daß sie eine Zeitlang in Schmach liegen, die wird darnach eine ewige Ehre, die niemand genug preisen kann. Wenn jetzt die allermächtigsten und reichsten Königinnen alle zusammen thäten, würden sie froh sein, daß sie sollten die Windeln waschen, da Isaak inne gelegen wäre.
\quellseite{177}
Solche Ehre kommt jetzt aus der vorigen Schande. Siehe, das sind Gottes Wunderzeichen. Wenn wir die Augen könnten zuthun, ließen die Welt lachen, und hielten an Gottes Worte, würden wir sehen, wie es die Gläubigen zu Ehren macht, die Welt aber zu allen Schanden mit ihrer Ehre. Weiter spricht sie:

V. 7. Wer durfte es auch Abraham selbst ansagen, daß Sara Kinder säugete, und hätte ihm einen Sohn geboren in seinem Alter?

9. Da wird sie fröhlich, und überwindet die Schande. Zuvor sagt sie: „Gott hat mir ein Gelächter zugerichtet.“ Als sollte sie sagen: Ei, wie thut es so wehe, daß man von mir sagen soll etc. Das ist noch eine weibische Krankheit, daß sie allezeit schwächer von Muth sind denn der Mann; darum hält ihr es auch\footnote{„auch“ fehlt in der Erlanger.} Gott zugute. Wiederum sieht sie an die große Gnade, und spricht: Es soll's Abraham selbst nicht glauben, daß ich ein Kind trüge. Das thut allemal Gottes Werk; wenn man stille hält, so folgt Lust und Freude hernach desto größer. Darum sagt der Text weiter:

V. 8. Und das Kind wuchs und ward entwöhnet; und Abraham machte ein groß Mahl am Tage, da Isaak entwöhnet ward.

10. Das ist seltsam Ding: Warum macht er nicht ein groß Mahl, da der Sohn geboren ward? oder da die sechs Wochen um waren (wiewohl sie da vielleicht nicht haben sechs Wochen gehalten), oder am achten Tage? Was das\footnote{So die Jenaer. Wittenberger: „dies“; Erlanger: es bedeutet.} bedeutet, kann ich nicht wissen. Ich halte, es sei zugegangen, wie die Weiber bei uns Gewohnheit haben, daß sie zuweilen zusammen gehen und Kinderbier halten,\footnote{Die Worte „und Kinderbier halten“ fehlen in der Erlanger; in der Jenaer und in der Wittenberger stehen sie im Text.} schlemmen und dennen. So scheint es, als sei es eine Landesitte und Weise gewesen, um die Zeit, wenn man das Kind entwöhnet hat.

11. Das will nun die Schrift anzeigen, daß er sich habe geschickt, wie sie sich hielten im Lande; weil er im Lande ist, muß er sich ihnen eben machen, und Landes Sitten halten, wo es nicht schädlich noch sündlich ist. Darum hat er sich billig und recht also gehalten, wie andere.\footnote{Wittenberger und Erlanger: wie ein anderer.}

12. [Es] ist ein Exempel der Liebe, daß man sich nach dem Nächsten richte, und eben mache, daß [es] die leiden können, bei denen wir wohnen. Sonst gefällt einem\footnote{„einem“ fehlt in der Erlanger.} jeglichen seine Weise, und anders nicht; [er] meint, es soll sich ein ganz Land nach seinem Kopf richten. Darum sollen wir uns halten nach Landes Sitten, und uns dieselben gefallen lassen, daß wir Frieden halten mit jedermann, wie Paulus Röm. 12, 18. lehrt, nicht eine eigene Weise aufwerfen und ausführen nach eigenem Kopf. Das hat er, achte ich, hier gemeint. Nun kommt das andere Stück, daß Sara den andern Sohn, Ismael, austreibt. [Es] spricht der Text also:

V. 9—21. Und Sara sahe den Sohn Hagar, der Egyptischen, den sie Abraham geboren hatte, daß er ein Spötter war, und sprach zu Abraham: Treibe diese Magd aus mit ihrem Sohne, denn dieser Magd Sohn soll nicht erben mit meinem Sohn Isaak. Das Wort gefiel Abraham übel, um seines Sohnes willen. Aber Gott sprach zu ihm: Laß dir’s nicht übel gefallen des Knaben und der Magd halben, alles, was Sara dir gesagt hat, dem gehorche. Denn in Isaak soll dir der Same genannt werden, auch will ich der Magd Sohn zum Volk machen, darum, daß er deines Samens ist. Da stund Abraham des Morgens früh auf, und nahm Brod und eine Flasche mit Wasser, und legte es Hagar auf ihre Schulter, und den Knaben mit, und ließ sie aus. Da zog sie hin, und ging in der Wüste irre bei Bersaba. Da nun das Wasser in der Flasche aus war, warf sie den Knaben unter einen Strauch, und ging hin, und setzte sich gegenüber von\footnote{„von“ fehlt in der Erlanger.} ferne, einen Armbrustschuß weit. Denn sie sprach: Ich kann nicht zusehen des Knaben Sterben. Und sie setzte sich gegenüber, und hub ihre Stimme auf, und weinete. Da erhörte Gott die Stimme des Knaben. Und der Engel Gottes rief vom Himmel der Hagar, und sprach zu ihr: Was ist dir, Hagar? Fürchte dich nicht, denn Gott hat erhört die Stimme des Knaben, da er liegt. Stehe auf, nimm den Knaben, und halte ihn mit deinen Händen, denn ich will ihn zum großen Volk machen. Und Gott that ihr die Augen auf, daß sie einen Wasserbrunnen sahe, da ging sie hin, und füllete die Flasche mit Wasser, und tränkte
\quellseite{178}
den Knaben. Und GOtt war mit dem Knaben, der wuchs, und wohnte in der Wüste, und ward ein Schützenmeister, und wohnete in der Wüste Pharan. Und seine Mutter nahm ihm ein Weib aus Egyptenland.

13. Ismael ist nun fast dreizehn Jahre alt gewesen, wie zu rechnen ist aus dem vorigen Capitel; ein Knabe, der zu seinen Jahren kommen war. So wird nun das lange hernach geschehen sein, daß die zween Söhne konnten mit einander reden. Denn er sagt, sie hatte Achtung auf ihn, und sahe, daß er ein Spötter war. Das verdroß sie, darum will sie ihn mit der Magd austreiben. Und da siehst du, wie sie eine züchtige, feine Hausmutter ist. Droben [Cap. 16, 5.] wollte sie die Magd auch züchtigen, da sie zu stolz war. Der Text drückt aber nicht aus, wie das zugegangen sei; es sind auch hier viele jüdische lose Fabeln erdichtet.

14. Aber St. Paulus rühret es zum Galatern [Cap. 4, 29. 30.] und legt es aus auf den Verstand: „Gleichwie zu der Zeit, der nach dem Fleisch geboren war, verfolgte den, der nach dem Geist geboren war; also gehet es jetzt auch. Aber wie spricht die Schrift: Stoß die Magd hinaus mit ihrem Sohne, denn der Magd Sohn soll nicht erben mit der Freien.“ Das deutet er selbst also: Wenn man predigt vom Glauben, daß die Werke nicht gelten, fromm zu machen, das ist eine geistliche Predigt; welche\footnote{In den Ausgaben: welche, die das.} das fassen und halten, das sind geistliche Leute, Abrahams Kinder. Die andern aber mögen das nicht leiden, so von Werken predigen; das ist, der große fleischliche Haufe spottet, verfolgt die andern. So bleibt ein Streit, weil die Welt steht, zwischen Glauben und Werken; wollen beide ihr Ding ungetadelt haben. Nun ist dieser viel, und sind stärker, die mit Werken umgehen; darum liegen sie zeitlich oben, und verfolgen jene. Doch kommt endlich das Urtheil, daß GOtt diese verwirft, und den andern das Erbe zutheilt.

15. So wird es nun zugegangen sein, wie wir noch sehen in den Werthheiligen. Die Mutter Sara wird das Kind gezogen haben, wie eine Mutter einen rechten Erben zieht. Der andere Sohn aber ist nun groß, verließ sich darauf, Abraham wäre sein Vater, und [er] wäre der erste Sohn, und gebührete ihm zweierlei Recht. So wird auch (dachte er) der Vater mich lieber haben, und hat immer darauf getroßt, daß er wollte obenan sein. Solches war nun nicht zu leiden, sie konnte es und sollte es auch nicht länger leiden, denn er hatte solches ohne Zweifel lange getrieben, und zu viel gemacht. Und was seine Mutter zuvor der Sara gethan hatte, will der Sohn jetzt auch thun. Zuvor wollte sie Frau sein, und achtete Sara für nichts. Nun, weil sie den Sohn hat, will er nach der Mutter geraten, will ihren Sohn auch unterdrücken. Darum denkt sie, sollte es so fortgehen, so müßte ich zuletzt mit meinem Sohn hinweg; darum wird nichts Besseres, denn ich theile sie ab, gebe ihm, was er haben soll, und lasse ihn laufen.

16. Aber das verdroß Abraham; es war auch Blut und Fleisch und väterliches Herz, das ihn rührte, und sahe seiner dazu, so lange, bis ihm GOtt sagte: „Laß dir es nicht übel gefallen, was dir Sara gesagt hat“; das Urtheil fälle ich, spricht er: Der Magd Sohn muß hinaus, denn „von Isaak soll dein Same genennet werden“. Da steht der Spruch, den Paulus hoch anzeucht. [Er] will so sagen: Du darfst nicht auf Ismael sehen, daß er der rechte Same sei, davon Christus kommen wird, und alle, die gen Himmel gehören, sondern Isaak ist es: was von ihm kommt, das soll der rechte Same sein. Also ist Ismael ausgeschlossen, und durch Gottes Wort Isaak eingesetzt, und in das Wort gefasset, daß Abraham sicher wird, daß er nicht allein ein Erbe von Fleisch und Blut ist, sondern von Gnaden. Doch sagt er: „Ich will der Magd Sohn auch zu einem großen Volk machen“, ohne daß er die Ehre nicht soll haben, daß Christus von ihm komme.

17. So haben wir nun die zween Söhne des Erzvaters Abraham, und wie sie deuten zweierlei Volk auf Erden, die rechtschaffen Gläubigen,\footnote{Erlanger: gläuben.} und Heuchler oder falsche Heilige: und wie Ismael mit der Mutter ausgestoßen wird, Isaak aber mit seiner bleibt und erbt. Daß das die Deutung sei, daß die Juden, so mit Werken umgingen, ausgeworfen werden, weil sie ohne Glauben sind; die andern aber, im Glauben, sollten Gottes Volk bleiben. Das soll und muß immer also gehen. Darum hält es uns die Schrift vor, daß sie uns des gewiß mache. Es ist nicht
\quellseite{179}
möglich, daß alle gleich Isaac seien; der größte und beste Haufe muß wider die andern sein. Folgt weiter:

V. 22—34. Zu derselbigen Zeit redeten Abimelech und Phichol, sein Feldhauptmann, mit Abraham, und sprach: Gott ist mit dir, in allem, das du thust. So schwöre mir nun bei Gott, daß du mich, noch meine Kinder, noch meine Kindeskinder, nicht \unsicher{vervortheilet}, sondern die Barmherzigkeit, die ich an dir gethan habe, an mir auch thuest, und an dem Lande, da du ein Fremdling innen bist. Da sprach Abraham: Ich will schwören. Und Abraham strafte Abimelech um des Wasserbrunnens willen, den Abimelechs Knechte hatten mit Gewalt genommen. Da antwortete Abimelech: Ich habe es nicht gewußt, wer das gethan hat, auch hast du mir’s nicht angesagt; dazu habe ich es nicht gehört, denn heute. Da nahm Abraham Schafe und Rinder, und gab sie Abimelech, und machten beide einen Bund mit einander. Und Abraham stellte dar sieben Lämmer besonders. Da sprach Abimelech zu Abraham: Was sollen die sieben Lämmer, die du besonders dargestellt hast? Er antwortete: Sieben Lämmer sollst du von meiner Hand nehmen, daß sie mir zum Zeugniß seien, daß ich diesen Brunnen gegraben habe. Daher heißt die Stätte Bersaba,\footnote{Randglosse: Bersaba heißt auf deutsch Schwörbrunn oder Eidsbrunn; möchte auch wohl Siebenbrunn heißen.} daß sie beide mit einander da geschworen haben, und also machten sie den Bund zu Bersaba. Da machten sich auf Abimelech und Phichol, sein Feldhauptmann, und zogen wieder in der Philister Land. Abraham aber pflanzte Bäume zu Bersaba, und rief daselbst an den Namen des Herrn, des ewigen Gottes. Und war ein Fremdling in der Philister Lande eine lange Zeit.

18. Zuletzt ist noch dies Stück da, wie Abraham mit Abimelech einen Bund macht des Wassers halben; ist auch darum geschrieben, daß immerdar angezeigt wird, wie die lieben Gottes Kinder auf Erden müssen viel böser Tücke leiden; denn da sitzt er im fremden Lande, da muß er ein Gast sein, wie der Prophet [Jer. 52, 4.] sagt, und sich nähren mit Kummer. So muß es uns auch gehen, wenn es recht soll gehen. Das sei genug von dem Capitel.

\section{Das zweiundzwanzigste Capitel}

V. 1. Nach diesen Geschichten versuchte Gott Abraham und sprach zu ihm: Abraham, Abraham! Und er antwortete: Hier bin ich.

1. Da haben wir abermal ein Stück von der Legende des heiligen alten Erzvaters. Ist ein recht fein Stück, wer es könnte mit Worten erlangen. Jetzt haben wir gehört, wie der gute Abraham bisher in mancherlei Anfechtung und Versuchung Gottes gestanden ist, und nie keine gewisse Stätte gehabt hat. Es ist schlecht beschrieben, aber so reichlich angezeigt, daß freilich nicht viel Legenden also geschrieben sind.

2. Wenn Gott einen Heiligen lobt, so lobt er ihn also, daß [es des] Lobens werth ist. So hat er diesem Manne gethan. [Er] ist ein recht Exempel des Glaubens. Sein Leben ist doch lauter Anfechtung, [es] muß alles gehen im Glauben; noch muß er Weib, Kind, groß Gesinde, dazu Geld und Gut haben, doch immer dahin stellen, daß es Gott hinnehme. [Er] hat sein Weib, Sara, oft müssen in die Fahr gehen, dazu hat ihm Gott die Magd, Hagar, mit dem Sohne hingenommen; das Land geht auch dahin, daß er nicht einen Fußbreit hat. Also spielt Gott mit ihm, wie mit einem Apfel; und er hält so stille, läßt ihn mit ihm machen und gehen, wohin es geht. Wenn er es hat, kann er es recht brauchen; wiederum kann er es entbehren, wenn es hinweg fährt.

3. Hier aber greift er ihn auf das höchste an, nimmt ihm das höchste Gut; denn er nichts Lieberes auf Erden hat; deß er auch Ursache hatte. Denn er war ein einiges Kind, und hatte von ihm Gottes Wort und Verheißung, daß von ihm sollte kommen der Same, dadurch alle Völker sollten gesegnet werden. Da läßt Gott die natürliche Begierde stecken, denn er hat ihm das Vaterherz nicht genommen, bricht die Natur
\quellseite{180}
nicht, läßt sie so bleiben, wie er sie geschaffen hat, hebt sie nicht auf, auch in den Allerheiligsten; wie du siehst, daß Abraham so ein großer Heiliger ist, noch läßt er den natürlichen Affect und Bewegung tief in ihm stecken.

4. Es ist nicht, wie etliche meinen, daß man so gar verachten soll alle Dinge, daß wir keine Schmerzen noch Leid sollen haben oder fühlen. Ich soll mich meines Nächsten Schaden herzlich annehmen und lassen leid sein, als meines eigenen; doch so geschickt sein, daß ich nicht darnach frage, wenn mir alles genommen wird, bis auf einen Heller. Die rechte Liebe nimmt sich des ganzen Menschen und aller seiner Güter an. Aber so thut sie: Das Gröbste behält sie zum Gröbsten; so hilft sie zum zeitlichen Leben auch. Also achtet ein Christ seinethalben nichts, sondern, was ihm Gott in die Hand gibt, nimmt er an, geht damit um; wenn er es wieder wegnimmt, läßt er es fahren.

5. Also nimmt Gott dem frommen Vater das allerhöchste Kleinod; nicht Holz noch Steine, weder Silber noch Gold, sondern den einigen Sohn, den er sonderlich darum lieb hat, über die gemeine Vaterliebe, daß er von dem heiligen Weibe geboren war, damit es Gott so wunderlich schaffte in den alten Tagen, dazu auch die Verheißung von dem herrlichen Samen auf ihn lautet, daß er große und viel Ursache hat, ihn zu lieben. Gott hat ihn auch selbst lieb; noch, als er sieht, daß er ihn so lieb hat, und recht daran thut, denkt er: Harre, ich muß mit ihm spielen, und spricht:

B. 2—8. Nimm Isaak, deinen einigen Sohn, den du lieb hast, und gehe hin in das Land Morija,\footnote{Randglosse: Morija heißt eine Erzeigung, und ist der Berg, da Salomon hernach zu Jerusalem den Tempel auf bauete, und heißt der Berg der Erzeigung, daß Abraham daselbst der Widder erzeigt ward, und darnach der rechte Widder, Christus, erzeigt und offenbar sollte werden von Gott dem Herrn.} und opfere ihn daselbst zum Brandopfer auf einem Berge, den ich dir sagen werde. Da stund Abraham des Morgens frühe auf, und sattelte seinen Esel, und nahm mit sich zween Knaben, und seinen Sohn Isaak, und spaltete Holz zum Brandopfer, machte sich auf, und ging hin an den Ort, davon ihm Gott gesagt hatte. Am dritten Tage hub Abraham seine Augen auf, und sahe die Stätte von ferne, und sprach zu seinen Knaben: Bleibet ihr hier mit dem Esel; ich und der Knabe wollen dorthin gehen, und wenn wir angebetet haben, wollen wir wieder zu euch kommen. Und Abraham nahm das Holz zum Brandopfer, und legte es auf seinen Sohn Isaak. Er aber nahm das Feuer und Messer in seine Hand; und gingen die beide mit einander.

6. Er greift ihn eben an mit rechten Worten, davon ihm das Lachen wohl vergangen mußte. Da mußte die Natur zurück prallen, und ihr wehe thun, und, das viel härter ist, daß er es selbst thun soll, und nicht ein anderer, sondern selbst das Schwert zücken, und schlachten zu einem Brandopfer; also, daß alles gar zu Pulver sollte brennen, daß nicht ein Härlein bliebe.

7. Nun, das sind je große, mächtige Versuchungen, daß Gott den Vater so bringt, seinen einigen, lieben Sohn zu würgen, und wiederum [zu] nehmen, den er ihm zuvor gegeben hatte, daß er freilich lieber zweimal selbst gestorben wäre, und alle sein Gut, Weib, und was er hatte, dahin gegeben. Denn er hat also müssen denken: Nun ist der Sohn schon dahin; hat nicht in Sinn genommen, daß er sollte lebendig bleiben. Darum muß es ihm ohne Zweifel über die Maße wehe im Herzen gethan haben.

8. Was thut aber der Herr? Er heißt ihn auf einen Berg gehen, den er ihm weisen will. Das waren drei ganze Tagereien. Es wäre ja noch leichter, wenn es ja sein sollte, daß er bald davon käme; aber er martert ihn noch weiter, nimmt noch Zeit dazu, daß er wohl gebraten wird, und sich durchfressen muß. Daß freilich ein Stück oder zehn sind gewesen, die ihm das Herz pochend gemacht haben, daß er hätte mögen sagen: Wie stellt sich doch Gott so närrisch? womit habe ich es verdient, daß er mich so martert? Solcher Gedanken hat er ohne Zweifel unzählig viel gehabt. Ueber das, da sie an den Berg kommen, muß Isaak sein eigen Kreuz tragen. Da gibt ihm der Sohn auch einen Stoß. Es hat ihm auch wehe gethan, daß er es allein mußte, darf dem Sohne nichts davon sagen, weder der Mutter, noch niemand. So spricht nun Isaak zum Vater:

B. 9—14. Mein Vater. Abraham antwortete: Hier bin ich, mein Sohn. Und er sprach: Siehe, hier ist Feuer und Holz, wo ist aber das Schaf zum Brandopfer? Abraham antwortete: Gott
\quellseite{181}
wird mir zeigen, mein Sohn, das Schaf zum Brandopfer. Und gingen die beide mit einander. Und als sie kamen an die Stätte, die ihm Gott sagte, bauete Abraham daselbst einen Altar, und legte das Holz darauf, und band seinen Sohn Isaak, legte ihn auf das Holz, und reckte seine Hand aus, und fassete das Messer, daß er seinen Sohn schlachtete. Da rief ihm der Engel des HErrn vom Himmel, und sprach: Abraham, Abraham! Er antwortete: Hier bin ich. Er sprach: Lege deine Hand nicht an den Knaben, und thue ihm nichts. Denn nun weiß ich, daß du Gott fürchtest, und hast deinen einigen Sohn nicht behalten vor mir. Da hub Abraham seine Augen auf, und sahe einen Widder hinter ihm, in der Hecke mit seinen Hörnern verfangen, und ging hin, und nahm den Widder, und opferte ihn zum Brandopfer, an seines Sohnes Statt. Und Abraham hieß die Stätte: Der HErr siehet.\footnote{Randglosse: Siehet, das ist, Gott sorgt für alles und wacht.} Daher man noch heutiges Tages sagt: Auf dem Berge, da der HErr siehet.

9. Das Wort wird ihm ein glühender Spieß im Herzen gewesen sein; als sollte er sagen: Ach! lieber Sohn, schweig stille, es möchte mir das Herz brechen. Das muß er auch noch fressen, und das Herz wohl kochen. Darnach fährt er fort, und schreibt, wie es gegangen sei, einfältig dahin, gibt aber einem jeglichen zu bedenken, wie das Vaterherz gestanden ist. Isaak ist dem Vater gehorsam, gibt sich darein, ist gleich so viel, als wäre er bereits todt. Wunderlich groß ist es, daß er das Herz so konnte überwinden.

10. Aber das Größte ist, daß Gott hier wider sich selbst redet. Das ist ein Räthsel, das niemand errathen kann denn der Heilige Geist. Denn Gott hat geboten, man soll nicht tödten [2 Mos. 20, 13.]. Nun heißt er es hier selbst, so doch Isaak nichts verschuldet hatte. Item, [er] hat ihm zuvor verheißen, von Isaak den Samen zu geben. Das Wort mußte Abraham glauben; also, daß sein Herz so stand: Der Sohn muß ein Vater vieler Kinder werden, und soll ausgebreitet werden in alle Welt. Wie stellt sich denn Gott also? Nun wird er wetterwendisch, und redet das Widerspiel, und muß der Sohn jetzt sterben; was will doch hieraus werden? Was kann die Vernunft da sagen? Sie ist ganz geschlagen, daß sie nicht weiß, wo hinaus, und muß sagen, es sei nun aus.

11. Nun mußte dennoch dies beides wahr sein: Abraham glaubte und wußte nichts anders, denn der Sohn müßte sterben; wiederum auch, daß er sollte ein Vater werden vieler Völker [Röm. 4, 17.]. Wie reimen sich die zwei zusammen? Also, wie es St. Paulus [Hebr. 11, 17.] auslegt, er hat müssen also denken: Gott ist allmächtig und wahrhaftig, der Sohn ist schon hin, den muß ich lassen fahren; aber Gott hat noch so viel, wenn ich und alle Welt todt wären, kann er ihn wieder auferwecken, auch über hundert Jahr, und zum Vater machen. Also läßt ihm Gott nicht mehr denn den einigen Trost, daran er sich in der höchsten Versuchung hält, daß ihn Gott wieder lebendig machen würde, wenn ihm es gefiele, darum, daß Gott nicht lügen kann. Das ist nun der allerhöchste Streit, wenn sich Gott also stellt, als lüge er, daß sein Wort wider einander lautet.

12. Solches lässet uns nun wohl lernen, ob wir auch also versucht würden, sonderlich wenn wir nun sterben sollen. Weil uns Gott das Evangelium gibt und tröstet, so ist es alles lieblich. Wie aber, wenn er spricht am Tode: Ich will dein nicht? Hast du nun fast gemeint, du stehest fest und glaubest an Christum, daß er dein ist, so kommt dir nun ein Spieß in das Herz, daß Gott sagt: Ich will dir Christum nicht geben, darum gib mir ihn her, und bleib du allein. Was kann da das Herz anders sprechen, denn: Es ist verloren? Und wenn es schon sagt: Hast du mir ihn doch geredet, so sagt er wohl wieder: Ich bin Gott, mag ich es nicht machen, wie ich will? Was fragst du darnach, bist du doch mein? Wes soll ich mich nun halten, wenn es so wider einander dringt? Den Spruch sollst du haben in 81. Psalm, V. 9. 10.: Israel, si audieris me, non erit in te Deus alienus, „höre, mein Volk, ich will dir zeugen, Israel, du sollst mich hören, daß unter dir kein fremder Gott sei, und du keinen unbekannten Gott anbetest“ etc. Darum, wenn du ihm gehorchst, so darfst du nicht denken, daß du einen Gott habest, der sich ändert. Semel loquitur Deus etc., „wenn Gott einmal redet (sagt Hiob [Cap. 33, 14.]), so wiederholet er es nicht.“ Item, Malachias [Cap. 3, 6.] spricht: Ego Deus et non mutor.

13. Darum hat Gott die Weise, daß er will über dem ersten Worte halten, das er geredet hat.
\quellseite{182}
Was nun hernach dawider lautet, will er, daß wirs uns nicht annehmen, sondern gewiß seien, es sei GOttes Versuchung. Das ist schon geboten im fünften Buche Mosis [Cap. 13, 1—3.]: „Wenn ein Prophet oder Träumer unter euch wird aufstehen, und gibt dir ein Zeichen oder Wunder, und das Zeichen oder Wunder kommt, davon er dir gesagt hat, und spricht: Lasset uns andern Göttern nachwandeln, die ihr nicht kennet, und ihnen dienen; so sollst du nicht gehorchen den Worten solches Propheten oder Träumers; denn der HErr, euer GOtt, versuchet euch, daß er erfahre, ob ihr ihn von ganzem Herzen und von ganzer Seele lieb habet.“ Das sind treffliche Worte: Wenn ein Prophet kommt, und dir ein Zeichen gibt, das muß je GOtt thun: noch mußt du es nicht glauben; und gibt Ursach, denn GOtt versucht euch, wenn er das Wort will wiederholen. Darum halte fest bei dem ersten Worte, und laß das andere alles gehen, allein das lasse nicht fahren. So strenge will das GOtt haben, daß er auch alle Zeichen, die dawider\footnote{„dawider“ fehlt in der Erlanger.} geschehen (wenn sie gleich von ihm kämen), nicht zu achten gebeut.

14. Solches hat auch St. Paulus vermahnt zum Thessalonichern [2. Ep. 2, 2.]: „Ich bitte euch, daß ihr euch nicht bald bewegen lasset von eurem Sinne, noch erschrecken, weder durch Geist, noch durch Wort, noch durch Briefe, als von uns gesandt, daß der Tag Christi vorhanden sei.“\footnote{Erlanger: ist.} Und bald darnach, V. 9., vom Antichrist: „Welches Zukunft geschieht nach der Wirkung des Teufels, mit allerlei lügenhaftigen Kräften, und Zeichen, und Wundern“; wie auch Christus [Matth. 24, 24.] warnt und spricht: „Es werden falsche Christi und Propheten aufstehen, und große Zeichen und Wunder thun, daß auch die Auserwählten, so es möglich wäre, verführet werden“ etc. Da ist nun jedermann hinan gegangen, und wenn jemand dawider predigt, sagen sie: Da ist das Zeichen vor Augen. So hat der Teufel bisher seinen Muthwillen getrieben, und [ich] glaube wohl, daß auch viel rechtschaffene Zeichen sind geschehen. Warum soll man denn ihnen nicht folgen? Hier hörst du es, er will nicht leiden, daß du ihn zum Lügner machest, sondern versuchst dich. Darum sollst du dich vom ersten Worte nicht lassen treiben; wenn er gleich einen Engel vom Himmel sendete, und hieße dich abtreten, sollst du es schlecht nicht glauben noch nachfolgen, sondern das erste gelegte Wort soll geschehen wider alle Wunderzeichen und Lehre, von GOtt oder dem Teufel, woher sie kommen.

15. Deß liest man auch ein Exempel im ersten Buche der Könige [Cap. 13, 9. ff.] von einem frommen Propheten, von GOtt geschickt gen Bethel, dem von ihm geboten war, er sollte nicht essen noch trinken an dem Orte, zu Bethel bei dem Könige; das er auch that. Aber als er wieder heimging, kommt ein anderer Prophet zu ihm, und sagte, er sollte mit ihm heimgehen und essen. Als er sich aber weigerte, sprach der andere zu ihm [V. 18.]: „Ich bin auch ein Prophet, wie du, und ein Engel hat mit mir geredet, durch des HErrn Wort, und gesagt: Führe ihn wieder heim, daß er Brod esse und Wasser trinke“; da geht er mit ihm, läßt das erste Gebot fahren, meinet, weil jener sagte, GOtt hat es geredet, müßte er es thun. „Aber da sie aßen, kam das Wort des HErrn (spricht der Text, V. 20—24.) zu dem Propheten, der ihn wieder umgeführet hatte, und schrie zu dem andern, und sprach: Darum, daß du dem Munde des HErrn bist ungehorsam gewesen, und hast nicht gehalten das Gebot, das dir der HErr geboten hat, soll dein Leichnam nicht in deiner Väter Grab kommen. Und da er weg zog, fand ihn ein Löwe auf dem Wege und tödtete ihn.“

16. Siehe, was wollte dieser Prophet sagen? Kann er nicht sprechen: Hast du mich doch selbst betrogen? Was wird aber GOtt sagen? So wird er sprechen: Du solltest mich für den Mann gehalten haben, der nicht lügt, und darauf geblieben sein, das ich dir gebot, wenn ich dir gleich anders gesagt hätte. Daher sagt der Prophet Ezechiel [Cap. 14, 9.]: „Den Propheten, der in meinem Namen kommt, und Lügen predigt, habe ich betrogen“, das ist, ich habe geschafft, daß er komme und die Leute verführe.

17. Also sollten wir nun auch thun. Wenn ich recht Christum empfangen habe im Sacramente, habe sein Wort und Zeichen, so steht der Glaube und Gewissen in gutem Muthe. Wo nun der Tod herfiele, und GOtt ließe sich hören oder fühlen, er wollte ungnädig sein und mich nicht haben: dennoch soll ich nicht wanken noch zurück treten; ob auch Moses oder ein Engel, oder gleich Christus käme, doch soll ich daran halten, denn das Wort GOttes kann nicht lügen. Spricht aber dein Herz: Ja, ist doch das auch
\quellseite{183}
Gottes Wort? Antworte du: Er wird es wohl machen und deuten, wie er es meint. Also müssen wir ihn überwinden mit seinem eigenen Worte. Das ist fast der höchste Kampf und Streit, welchen die heiligen Väter wohl gehabt haben.

18. Also ist es hier mit Abraham gegangen. Das erste Wort mußte bleiben; wie aber das andere auch sollte wahr sein, und mit dem andern überein kommen, mußte er Gott befehlen. Derhalben, wer in der Ausrichtung will gelehrt sein, der ergreife die Sprüche, daß Gott nicht leugt, noch wandt, was er einmal sagt; was aber darnach dawider lautet, ist darauf zu weisen und [zu] deuten, daß er den Glauben versucht. Das ließ ihn ausdeuten; es ist aber ein schwer Stück.

19. Also sollten wir gethan haben, da das Papstthum und Möncherei aufgekommen sind; darauf hat jedermann gesehen, wie gewaltig es zunahm, und sagten: Ei, das ist recht und köstlich. Da war niemand so klug, der da sagte: Ob gleich das Papstthum so glücklich fortfährt, muß ja das erste Wort nicht fallen, sondern dennoch bleiben, daß der Glaube allein selig mache, und alle Lehre und Wesen, das auf Werke bauet, verdammlich seien. Aber das hat man lassen fahren; so ist der Teufel eingerissen mit falschen Zeichen und hat genommen, daß sich jedermann davor fürchtete, und konnte niemand wehren, bis es mit voller Gewalt überhand nahm, und des Glaubens Lehre ganz austilgte.

20. Dies ist die große Historie von dem starken Glauben des Abraham, darinne du siehst, wie ihm Gott so reichlich seinen Sohn wieder gibt, und ihn mit Segen überschüttet. Denn dieweil er fest hält am ersten Worte, so fährt er das Wort wieder um, und zeigt ihm, daß er ihn nur versucht habe. Da folgen die Widersprüche, die Gott thut. Nun, diesem Exempel nach, haben wir viel Geschicht in den Historien, wie auch unter den Juden viele ihre eigenen Söhne verbrannt haben [5 Mos. 12, 31.]; das war ihr geistlicher Orden, hielten es für das allerköstlichste Werk, hatten einen Abgott dazu vor der Stadt Jerusalem, mit Namen Moloch [2 Kön. 23, 10.], und ist viel greuliches Dinges daraus kommen.\footnote{Vgl. Tischreden Cap. 5, § 5 und Cap. 24, § 2. Malch, St. Louiser Ausgabe, Bd. XXII, 243 f. und 704 f.}

21. Denn weil dies Werk, so Abraham thut, ein köstlich Werk ist, ahmen sie ihm nach, bis an das babylonische Gefängniß. Manche feine Leute meinten, sie wollten Gott auch einen großen Dienst thun; auch so waren Propheten dabei, und stärkten sie mit diesem Exempel; da plumpte der Haufe hernach [Jer. 32, 35.]. Dazu kamen denn auch die rechten Propheten, und sprachen: Ihr Mörder und Bluthunde, ihr\footnote{„ihr“ fehlt in der Erlanger.} würgt eure Kinder dem Teufel. So fuhren jene wieder hervor: Wie dürft ihr dawider reden, das Gott so hoch gepreiset hat? Also ist manch edel Blut über dem Greuel vergossen, allein dadurch, daß sie auf dies Exempel drangen.

22. Nun, woran fehlt es denn? Warum gefällt es denn Gott nicht auch, und gefällt ihm hier von Abraham wohl? Das ist es, daß ich gesagt habe, wenn es Gott nicht Abraham geheißen hätte, würde er es ihm nicht wohl gedankt haben. Es ist ein groß Werk, aber Gott achtet es doch nicht. Es ist nur um deßwillen zu thun, daß es Abraham aus seinem Worte und Befehle thut; wo das nicht wäre, so hätte er greuliche Sünde damit gethan. Also sagt der Prophet Micheas [Cap. 6, 7.]: Was soll man Gott thun, daß ihm wohl gefällt, soll man ihm einen Sohn opfern? Nein, sagt er, Gott will sein nicht. Warum nicht, hat es doch Abraham gethan? Noch will er es nicht. Warum das? Wenn er es geboten hätte, daß man das Wort hätte, so wäre es wohl gethan; wo nicht, so stehe nur beileibe davon ab. Denn die Narren haben dem Exempel gefolgt, und nicht gesehen, daß hier Gottes Wort und Glaube ist; aber mit ihnen ist nichts denn Vermessenheit und eigener Dünkel, damit sie wollen ein eigen Werk und Dienst thun, der ihm wohlgefalle, dadurch die Sünde zu tilgen und den Himmel [zu] verdienen; ist eine lautere eigene erdachte Andacht, welche vor Gott ein Greuel ist; aber was nach seinem Worte gethan wird, das hat er gerne.

23. Also haben wir auch bisher gethan, daß wir der Heiligen Exempel genommen haben, und mit Haufen sind hinnach gefahren, welches Gott nicht haben will. Der König Salomon hat Gott einen Tempel gebauet, und ist darum gelobt [1 Kön. 6, 1. ff.], so sind wir auch darauf gefallen, und [haben] gemeint, wer eine Kirche stiftete, hätte großes Verdienst bei Gott. Nicht also; sondern, wenn du nun Salomon bist, so
\quellseite{184}
baue auch eine. Er hatte Gottes Wort dazu, so hast du nichts, und willst Gott meistern, was er ihm soll gefallen lassen. Auch hat er es umsonst gethan und im Glauben, nicht, daß er ihm dadurch Gott gnädig machte; du kommst aber und willst Gott mit dem Werke kaufen, und dir zu eigen machen. Derhalben ist es verkehrt Ding, wenn du sie, die Heiligen, ergreifst mit den Crempeln.

24. Darum ist es gefährlich, Heiligen-Legenden predigen, ausgenommen, die in der Schrift stehen, welche alle auf den Glauben gerichtet sind. Zu andern Legenden ist nicht so der Glaube angezeigt, sondern die Werke aufgeworfen, wie viel und lange sie gefastet, gewacht und streng Leben geführt haben. Darum hält man sie für heilig, und will jedermann ihnen so nachahmen; damit geht der Glaube zu Boden.

25. So lasse nun solches fahren. Siehe, was diese Crempel lehren. Die sagen nicht, ob die Väter viel gefastet haben oder nicht, sondern daß sie Gottes Wort gehört haben, darinne gewandelt, und sie das gehalten. Wo das nicht vorgegangen ist, ist Abraham auch nicht gegangen. So lehrt auch die Epistel an die Hebräer [Cap. 13, 7.], wir sollen dem Glauben der Heiligen nachfolgen, spricht er, nicht ihren Werken. Es muß ein jeglicher sein Werk haben, aber alle müssen wir einerlei Glauben haben.

26. Das sei darum gesagt, daß man müsse davon zu richten, und niemand keinen Gottesdienst anrichte, was Gott nicht befohlen hat, sonst geht man in der Blindheit, und hält es für gute Werke. Denn wenn man schon sonst grobe Sünde thut, sieht doch jedermann, daß [es] nicht recht ist; hier will man sich noch davon rühmen. Darum ist Gott keinem Dinge so feind, als der eigenen Andacht; [sie] ist auch die schändlichste Plage auf Erden. Derhalben siehe nun, ist das Exempel des heiligen Vaters so übel gerathen, daß doch so rein beschrieben ist, was will mit unsern Werken werden? Was wollten unsere tollen Narren thun, die nicht halb so klug oder fromm sind, als jene? Noch dürften sie sich rühmen, Gott lasse sie nicht irren.

27. Bisher\footnote{Erlanger: Bis hieher.} haben wir nun gehabt das edle und feine Exempel des Glaubens, wie der fromme Vater Abraham so fest auf das Wort hat gehalten, daß er auch wider Gott selbst gestritten und gewonnen hat. Nun sagt er hier weiter in der Historie, daß ihm Gott aufs neue wiederum verheißt mit einem Eide, daß er seinen Samen segnen wolle etc., und lautet also:

V. 15—17. Und der Engel des Herrn rief Abraham abermal vom Himmel, und sprach: Ich habe bei mir selbst geschworen, spricht der Herr, dieweil du solches gethan hast, und hast deinen einigen Sohn nicht behalten, daß ich deinen Samen segnen und mehren will, wie die Sterne am Himmel, und wie den Sand an dem Ufer des Meeres; und dein Same soll besitzen die Thore seiner Feinde.

28. Hier sagt er unter andern, daß sein Same soll besitzen, das ist, einnehmen die Thore, das ist, Städte, Land und Leute seiner Feinde, und wiederholt die Worte und Verheißung, die er ihm zuvor gethan hatte, ehe denn Isaak geboren war. Und spricht also:

V. 18. 19. Und durch deinen Samen sollen alle Völker auf Erden gesegnet werden, darum, daß du meiner Stimme gehorchet hast. Also kehrete Abraham wieder zu den Knaben, und machten sich auf, und zogen mit einander gen Bersaba, und wohnete daselbst.

29. Diesen Spruch soll man lassen stehen nach der Historie, wie er klingt und lautet. Schlechte, geringe Worte sind es, wie du siehst; aber Gott hat eine andere Weise zu reden, denn die Menschen. Wenn er das Maul aufthut, und ein Wort läßt gehen, so gilt es. Man muß seine Predigt nicht so stückeln, wie Menschenlehre, hier ein Stück, und da ein Stück; da sei keusch, da gehorsam, da faste, da bete, daß es getheilt Wesen ist. Er aber, wenn er redet, redet er eben als er es fasset und hat; wie er es allzumal auf einen Haufen in der Hand hat, so gibt er es auch mit Haufen heraus.

30. Also hat er mit diesen kurzen Worten das ganze Evangelium und Reich Christi gefaßt, also, daß es niemand kann genug ausstreichen. Abraham hat es wohl verstanden, die Propheten auch, aber wenig Juden, wie sie es auch heutiges Tages nicht verstehen. Wenn wir nicht mehr hätten, hätten wir dennoch Schrift genug daran. Ich halte auch, die heiligen Väter haben an dem Spruche alle Bücher gehabt. Was die Schrift fassen will, das faßt sie alles auf einen Haufen,
\quellseite{185}
daß man nicht darf viel Auszüge, sondern mehr Einzüge machen.\footnote{Das heißt: Man soll es nicht einschränken, sondern weit spannen. Vergleiche § 34.}

31. Was ist es denn gesagt: „In deinem Samen sollen alle Völker gesegnet werden auf Erden“? Denn segnen muß man nicht verstehen, wie es die Juden deuten, daß es so viel heiße, er werde das jüdische Volk preisen unter den Heiden; also, daß sie nur sehen auf die Zunge, als sei es wohl reden, etwas Gutes wünschen; sondern Gottes Werke sind rechtschaffen, gehen zu Grund und Boden, ist alles Ernst und vollkommenlich Ding. Darum hat das Wort segnen alle seine Werke mit einander in sich; denn [Segen] will er lassen ausgehen, nicht allein über die Juden, sondern über alle Welt, und soll geschehen durch Isaaks Samen.

32. Wie will nun das zugehen? Also: Wir sind durch und in Adam alle gefallen in Fluch und Sünde, und auch verdammt durch die Sünde [Gal. 3, 10.] zum ewigen Tode. Der Fluch liegt uns auf dem Halse, darinne werden wir geboren. Alles, was von einem Weibe kommt, ist ein Kind des Fluches und Zornes [Eph. 2, 3.], des Todes und ewiger Verdammniß, unterthan dem Teufel. Das ist ein Fluch über die ganze menschliche Natur. Denn wie Adam ist, so zeugt er Kinder, eitel verdammte und verfluchte. Das bezeugt er hier erstlich mit dem Worte „segnen“. Denn, sollen sie gesegnet werden, so müssen sie vorhin nicht gesegnet sein. Wenn der Segen vorhin da wäre, was dürften wir dieses Samens? sondern, weil er Abraham solchen Segen zusagt, so ist es so viel geschlossen: Ihr seid allzumal verflucht.

33. Also stößt das Wort Segen allein nieder alles, was Mensch und Natur ist und vermag. Laß nun auftreten einen heidnischen Menschen, oder einen von uns, die sich Christen nennen, und doch Christum und den Glauben nicht haben, mit aller Klugheit, guten Werken, Tugenden und freiem Willen und ganzem Leben, so ist es alles nichts, und verflucht vor Gott. Ursache, es kommt aus dem Segen nicht; geht es nicht hieraus, so ist es nicht recht, ja, es ist alles des Teufels. Er spricht nicht: Die Heiden werden sich selbst segnen, oder ihnen den Segen erlangen, der Segen muß aus ihnen machen. Weil er nun nicht aus ihnen wächst, so müssen sie verflucht sein.

34. Darum sage ich: Wenn Gott redet, thut er den Mund weit auf, ist nicht enge gespannt, sondern trifft die ganze Welt, und schließt sie ganz und gar in Fluch; wie Paulus sagt Röm. 11, 32. und Gal. 3, 22.: „Die Schrift hat es alles beschlossen unter die Sünde, daß er sich aller erbarme, und die Verheißung käme durch den Glauben an Jesum Christum.“ Er flucht zwar nicht, sondern zeigt uns an, was wir sind, nämlich, daß wir gar im Fluche stecken: er will aber segnen, nicht zwei oder drei Völker, sondern die ganze Welt. Daran sollten wir je Zeugniß genug haben, daß alle Menschen mit allen ihren Werken und Gottesdienste des Teufels sind. Abraham hat daraus können richten alles Wesen und Lehre, wie es ihm vor Händen kommen ist, und hat so gesagt: Ist der Segen da, so ist es gut; ist es aber ohne den Segen, so ist es verloren, es gleiche, wie es wolle.

35. Der Segen hat nun sollen gehen über alle Heiden, doch also, daß es geschehe in Abrahams Namen. Das ist das Evangelium. Aber der Spruch will nicht sagen, daß alle Heiden werden den Segen empfahen. St. Paulus ist ein Meister, solche Sprüche auszulegen. Es ist hierinne verheißen (sagt er [Röm. 10, 18.]), daß Gott wollte das Evangelium lassen ausgehen. [Er] spricht nicht, alle werden es annehmen. Wie so? Also: Das Evangelium ist ein solches Wort, welches uns Gottes Gnade und Barmherzigkeit anbietet, durch Jesum Christum verdient und mit seinem Blute erworben, welcher ist das Lamm Gottes, das die Sünde trägt [Joh. 1, 29.], und läßt aller Welt verkündigen: Glaube und halte [dich] an mich, lasse die Werke fahren, so bist du gesegnet, und los vom Tode und allem Unglück. Der Segen ist gegangen so weit die Welt ist, gekommen beide über Heiden und Juden, und geht noch immer weiter.

36. So viel will dieser Spruch: Ich will einen Segen ausgehen lassen über alle Völker unter der Sonne. Denn man muß zulassen, daß das Wort: alle Heiden, hebräisch [כל גוים] Col Gojim, auch die Juden begreife; wie Moses zu ihnen sagt im 5. Buche [Cap. 4, 7.]: „Wo ist ein solch [גוי] Goi, dem Gott so nahe ist“, als dir? Darum soll der Zusatz „über alle Heiden“ heißen, was für Volk unter der Sonne ist, Gott gebe, es sind Juden oder Heiden.
\quellseite{186}

37. Siehe, so weit streckt sich der Spruch, daß GOtt das Evangelium wollte lassen kund werden über alle Welt. Das heißt er „gesegnet“. Der Segen ist da, wird jedermann angeboten, aber es nimmt ihn nicht jedermann an, und ist ein ernstlicher Segen, der da geht über den ganzen Menschen. Glaubst du an den Samen, so sind dir vergeben deine Sünden mit einander, und der Tod ist aufgehoben, die Hölle überwunden, und der Himmel dein. Das alles wünschst und bringt er dir auf einmal, aber wenig sind ihrer, die den Segen annehmen. GOtt grüßt alle Welt, aber wenig danken ihm. Noch bleibt es wahr, sie werden alle gesegnet, das ist, es wird ihnen allen angetragen und vorgelegt; daß aber nicht Alle Christen werden, macht, daß sie solchen heilsamen Segen ausschlagen, ja, zum meisten Theil verfolgen. Es ist nicht darum gesagt, daß GOtt wollte jedermann bekehren. St. Paulus [Eph. 1, 13.] deutet es allein auf das Evangelium, daß es sei ein Geschrei, das er läßt ausgehen über jedermann, es soll eitel Segen sein; wen es trifft, den trifft es, wer es fähet, der hat es.

38. Weiter ist den Juden hier auch ein Ziel gesteckt, daß sie sich nicht können rühmen, sagt Paulus [Gal. 3, 18.], als hätten sie es irgend mit verdient. Denn zu der Zeit, da GOtt den Segen Abraham verhieß, sind wir nicht gewesen, die Juden auch nicht; da war weder Volk noch gute Werke, weder Mosis noch Propheten noch jemand, daß sich nicht jemand dürfe rühmen, als habe er es erworben. Daß uns das Evangelium gegeben wird, ist allein GOttes bloße Güte, die er uns ausschüttete, ehe jemand darum bat oder darauf dachte. Denn er hatte es im Sinne zu geben, ehe wir geboren waren. Hätte er uns zuvor lassen fromm sein, und hätte vorhin den Juden Gesetze gegeben, ehe er Abraham die Zusagung that, möchten sie sagen: Das haben wir verdient. Wäre das Gesetz nicht, und hätten wir nicht gute Werke gethan, so wäre Christus nicht gekommen. Den Ruhm will er legen, kommt zuvor; ehe die Heiden da sind, die gesegnet sollen sein, ist die Zusagung und Gnade da.

39. Darum, wer GOttes Wort hört und erkennt, der muß sagen: So hätte ich es nimmermehr verdient, noch gebeten, oder in Sinn genommen. Darum sind die Worte so tief, daß sie nicht zu [er]gründen sind. Wo von Segen gesagt wird, da ist das Evangelium; wo das Evangelium ist, da ist GOtt mit Christo und allen Gütern. Also kann GOtt mit Einem Worte alle Dinge fassen. So ist aufgehoben aller Juden Pracht und Heiligkeit, und der Heiden Klugheit; [es] müssen alle solche sagen: GOtt habe Lob und Dank, der uns gnädige Verheißung gegeben hat.

40. Und dieser Text wird nun die Glosse und das Licht sein aller Gesetze, die hernach geschrieben sind. Wenn die Juden sagen: Wir haben von GOtt empfangen die zehn Gebote, welche so wir halten, sind wir fromm, und dürfen Christi nicht, so ist dieser Spruch da, und schließt, daß GOtt aus lauter Gnade zusagt und verheißt. Wo bleiben aber die zehn Gebote? Antwort: Das ist das erste Wort, ehe die zehn Gebote gegeben und geschrieben sind; da gilt nicht sagen, daß GOtt gelogen habe. Der Same und Segen über alle Welt ist da; wer den nicht hat, der ist verloren.

41. Warum sind denn die Gebote gegeben? Nicht darum, daß du durch deine Werke fromm sollest werden, oder könntest das Gesetz erfüllen, sondern, daß du den Leib zähmen sollst, ehrbarlich leben, und das gelobte Land einnehmen; und vornehmlich, daß man daran lerne erkennen, wer wir sind, und wo es uns fehlt, daß wir desto mehr Durst nach dieser Verheißung haben. Es ist eben darum gegeben, daß dieser Same gefördert, und der Segen getrieben würde.

42. So haben die Propheten gethan und geschrieben: Herr GOtt, wer kann es halten? Je mehr du gebietest, je weniger man thut. Wir sollten GOtt trauen, und sein Gebot halten; das thun wir nicht, finden nichts mehr am Gesetze, denn daß nichts Gutes an uns ist. Darum hat es Moses gegeben, daß es den Fluch offenbar machen soll; und wenn wir denn solches sehen und fühlen, daß wir dann anheben und sagen: Nun komm, Herr, gib uns den Segen, erlöse uns von diesem Fluche. Derhalben hilft das Gesetz nichts zum Gewissen, der Segen muß es allein thun.

43. Also sind alle Schrift dahin geordnet, daß sie diesen Spruch immer treiben. [Es] ist auch im ganzen Mose der Hauptspruch, und geht alles darauf, was zuvor und hernach steht, auf daß alle Welt müsse sagen: Es wird nichts Gutes daraus, hier ist eitel Tod und GOttes Zorn: nun ist das dein Wort, darum halte, was
\quellseite{187}
du gesagt hast, daß du uns den Segen gebest, und den Samen kommen lässest, der uns solches gebe, und helfe vom Tod durch Leben, von Sünde durch Gerechtigkeit.

44. Siehe, das ist das Evangelium, dadurch danieder liegt und verdammt ist alle Welt unter Sünde, Tod und Teufel, mit aller ihrer Herrlichkeit, Frömmigkeit und guten Werken. Denn es findet sich, daß es nicht der Segen ist. Wiederum ist aufgerichtet vor und ohne alle Werke lauter Gnade durch Christum, daß sich niemand rühme, sondern jedermann Gott danke, daß er den Samen erweckt hat, durch welchen der reiche und ewige Segen kommt. Das ist die ganze Theologie auf einen Haufen, davon bisher keine Gelehrten noch hohe Schulen ein Wort verstanden haben.

45. Nun müssen wir auch das Wörtlein „Samen“ handeln. Hier hätten die Juden noch eine Ausflucht, daß sie möchten sagen: Warum soll solches eben auf Christum gezogen werden? Sind doch Isaak, Jakob und alle Propheten und Heiligen Abrahams Same? Antwort: Wir können dazu nicht weiter, denn daß wir Trotz bieten, nicht allein den Juden, sondern aller Welt, Engel, Tod, Teufel, daß sie uns einen Samen geben, der mit solchen Wundern bestätigt ist, als unser Christus. Das Werk ist da, davon die Worte lauten, und ist bekannt, daß er Abrahams Same ist; so muß alle Welt, ohne ihren Dank, bekennen, daß er so weit gepredigt ist, als die Welt reicht. Das ist und wird nimmermehr keinem widerfahren; dazu auch, daß ihn die Heiden haben angenommen, wie es zuvor verkündigt ist, und ihn jedermann dafür hält, und das Evangelium für Wahrheit hält. Alle Welt hat mit ihm zu schaffen, und sind dennoch Heiden, die billig dem sollten feind sein; bekennen dazu, daß er nicht allein Abrahams Same ist, sondern hören und predigen auch, daß er den Segen gebracht hat, den Tod überwunden samt der Sünde und Hölle. Das hat er auch bewiesen mit großen Wundern vom Himmel durch das Evangelium und durch die Apostel. Da war eitel Wunder, damit Gott die Predigt bestätigte. Nun gib uns einen andern, der das alles gethan hat, und noch immerdar geht durch die Welt, also, daß sich alle Gewalt dawider legt, die Juden auch; noch hat es nicht geholfen, es hat müssen durchbrechen; wenn noch so mächtige Könige auf einem Haufen zusammen thäten, so könnten sie es nicht so weit bringen.

46. Weiter ist in den Worten auch begriffen alles, wie wir glauben von Christo, daß er mußte von einer Jungfrau geboren werden, leiden und sterben, vom Tode auferstehen, gen Himmel fahren, und zu der Rechten Gottes sitzen und regieren. Denn, daß seine Mutter sollte eine Jungfrau sein, folgt daher: Alle Heiden sind verflucht; was vom Fleisch geboren, ist schon ein Kind des Zorns und Todes [Eph. 2, 3.], so muß er je nicht von Fleisch und Blut kommen, oder von Mann und Weib geboren werden, soll er nicht verflucht sein, noch gesegnet werden, sondern den Segen bringen, daß alle, so gesegnet werden, durch ihn gesegnet werden. Wiederum muß er dennoch ein natürlicher Mensch sein, Fleisch und Blut haben wie Abraham. Denn die Schrift heißt gemeiniglich die Kinder Samen. Soll er nun ein natürliches Kind und Sohn sein, so muß er von einem Weibe kommen, denn es heißt niemand ein Kind, es bringe denn sein Wesen von der Mutter; doch, soll es nicht verflucht sein, so muß er nicht natürlicher Weise empfangen werden. Darum erzwinget es sich, daß die Mutter nicht von einem Manne muß schwanger sein, und doch recht schwanger, daß sie eine rechte natürliche Mutter heiße, und das Kind ihr Blut und Fleisch sei. Darum ist kein ander Mittel, denn daß er ohne alle Verrückung aus dem Heiligen Geiste empfangen wird. So bringt es mit sich, daß sie sei eine Jungfrau, und des Kindes Mutter, in rechtem natürlichem Wesen.

47. Solches haben auch aus diesem Spruche geführt die Propheten, als Jesaias [Cap. 7, 14.]: „Siehe, eine Magd oder Jungfrau wird schwanger werden, und einen Sohn gebären.“ Also nennt ihn die Schrift allein fructum ventris, eine Frucht des Mutterleibes, nicht eine Frucht der Lenden; das ist, ein Kind allein von einem Weibsbild, nicht von einem Manne. Siehe, die Sprüche sind alle daher geflossen, daß sie solches gesehen haben; also, daß der Spruch nicht auf Christum gehe, wie im 51. Psalm, V. 7., steht: „Siehe, ich bin in Unmugend gemacht, und meine Mutter hat mich in Sünden empfangen.“ Weil er nun ohne den Fluch ist, so ist er auch ohne Sünde und Tod. Also findet sich’s, daß er nicht von Mannes Samen empfangen ist.
\quellseite{188}

48. Weiter: Alle, die geboren werden vom Weibe, müssen natürliche Menschen sein, und sterben, wie andere Weiber Kinder; so muß dieser Same auch, wie alle andere Menschen, sterben. Wiederum, sollte er der sein, der alle Welt segnete, [so] kann er nicht leiblich auf Erden bleiben, daß er König zu Jerusalem wäre, und regierete die ganze Welt; das Reich ist zu groß. Soll er selbst gegenwärtig regieren, muß er in solchem Wesen sein, daß er bei jedermann und an allen Orten sei; darum muß er ein leiblicher Mensch geboren sein, wie andere, durch den Heiligen Geist, daß er im sterblichen Leben wandele, aber durch und aus demselbigen Wesen in ein ander, unsterblich, geistlich Wesen treten, daß er bei jedermann im Herzen regiere; sonst könnte er es nicht ausrichten, daß er in jedermanns Herzen sei und mit ihm rede, und schaffe über tausend Meilen Weges. Darum müßte er vom Tode auferstehen, gen Himmel fahren, und sitzen an dem Orte, da er alle Creaturen sehen, erfüllen, und in der Hand haben kann, kriegt Gewalt über alle Dinge, Engel und Teufel.

49. Auch muß er GOtt sein. Denn soll er so groß Gut verheißen, und den Segen geben, muß er auch die Gewalt haben; sonst wäre es lächerlich, wenn er etwas zusagte, und könnte nicht helfen. Weil er nun uns läßt zusagen, daß er uns vom Tode und des Weltfürsten Reiche erledigen, und ewiges Leben geben will, muß er Gewalt haben und allmächtig sein, dazu auch Weisheit, daß er uns auch könne und müsse zu erhalten im Leben, wider Tod und Sünde. Sonst, wenn er nur gleich wünschte, und doch nicht bei mir wäre, könnte er mir nicht helfen. Darum, weil er, so wir mitten im Tode und Hölle stecken, kann selig machen, muß er die allerhöchste, göttliche Macht haben, und ihm alles unter den Füßen liegen [Ps. 110, 1.], welche Gewalt niemand hat denn Gott allein. Weil sie denn dem Samen gegeben ist, muß er unwidersprechlich ein wahrhaftiger GOtt sein, und doch der Same bleiben, das ist, beide wahrer GOtt und Mensch sein. So haben wir alles mit einander in Christo, Geburt, Sterben, Auferstehung und sein ewiges Regiment. Wer möchte solches alles mit so kurzen Worten fassen, wenn nicht der Heilige Geist selbst redete?

50. Das sei genug geredet von dem Texte, der wohl werth wäre, daß man ihn zeichnete mit güldenen Buchstaben. Es liegt aber allein daran, daß man den Segen recht erkenne. Es soll ein Segen sein, den GOtt gibt; was GOtt aber segnet, ist mit der That und Werk gesegnet. Denn er segnet nicht wie die Leute, allein mit Worten, da\footnote{Erlanger: daß.} nichts mehr nachfolgt, sondern ist so ein lebendiger Segen, der alles Gut mitbringt. Es ist nicht ein Gruß allein, sondern ein Werk und geschäftig Ding, das bald von statten geht. Die Frucht, so solchem Segen folgt, ist, daß es niemand wehren noch hindern kann. Wenn GOtt segnet, so ist reden und thun Ein Ding, wie im 33. Psalm, V. 9., steht: Dixit, et factum est, „er sagt, so stehet es da.“ Darum muß der Segen ausrichten, was er redet, und thätig sein. Geht es von GOtt, so muß es treffen Sünde, Tod, Teufel, Hölle; so tief geht sein Wort.

51. Daher wirst du so viel Sprüche finden in den Propheten, hieraus gezogen. Es ist eine große, reiche Fundgrube, daraus sie gegraben und genommen haben alles, was sie geschrieben haben von dem zukünftigen Christo und seinem Reiche. Nun, das sehen die Juden nicht, sind verblendet, hängen an Mose, und stehen auf ihren Werken, verstehen nicht, daß er zuvor durch den Spruch so große Dinge zusagt, ehe er das Gesetz gegeben hat.

52. Das ist die Historie, darinne wir lernen, wie gewaltig GOtt redet. Und wenn gleich die Worte solches alles nicht geben, so müßte es doch die Historie durch die geistliche Deutung geben. Also werden wir an Isaac auch sehen, wie Christus darinne abgemalt ist. [Es] ist so ein weites Wort, daß [es] über alle Welt geht, und faßt Christum und sein Reich gar. So wollen wir nun die Deutung handeln.

53. Ich habe oft gesagt, wer in der Bibel wohl studiren will, sonderlich die geistliche Deutung der Historien, daß er es alles ziehe auf den Herrn Christum. Denn es ist alles zu thun mit den Deutungen um das Predigtamt, oder Evangelium und Glauben, in welchen Christi Reich steht, wie wir hier auch sehen werden. Nun hält unser Glaube also, daß wir fest glauben, daß der Same Abrahams wahrhaftig GOtt und Mensch ist; und wiewohl es unmöglich war, daß die Person, als GOtt, sollte geboren werden, leiden, sterben, auferstehen, und andere Dinge thun, die er auf Erden gethan hat, darum, daß
\quellseite{189}
die Gottheit gar unleidlich und unsterblich ist: dennoch ist es wahr, daß man sagt, Gottes Sohn hat gelitten und ist gestorben.

54. Als man auch von einem Menschen sagt, welches Seele niemand verwunden kann mit Eisen, denn sie ist ein Geist: dennoch, wenn einer in [den] Kopf gehauen ist, spricht man recht: Der Mensch ist wund. Item, so sagt man auch: Der Mensch ist todt, so doch die Seele nicht sterben kann. Also schreibt man alle Werke der ganzen Person zu, ob es wohl allein der Leib thut oder leidet, darum, daß beide zusammen \unsicher{geeinet} und Eins werden. Also auch hier; wiewohl der Mensch Christus alles gethan und gelitten hat, doch, weil er mit der Gottheit Eine Person ist, spricht man auch, GOtt habe gelitten, sei gestorben, und auferstanden. Das ist unser Glaube, das lasset uns nun hier sehen.

55. Isaak wird auf den Altar gelegt, und gebunden, der Vater nimmt ihn bei dem Kopf, will ihn hauen und schlachten, also, daß er in der Wahrheit geopfert ist, wenn man des Vaters Herz ansieht. Denn da geht frei, beide Wille und Hand, [er] hat sich's allerdings erwogen, der Tod ist ganz in seines Vaters und seinen Augen, daß ihrer keiner anders weiß noch sieht: noch bleibt er lebendig, und wird für ihn der Bock geopfert. Also stirbt er im Schein und Ansehen, und bleibt doch im Tode leben. Das ist der Herr Christus, der scheint und stellt sich in aller Weise, als sterbe Gottes Sohn; doch stirbt er nicht, sondern der Mensch stirbt wahrhaftig: das ist der Tod in der Dornhecke. Also sagt nun die Schrift, daß GOtt wahrhaftig seinen Sohn hingegeben hat; noch ist er nicht gestorben, sondern der Mensch. So ist Isaak eine Figur der Gottheit, der Tod der Menschheit Christi.

56. Zum andern ist hier bedeutet die Auferstehung von den Todten. Weil Isaak dem Abraham verheißen war, er sollte werden ein Vater vieler Heiden, und wird doch zum Tode geantwortet, daß die Sprüche mußten beide wahr sein, daß er ein Vater vieler Völker sollte werden, und doch dem Vater unter den Händen sterben. Da ist das Reich Christi gar abgemalt. Also ist der Spruch wahr, daß er sollte ein König und Heiland sein aller Welt [Sach. 9, 9. Joh. 4, 42. 1 Tim. 4, 10.] Wie stellt er sich aber dazu? [Er] läßt sich an das Kreuz schlagen und stirbt; schickt sich eben wie Isaak, als er sollte ein Vater werden. Nun hat das die Schwachgläubigen feindlich geärgert.

Die Apostel fahren zu, meinen, das wird der König werden, und das Volk erlösen [Luc. 24, 21.], so wird er auf das schändlichste gemartert und gekreuzigt. Wo bleibt nun das Königreich? Sie waren zu schwach dazu, daß sie also hätten gedacht: Christus wird ein König und Heiland sein; ob er gleich jetzt stirbt, so ist der Vater mächtig, daß er ihn wieder lebendig mache. Also mußte es auch geschehen. Christus fällt und stirbt, steht aber wieder auf, wird ein König, wie geschrieben ist; ist König auf Erden, sitzt aber nicht mehr also, wie zuvor, doch regieret er in unsern Herzen, hat allein darinne zu schaffen, predigt, und erleuchtet sie; und Summa, was in der ganzen Christenheit geschiehet, das thut er.

57. So ist es nun wahr, Christus ist ein König, und regieret, so weit die Welt ist [Ps. 72, 8.], daß ihm alles unterthan muß sein; lebt in \unsicher{Wahrheit} auch leiblich, aber in unsterblichem und geistlichem Leben. Welche Vernunft nun solches könnte ausdeuten, die könnte nicht weiter sehen, denn daß er sollte mit großem Pracht auf Erden sichtlich regieren. Wenn das wäre, so wäre kein Glaube in der Welt. So besteht nun die Schrift: Isaak bleibt leben, [es] ist nur eine Gestalt und Schein des Todes da. So auch mit Christo. Er stirbt frei dahin in seiner und aller Welt Augen, kommt aber durch den Tod in das Leben, wird ein König und regieret ewiglich. Also ist hier das ganze Königreich gefaßt. Denn wo die Auferstehung steht, da steht es alles mit einander.

58. Weiter, der Bock in der Hecke mit den Hörnern, wie gesagt, ist Christus nach der Menschheit. Nun sehen wir, wie das Evangelium geht, und die ganze Schrift zuvor sagt, daß es Rumor anrichtet, hängt nicht zwischen Rosen und Lilien, Purpur noch Wolle; da ist nichts Weiches, sondern eitel Dornen; darinne ist der Bock verwirrt, kann nicht heraus kommen. Das bedeutet alles das Predigtamt. Horn heißt in der Schrift: mysterium verbi evangelii, welches ist die Kraft Gottes. Denn wie der Bock mit den Hörnern stößt, so stößt auch die Predigt des Evangelii vor den Kopf, was groß ist; [es] kann nicht leiden, daß jemand klug und fromm sei, und sich darauf verlasse. Wenn das Thier nun geräth unter die Dornhecken, da verwirrt sich's darinne.

59. Dornen sind die Gottlosen, sonderlich die das Evangelium verfolgen, stachlicht und unschlachtiger
\quellseite{190}
Art. Denn die frommen Christen lassen sich gerne stoßen. Aber das andere ist ein zänkisches, habersüchtiges Volk, wie Paulus sagt Röm. 2, 8.: His, qui sunt ex contentione, wie man ihnen predicat, so ist es nicht recht; [sie] sind eitel Hecken, die endlich in das Feuer gehören. So siehst du in der Figur, daß es so sei und gehen muß. Darum soll es niemand wundern, ob sich Hader über dem Evangelio [er]hebt, daß man die Christen verjagt und verfolgt; sprich nur: Es muß so sein, und nicht anders, der Widder muß in den Hecken hangen. Was der Teufel groß achtet auf Erden, das muß alles Christum verfolgen; so stößt es sich mit ihm, schadet ihm aber nichts; er hat harte Hörner im Kopfe, doch verwirrt er sich darin, wird gerissen und geopfert.

60. Was deutet es aber, daß Abraham die Esel und Knaben weit herunter vom Berge läßt stehen, und er geht mit dem Sohne allein hinan? Unsere Legales\footnote{Das ist: Gesetzlehrer.} und Wertprediger haben es also gedeutet: Wenn man zur Kirche geht, Vesper, Metten 2c. betet und singt, soll man die fleischlichen Gedanken draußen lassen, und sich bereiten mit guten Gedanken, daß es ein gut Werk sei; und dabei bleibt es. Wir aber haben so gesagt, daß der äußerliche Gottesdienst dazu diene, daß man immer die Schrift lehre und treibe, und baue den Glauben damit, und daß niemand so beten und lesen soll, daß er wolle damit ein gut Werk thun, sondern das Gewissen zu bauen, und den Glauben zu stärken. Darum taugt das Gauklerwerk gar nichts.

61. Der Esel ist der faule Schelm, der alte Adam, Blut und Fleisch, der muß auch da sein, wenn man Gott dienen soll; denn Christi Reich muß auf Erden regieren in unsern Herzen. Die Knaben aber, die den Esel führen, sind die Gesetztreiber. Darum muß man die zwei Parteien von einander scheiden. Die zwei gehen hinan auf den Berg; der Esel mit den Knechten bleibt weit davon. Mit dem Evangelio führt man nicht das Maul oder Füße zu Gott, sondern das Herz; man kann es nicht in die Ohren, Mund, oder Zunge fassen, sondern im Herzen muß es leben. Das sieht die Welt nicht, der Knecht und Esel werden es nicht gewahr, es siehet's niemand denn Abraham. Der Esel aber muß dennoch ein Regiment haben; der Leib ist im äußerlichen Regimente dem Gesetze unterworfen, wie der Esel bewahrt wird von den Knechten.\footnote{Wittenberger und Erlanger: vom Knechte.} Daß damit bedeutet sind alle, die das Volk äußerlich regieren, und das Fleisch im Zwange halten. Der faule alte Mensch kann nicht vor Gott kommen, dennoch ist Isaac da, das ist, die Seele steckt im Leibe, die kommt vor Gott.

62. So soll man es nun scheiden. Durch den Glauben an das Evangelium treten wir zu Gott, opfern uns Gott selbst, wie hier Isaac, und werden ihm zu eigen gegeben. So ist das Herz mit dem Glauben geopfert, aber der Esel bleibt hinter dem Berge, kann nicht herzu, denn er ist noch voll Sünde. Das Herz wird wohl durch den Glauben rein, ist ein neuer Mensch und heiliges Opfer, aber der alte Adam ist noch voll böser Lust, den muß man zähmen, bringen und zwingen, daß er von Tag zu Tage gefegt werde,\footnote{Erlanger: wird.} und immer Meister habe, die ihn züchtigen.

63. Darum muß man das Gesetz nicht führen auf das Gewissen, sondern demselbigen\footnote{„demselbigen“ fehlt in der Erlanger.} nichts predigen denn den Glauben, ohne alles Gesetz, daß es frei gehe zu Christo. Danach soll man auch predigen auf den faulen Esel, daß man ihn im Zwang halte und dämpfe. Dazu soll das Gesetz allein dienen. Seele und Gewissen soll keinen Knecht haben, der sie führt, sondern frei bleiben. Nun sind hier zweien Knechte, das sind Moses und die\footnote{„die“ fehlt in der Erlanger.} Propheten, die das Gesetz gelehrt und getrieben haben in das Volk. Das ist die Deutung, das ist das Predigtamt gar, Gesetz und Evangelium, trifft beide, innerliche und äußerliche Menschen, und die ganze Welt.

64. Das Letzte ist vom Berge Moria. Abraham hat gewohnt im Philisterlande, bei Gaza, von dannen ist er hinauf gezogen drei Tagereisen; am dritten Tage hat er den Berg von ferne gesehen; ohngefähr einen Armbrustschuß oder zween, da ließ er die andern hinter sich, und ging mit Isaac allein hinauf. Der Berg heißt Moria, und ist eben der, darauf hernach der König Salomo den Tempel bauete [2 Chron. 3, 1.], nicht, wie etliche sagen, da Christus gekreuzigt ist. Und heißt darum Moria, daß er spricht: Der Herr siehet darauf, denn \unsicher{יראה} Moria heißt auf deutsch so viel, als, des
\quellseite{191}
Herrn Gesicht, daß der Herr darein sieht. Und ist des Berges Name also geblieben darum, daß Gott so genau auf Abraham sah, wehrte ihm, daß er den Sohn nicht würgte, und zeigte ihm den Widder; also, daß es auch heißen mag, der Herr zeigt. Das deutet man dahin, daß Christus sollte kommen auf den Berg, da der Tempel auf gebauet ward, und das Evangelium predigen, und da lassen ausgehen in die ganze Welt; wie in den Geschichten der Apostel [Cap. 1, 2.] und im Jesaia [Cap. 2, 2.] steht.

65. Wir bleiben aber auf dem geistlichen Moria, weil wir das andere alles geistlich machen, darum lassen wir es sein die christliche Kirche, welche ist das Reich Christi, oder den Glauben. Im Alten Testament war nun geboten, daß man kein Opfer thun sollte an irgend einer Stätte, denn da es Gott erwählte; er zeigte aber nirgend hin, denn nur zu Jerusalem. Damit ist bedeutet, daß es nicht gilt, was wir für Werke thun, sie geschehen denn im Glauben; [es] ist schon das Urtheil gefällt, wie Paulus [Röm. 14, 23.] sagt: „Alles, was nicht aus dem Glauben ist, das ist Sünde“; ob du gleich dich todt marterst, da fragt Gott nichts\footnote{Jenauer: nicht.} nach, sieht auch alle dein Leiden nicht an, wie auch Paulus [1 Cor. 13, 3.] sagt, wo es nicht im Glauben geht. So thaten die Juden auch große Werke und Opfer; aber Gott fragte nichts darnach, darum, daß [es] nicht an dem Ort geschah, den ihnen Gott erwählte. Also soll das Wort niederschlagen alles, was nicht aus Gottes Wort und aus dem Glauben ist, wie jetzt sind die Stifte, Klöster, und allerlei geistlich Wesen.

66. Diesen Berg sahe Abraham noch von ferne. [Es] war noch nicht recht Moria, gehörte noch weit hin, ja viel hundert Jahre, ehe Christus kommen, und sein Reich durch das Evangelium sollte anheben, noch hat er es gesehen; wie Christus sagt im Johanne [Cap. 8, 56.]: „Abraham, euer Vater, ward froh, daß er meinen Tag sehen sollte, und er sahe ihn, und freuete sich.“ So steht nun das Opfer alles im Glauben, davon wird es ein heilig, rein und angenehm Opfer.

67. Das sei genug von der geistlichen Deutung dieser Geschichte, die beide, den Glauben und das Kreuz, predigt, daß man den alten Menschen tödte. Summa, es ist alles zu thun um Christum. Die Schrift will nichts wissen, denn von Christo; noch hat uns der Teufel genarrt, als wäre es nicht genug daran, und müßten\footnote{Jenauer: müßten.} dieweil eitel Menschenträume lernen und predigen. Zuletzt ist noch ein Stück in dem Capitel, wie Abraham gefragt ward, daß sein Bruder auch Kinder hatte von seiner Milla, nämlich acht Söhne, und von dem Kebsweibe vier; wie folgt.

. V. 20—24. Nach diesen Geschichten begab sich’s, daß Abraham angesagt ward: Siehe, Milca hat auch Kinder geboren deinem Bruder Nahor, nämlich Uz, den Erstgebornen, und Bus, seinen Bruder, und Kemuel, von dem die Syrer kommen, und Chesed,\footnote{Jenauer: Sabed; Wittenberger und Erlanger: Sebed.} und Haso, und Pildas, und Jidlaph, und Bethuel. Bethuel aber zeugete Rebekka \&c.

68. Das ist darum geschrieben nach der Historie, daß wir hernach hören werden, im vier und zwanzigsten Capitel, wie Abraham seinem Sohne ein Weib hat genommen, von demselbigen Geschlechte, nämlich Rebekka, seines Bruders Sohnes, Bethuel, Tochter, daß man wisse, wo sein Weib herkomme. [Es] ist aber auch in der geistlichen Deutung willen hier angezeigt, daß damit bedeutet wird, wie Christus, als er wieder auferstanden ist von [den] Todten, darnach ein Weib, das ist, die Heidenschaft, zu sich nimmt; [er] nimmt sie nicht aus seinem Lande, sondern läßt sie holen aus einem fremden Lande, hat die Juden verworfen, die Heiden lassen zu sich fordern und holen, durch die Apostel, wie wir hernach hören werden.
\quellseite{192}

\section{Das dreiundzwanzigste Capitel}

V. 1—20. Sara ward hundert sieben und zwanzig Jahr alt, so lange lebte sie, und starb in der Hauptstadt, die heißt Hebron,\footnote{Randglosse: Hebron. Hebron ist KirjathArba, spricht Moses, das ist, die Vierstadt. Denn die hohen Hauptstädte waren vor Zeiten alle Arba, das ist, in vier Theile getheilt, wie Rom, Jerusalem und Babylon auch. 1 Mos. 10, 10.} im Lande Canaan. Da kam Abraham, daß er sie klagte und beweinete. Darnach stund er auf von seiner Leiche, und redete mit den Kindern Heth, und sprach: Ich bin ein Fremder und Einwohner bei euch, gebet mir ein Erbbegräbniß bei euch, daß ich meinen Todten begrabe, der vor mir liegt. Da antworteten Abraham die Kinder Heth, und sprachen zu ihm: Höre uns,\footnote{„uns“ fehlt in der Wittenberger und in der Erlanger.} lieber Herr, du bist ein Fürst Gottes unter uns, begrabe deinen Todten in unsern köstlichen Gräbern, kein Mensch soll dir unter uns wehren, daß du in seinem Grabe nicht begrabest deinen\footnote{Wittenberger und Erlanger: den.} Todten. Da stund Abraham auf, und bückte sich vor dem Volk des Landes, nämlich, vor den Kindern Heth, und er redete mit ihnen, und sprach: Ist es euer Gemüth, daß ich meinen Todten vor mir begrabe, so hört mich, und bittet für mich gegen Ephron, den Sohn Zoar, daß er mir gebe seine zwiefache Höhle, die er hat am Ende seines Ackers. Um ein gleiches Geld gebe er mir sie unter euch zum Erbbegräbniß; denn Ephron wohnte unter den Kindern Heth. Da antwortete Ephron, der Hethiter, Abraham, daß zuhörten die Kinder Heth, vor allen, die zu seiner Stadt Thore aus- und eingingen, und sprach: Nein, mein Herr, sondern höre mir zu. Ich gebe dir den Acker, und die Höhle darinnen dazu, und gebe dir es vor den Augen der Kinder meines Volks, zu begraben deinen Todten. Da bückte sich Abraham vor dem Volk des Landes, und redete mit Ephron, daß zuhörte das Volk des Landes, und sprach: Willst du mir gehorchen, so bitte ich,\footnote{„bitte ich“ fehlt in der Erlanger.} nimm von mir des Ackers Geld, das ich dir gebe, so will ich meinen Todten daselbst begraben. Ephron antwortete Abraham, und sprach zu ihm: Mein Herr, höre doch\footnote{„doch“ fehlt in der Erlanger.} mich, das Feld ist vierhundert Seckel\footnote{Randglosse: Seckel. Seckel ist ein Gewicht an der Münze, ein Ort eines Guldens, denn vor Zeiten man das Geld also wog, wie man jetzt mit Golde thut.} Silbers werth; was ist das aber zwischen mir und dir? Begrabe nur deinen Todten. Abraham gehorchte Ephron, und wog ihm das Geld dar, das er gesagt hatte, daß zuhörten die Kinder Heth, nämlich, vierhundert Seckel Silbers, das im Kauf gang und gäbe war. Also ward Ephrons Acker, darinne die zwiefache Höhle ist gegen Mamre über, Abraham zum eignen Gut bestätiget, mit der Höhle darinnen, und mit allen Bäumen auf dem Acker umher, daß die Kinder Heth zusahen, und alle, die zu seiner Stadt Thore aus- und eingingen. Darnach begrub Abraham Sara, sein Weib, in die Höhle des Ackers, die zwiefach ist, gegen Mamre über, das ist Hebron, im Lande Canaan. Also ward bestätiget der Acker und die Höhle darinnen Abraham zum Erbbegräbniß von den Kindern Heth.

1. Das ist je ein närrisch Capitel anzusehen. Was hat er so viel Worte zu machen über solchen geringen Dinge, wie Abraham eine Grube kauft, da er einen Todten einlegt? Droben [Cap. 22, § 30 ff.] haben wir gehört, daß er das ganze Königreich Christi in drei oder vier Worten faßt, da es wohl noth wäre, daß er hunderttausend Worte machte, und macht so wenig, daß er nicht weniger machen kann, sagt es alles in Einem Worte: Segen. Hier wäscht er ein ganz Capitel von dem, da doch nichts an gelegen ist. Es ist vor der Vernunft doch ein lauter närrisch Ding, und unnützes Geschwätz. Was wollen wir nun daraus machen? Nach der Historie weiß ich nichts daraus zu machen, denn daß es zuwider den abergläubigen und hoffährtigen Heiligen geschrieben ist, welche die Gewissen gerne enge spannen, und meinen, wer Gott dient, müsse nicht mit solchen Weltlichen umgehen, daß man hieraus lerne, sich schicken in die Weise der Leute, mit denen man\footnote{„man“ von uns herübergenommen aus der alten Ausgabe Walch; in den andern Ausgaben: er.} umgeht; also,
\quellseite{193}
daß es sei ein Exempel der Liebe. Wenn man so ferne kommt, daß der Glaube und Evangelium bleibe, soll man sich nicht daran kehren, wie die Leute leben, was sie für Gesetze haben, sondern sich immerdar schicken nach ihrem Wesen; und daß nicht Sünde ist, daß ein Christ mit Leuten handelt und umgeht, wird ein Kaufmann, hält die Weise, wie die andern, und drückt sonderlich aus das Geld geben,\footnote{Jenaer: und.} das gänge und gäbe ist.

2. Es ist nicht eine geringe Tugend der Liebe, wie Paulus sagt, daß man sich schicke in jedermanns Weise, denn wir werden es nicht dahin bringen, daß sich eine ganze Welt nach uns richte. [Es] ist genug, wenn wir das erhalten, daß [das] Evangelium und der Glaube bleibe. Darnach müssen wir denken, die Welt wird doch nicht thun, was ich will, darum will ich mich darein schicken, und thun, was sie will. Darum sagt Paulus [Röm. 12, 11.]: „Schicket euch in die Zeit“, und machet euch eben jedermann. Item [Phil. 4, 5.]: *\unsicher{Modestia} vestra nota sit omnibus hominibus,* „eure Lindigkeit lasset allen Menschen kund sein“, das ist, daß ihr es also haltet in eurem Wesen, daß niemand über euch klagen mag; also auch 1 Cor. 10, 33.: Thut, wie ich thue; „ich mache mich jedermann gefällig“, das ist, ich weiche und gebe jedermann nach, lasse mein Recht und Sinn fahren, wo ich sehe, daß ein anderer dawider ist, daß ich mich je zu ihm thue, auf daß er Gefallen an mir habe.

3. Es ist eine feine Regel: Wer da will jedermann gefallen, der lasse ihm jedermann gefallen; und wiederum: Wer ihm niemandes Weise will gefallen lassen, wird auch wieder niemand gefallen. Also soll ein Christ thun. Das thun jene nicht. Wenn man ein wenig zu viel lacht, und über die Schnur fährt, so schnurren und purren sie; [es] sagt jedermann, daß sein Sinn der beste sei, das andere tauge alles nichts. So ist ein jeglicher, der in seiner eigenen Weise geht, läßt sich immer dünken, der andern sei nicht so rein als seine, ihm gefällt niemand wohl; und will gleichwohl, daß seine Weise jedermann gefalle. Das sind falsche Gleisnerwerke, die nicht in der Liebe gehen.

4. Darum muß das Werk Abrahams gepriesen werden, ob es auch nicht werth scheint zu schreiben. Könnte er keine andere Liebe finden, nach dem hohen Exempel des Glaubens? Ja, es ist freilich das größte, höchste Werk, das hierinnen nach der Historie angezeigt ist. Wer die Liebe hat, darf nicht viel mehr.

5. Daß Abraham den Acker um vierhundert Seckel kauft, muß man wissen, daß siclus ist zweierlei, gülden und silbern. Der silberne ist an der Münze fast so viel als ein Ortugulden, so machen vier Seckel einen Gulden. So geben die vierhundert Seckel eben hundert Gulden; so theuer hat er den Acker gekauft mit Bäumen, und alles, was darauf war; das war sein Schatz. Dabei ist das Kreuz angezeigt, wie er ein Gast gewesen ist, und nicht eines Fußes breit gehabt in dem Lande, das ihm Gott zugesagt hatte, wie Stephanus sagt in der Apostelgeschichte [Cap. 7, 5.].

6. Das Grab in der Höhle ist gewesen gegen Mamre über, welche ist Hebron (sagt der Text); [das] ist ihr alter Name gewesen, vielleicht von einem großen Manne, der da gewohnt hat, wie oben gesagt ist im vierzehnten Capitel, V. 24., von den drei Männern, Aner, Escol und Mamre. Diese Stadt nennt auch der Text KirjathArba, das ist, eine Vierstadt, auf griechisch τετράπολις, das ist, eine von den großen Hauptstädten, die gemeiniglich in vier Theile getheilt waren, wie Jerusalem und Babel, als oben im zehnten\footnote{In den Ausgaben: „elften“. Luther weist hierdurch hin auf das, was in 2 § des zehnten Capitels gesagt ist.} Capitel gesagt.

7. Hier ist nun unter andern Stücken auch angezeigt, wie man sich christlich halten soll gegen den Todten, daß es ein gut Werk sei, Todte zu begraben, sonst hätte er nicht so viel Worte davon gemacht, wo und wie Sara begraben ist.

8. Item, daß Abraham sie beklagt und beweint habe, ist auch darum geschrieben, daß es nicht böse ist, daß man sich bekümmert, betrübt, und Leide trägt um der Leute willen, die uns absterben. Wiewohl wir allzumal sterben müssen, doch sind wir also verfast durch die Liebe, daß ein jeglicher dem andern soll das Leben gönnen, wie wir auch alle verbunden sind, daß wir arm sind, und im Schweiße leben: noch soll die Liebe da zu schaffen haben, und sich um\footnote{„um“ fehlt in der Erlanger.} des Nächsten Armuth und andere Noth bekümmern. Nach dem Glauben dürfen wir weder Lebens noch Todes, aber nach der Liebe müssen wir uns des Nächsten annehmen; wie es dies Leben trägt,
\quellseite{194}
ist’s gut, daß wir uns mit ihm freuen, und wiederum, wie Paulus Röm. 12, 15. lehrt: „Freuet euch mit den Fröhlichen, und weinet mit den Weinenden.“

9. So soll sich ein jeglicher des andern annehmen, auch des Leibes Noth halben, wie es sein mag, auf daß nicht die hoffärtigen Heiligen kommen, und denken es in Wind zu schlagen, ein Christ möge guter Dinge sein, und leben, wie er wolle. Es soll keiner so ein großer Heiliger sein, der sich nicht treffen lasse und annehme der Nächsten Noth. Christus hatte auch genug vor sich; und wiewohl unter allen Wundern, die er that, war das größte, daß er das Evangelium predigte; nichtsdestoweniger half er allen Kranken, wie Matthäus [Cap. 8, 17.] schreibt: „Er hat unsere Schwachheit auf sich genommen, und unsere Seuche hat er getragen“ [Jes. 53, 4.].

10. So mußt du nicht anders richten von\footnote{Erlanger: mit.} einem Christen denn nach dem Glauben und der Liebe. Nach dem Glauben nimmt er sich keines Dinges an auf Erden; nach der Liebe nimmt er sich alles Dinges an. GOtt will nicht durch das Evangelium die Natur ausreißen, sondern läßt bleiben, was natürlich ist, richtet es aber auf die rechte Bahn. Natürlich ist es, daß ein Vater sein Kind, das Weib den Mann lieb hat, und fröhlich sei, wenn es ihm wohl geht, und wiederum. Deß nimmt sich ein Christ vor GOtt nach dem Glauben nicht an, wenn gleich alles zu Boden ginge, aber nach der Liebe soll man sich deß annehmen, als unserer eigenen Noth, und nach der Liebe richten. Sonst hätte er nicht lassen schreiben, daß der große Patriarch sein Weib beweinte, wo es nicht daher sollte dienen. Also schickt uns GOtt zu, daß solche Bewegungen des Herzens bleiben, und will dennoch, daß wir sie überwinden mit dem Glauben, und darum nicht verzagen, noch von GOtt fallen. So bist du denn ein recht christlicher Mensch, der GOtt in allen Sachen gehorcht, und sich\footnote{„sich“ fehlt in der Erlanger.} doch bekümmert, wenn einem Andern Ungemach widerfährt. Darum siehst du, daß nichts ist, daß etliche Geister vorgeben, man solle der Creaturen so gar los sein, daß man sich gar nichts annehme, wie es andern geht.

11. Das ist nun von der Liebe Abrahams

geredet, daß er so große Sorge trägt, daß er ein Eigenthum habe für den Todten. Vor der Welt ist es also gethan, daß jedermann denkt, wo er bleiben, und was er Eigenes haben wolle im Leben, vor dem Tode. Hier ist es ungefähr. Im ganzen Leben ist er ein Fremdling, daß er nicht eines Fußes breit hatte für seinen und seines Weibes Leib; jetzt, weil sie todt ist, und es nimmer genießt, will er erst eigen Gut haben; daß man ihn wohl möchte für einen Narren halten, daß er jetzt will sorgen, und hat es zuvor nicht gethan. Es ist alles hergegangen aus dem Glauben, wie auch dies hernachhin.

12. Welches, als ich halte, aus der heiligen Schrift hergekommen ist. Denn die Heiden, sonderlich die Römer, haben ihre Todten verbrannt; die Schrift aber heißt das begraben, und fließt, halte ich, aus dem dritten Capitel dieses Buches, da GOtt [Cap. 3, 19.] spricht zu Adam: „Bis daß du wieder zur Erde werdest, davon du genommen bist, denn du bist Staub und sollst zu Staub werden.“ Da steht das Wörtlein [עפר] Aphar, das heißt solche Erde, die lose ist, als wäre sie zermalmt, wie die Maulwürfe aufwerfen, oder so man aus dem Grabe wirft, die zerstoben ist. Wenn man nun den Menschen unter die Erde begräbt, wird wieder Erde daraus; so geht es alles aus GOttes Worte, wie er gehört hatte von seinen Vätern.

13. Nun hat er gewisse Zusagung, daß das Land, da er inne war, sollte seinen Erben gegeben werden. Auf das Wort verließ er sich, achtete nicht, ob er es beim Leben nicht hätte, mußte wohl, daß sie es nach seinem Tode würden einnehmen. So hat er gesehen auf die Auferstehung der Todten, und gedacht, es sollte ein ander Leben und Land sein, das ihm GOtt geben würde; wie es die Epistel an die Hebräer [Cap. 11, 10.] deutet. Das hat er damit bedeutet, daß er dies Leben verachtet, und so groß achtet auf den Todten, als der da sieht auf ein ander Leben. Auch zeigt GOtt damit an, wie er die Todten in seine Sorge und Hut nimmt, also, daß diese Sara jetzt besser versorgt ist, denn da sie lebte. Als sollte er sagen: Weil ihr lebt, sorgt ihr immer; so müßt ihr doch das Sorgen lassen, wenn ihr sterbet,\footnote{In den alten Ausgaben: gestorben.} und mich für euch sorgen lassen. Also liegt die Sara jetzt auf ihrem eigenen Gute und Lande, das ihnen
\quellseite{195}
Gott zugesagt hatte, welche vor gar nichts Eigenes davon hatte. Damit tröstet Gott uns alle, und zeigt heimlich mit an die Auferstehung des Fleisches, weil er sich mehr der Todten annimmt denn der Lebendigen; das würde er nimmer thun, wenn er nicht im Sinne hätte, die Todten wieder zu erwecken. Das ist der heiligen Väter Trost gewesen, darauf sie gestorben sind, daß sie jetzt liegen und ruhen in Gottes Schooße.

14. Nun, die geistliche Deutung zu rühren, was das Begräbniß bedeutet, haben wir aus St. Paulo, Röm. 6, 3. f.: „Wisset ihr nicht, daß alle, die wir in Jesum Christum getauft sind, die sind in seinen Tod getauft? So sind wir je mit ihm begraben durch die Taufe in den Tod“ etc. Item, V. 5.: „So wir aber sind sammt ihm gepflanzet worden zu gleichem Tode, so werden wir auch der Auferstehung gleich sein.“ Wenn man einen Christen tauft, so tauft man ihn auf den Tod, daß er solle sterben. Gleichwie man ein Korn in die Erde wirft, so man säet; das wird begraben, also, daß es in der Erde sterben muß, und wiederum heraus wachsen ein neues Korn. So werden wir in der Taufe gepflanzet mit dem Herrn Christo, gehen durch und in der Taufe in sein Sterben und Grab, und durch den Tod in ein neues und ewiges Leben. Das ist nun die Summa: Wer getauft ist, hat nicht mehr zu thun, denn daß er begraben werde. Denn ob ich gleich im Leibe bin, gehe ich doch damit um, daß ich ihn immer zuschurre, krieche je tiefer und tiefer in das Grab hinein, bis ich verfaule; also, daß ich täglich der Begierde abbreche, die Lust dämpfe und würge, und der Welt absterbe, so lange bis ich gar verwese unter der Erde; so muß ich der Welt zugescharrt werden, soll ich mit Gott leben. Das sei das Mysterium.

\section{Das vierundzwanzigste Capitel}

V. 1—67. Abraham war alt und wohl betaget, und der Herr hatte ihn gesegnet allenthalben. Und sprach zu seinem ältesten Knecht seines Hauses, der allen seinen Gütern vorstand: Lege deine Hand unter meine Hüfte, daß ich dich lasse schwören bei dem Herrn, dem Gott des Himmels und der Erde, daß du meinem Sohne kein Weib nehmest von den Töchtern der Canaaniter, unter welchen ich wohne, sondern daß du ziehest in mein Vaterland und zu meiner Freundschaft, und nehmest meinem Sohn Isaak ein Weib. Der Knecht sprach: Wie, wenn das Weib mir nicht wollte folgen in dies Land, soll ich denn deinen Sohn wieder bringen in jenes Land, daraus du gezogen bist? Abraham sprach zu ihm: Da hüte dich vor, daß du meinen Sohn nicht wieder dahin bringest. Der Herr, der Gott des Himmels, der mich von meines Vaters Hause genommen hat, und von dem Lande meiner Freundschaft, der mir geredet hat, und mir auch geschworen und gesagt: Dies Land will ich deinem Samen geben, der wird seinen Engel vor dir her senden, daß du meinem Sohn daselbst ein Weib nehmest. So aber das Weib dir nicht folgen will, so bist du dieses Eides quitt; allein, bringe meinen Sohn nicht wieder dorthin. Da legte der Knecht seine Hand unter die Hüfte Abrahams, seines Herrn, und schwur ihm solches. Also nahm der Knecht zehn Kameele, von den Kameelen seines Herrn, und zog hin, und hatte mit sich allerlei Güter seines Herrn, und machte sich auf und zog gen Mesopotamiam zu der Stadt Nahor. Da ließ er die Kameele sich lagern außen vor der Stadt bei einem Wasserbrunnen, des Abends um die Zeit, wenn die Weiber pflegten heraus zu gehen und Wasser zu schöpfen, und sprach: Herr, du Gott meines Herrn Abraham, begegne mir heute, und thue Barmherzigkeit an meinem Herrn Abraham. Siehe, ich stehe hier bei dem Wasserbrunnen, und der Leute Töchter in dieser Stadt werden heraus kommen, Wasser zu schöpfen. Wenn nun eine Dirne kommt, zu der ich spreche: Neige deinen Krug, und laß mich trinken, und sie sprechen wird: Trinke, ich will deine Kameele auch tränken; daß sie die sei, die du deinem Diener Isaak bescheret habest, und ich an derselben erkenne, daß du Barmherzigkeit an meinem Herrn gethan hast. Und ehe er ausgeredet
\quellseite{196}
hatte, siehe, da kam heraus Rebekka, Bethuels Tochter, der ein Sohn der Milka war, welche Nahors, Abrahams Bruders Weib war, und trug einen Krug auf den Achseln, und sie war eine sehr schöne Dirne von Angesicht, noch eine Jungfrau, und kein Mann hatte sie erkannt, die stieg hinab zum Brunnen, und füllete den Krug und stieg herauf. Da lief ihr der Knecht entgegen, und sprach: Laß mich ein wenig Wassers aus deinem Kruge trinken. Und sie sprach: Trinke, mein Herr. Und eilends ließ sie den Krug hernieder auf ihre Hand, und tränkte ihn. Und da sie ihn getränket hatte, sprach sie: Ich will deinen Kameelen auch schöpfen, bis sie alle trinken.\footnote{1, Erlanger: getränken. In der Bibel: getrunken.} Und eilte und goß den Krug aus in die Tränke, und lief aber zum Brunnen zu schöpfen, und schöpfete allein seinen Kameelen. Der Mann aber wunderte sich ihrer und schwieg stille, bis er erkennete, ob der Herr seine Reise hätte lassen wohl gerathen, oder nicht. Da nun die Kameele alle getrunken hatten, nahm er eine güldene Stirnspange, eines halben Sekels schwer, und zween Armringe an ihre Hände, zehn Sekel Goldes schwer, und sprach: Meine Tochter, wem gehörst du an? Das sage mir doch. Haben wir auch Raum in deines Vaters Hause zu herbergen? Sie sprach zu ihm: Ich bin Bethuels Tochter, des Sohnes Milka, den sie dem Nahor geboren hat. Und sagte weiter zu ihm: Es ist auch viel Stroh und Futter bei uns, und Raums genug zu herbergen. Da neigte sich der Mann, und betete den Herrn an, und sprach: Gelobet sei der Herr, der Gott meines Herrn Abrahams, der seine Barmherzigkeit und seine Treue nicht verlassen hat an meinem Herrn; denn der Herr hat mich den Weg geführet zu meines Herrn Bruders Hause. Und die Dirne lief, und sagte solches alles an in ihrer Mutter Hause. Und Rebekka hatte einen Bruder, der hieß Laban. Und Laban lief zu dem Manne draußen bei dem Brunnen, und als er sahe die Stirnspangen, und Armringe an seiner Schwester Händen, und hörte die Worte Rebekka, seiner Schwester, daß sie sprach: Also hat mir der Mann gesagt, kam er zu dem Manne, und siehe, er stand bei den Kameelen am Brunnen. Und sprach: Komm herein, du Gesegneter dem Herrn, warum stehst du draußen? ich habe das Haus geräumet, und für die Kameele auch\footnote{„auch“ fehlt in der Erlanger.} Raum gemacht. Also führete er den Mann in das Haus, und zäumete die Kameele ab, und gab ihnen Stroh und Futter, und Wasser zu waschen seine Füße, und der Männer, die mit ihm waren, und setzte ihm Essen vor. Er sprach aber: Ich will nicht essen, bis daß ich zuvor meine Sache geworben habe. Sie antworteten: Sage her. Er sprach: Ich bin Abrahams Knecht, und der Herr hat meinen Herrn reichlich gesegnet, und ist groß worden, und hat ihm Schafe und Ochsen, Silber und Gold, Knechte und Mägde, Kameele und Esel gegeben, dazu hat Sara, meines Herrn Weib, einen Sohn geboren meinem Herrn in seinem Alter, dem hat er alles gegeben, was er hat. Und mein Herr hat einen Eid von mir genommen, und gesagt: Du sollst meinem Sohne kein Weib nehmen von den Töchtern der Cananiter, in deren Lande ich wohne, sondern zeuch hin zu meines Vaters Hause, und zu meinem Geschlecht, daselbst nimm meinem Sohne ein Weib. Ich sprach aber zu meinem Herrn: Wie, wenn mir das Weib nicht folgen will? Da sprach er zu mir: Der Herr, vor dem ich wandele, wird seinen Engel mit dir senden, und deinen Weg fertigen, daß du meinem Sohne ein Weib nehmest von meiner Freundschaft und meines Vaters Hause. Dann sollt du meines Eides quitt sein, wenn du zu meiner Freundschaft kommst; geben sie dir sie\footnote{„sie“ fehlt in der Erlanger.} nicht, so bist du meines Eides quitt. Also kam ich heute zum Brunnen, und sprach: Herr, GOtt meines Herrn Abraham, hast du meinen Weg gefertiget, daher ich gereiset bin: siehe, so stehe ich hier bei dem Wasserbrunnen. Wenn nun eine Jungfrau heraus kommt zu schöpfen, und ich zu ihr spreche, gib mir ein wenig Wasser zu trinken aus deinem Kruge, und sie wird sagen: Trinke du, und ich will deinen Kameelen auch schöpfen, daß die sei das Weib, das der Herr meines Herrn Sohne bescheret hat. Ehe ich nun solche Worte ausgeredt hatte in meinem Herzen: siehe, da kommt Rebekka heraus mit einem Kruge auf ihrer Achsel, und gehet hinab zum Brunnen und schöpfet. Da sprach ich zu ihr: Gib mir zu trinken, und sie nahm eilend den Krug von ihrer Achsel, und sprach: Trinke, und deine Kameele will ich
\quellseite{197}
auch tränkten. Also trank ich, und sie tränkte die Kameele auch. Und ich fragte sie, und sprach: Weß Tochter bist du? Sie antwortete: Ich bin Bethuels Tochter, des Sohnes Nahor, den ihm Milca geboren hatte. Da legte ich eine Stirnspange an ihr Antlitz, und Armringe an ihre Hände, und neigte mich und betete den HErrn an, und lobete den HErrn,\footnote{„den HErrn“ fehlt in der Erlanger.} den GOtt meines Herrn Abraham, der mich den rechten Weg geführet hat, daß ich seinem Sohne meines Herrn Bruders Tochter nehme. Seid ihr nun die, so\footnote{„ich“ fehlt in der Erlanger.} an meinem Herrn Barmherzigkeit und Treue thun wollt,\footnote{„wollt“ fehlt in der Erlanger.} so saget mir es; wo nicht, so saget mir’s aber, daß ich mich wende zur Rechten oder zur Linken. Da antwortete Laban und Bethuel, und sprachen: Das ist von dem HErrn ausgegangen; darum können wir nichts dawider reden, weder Böses noch Gutes. Da ist Rebekka vor dir, nimm sie, und ziehe hin, daß sie deines Herrn Sohns Weib sei, wie der HErr geredet hat. Da diese Worte hörte Abrahams Knecht, bückete er sich dem HErrn zu der Erde, und zog hervor silberne und güldene Kleinode und Kleider, und gab sie Rebekka; aber ihrem Bruder\footnote{Erlanger: ihren Brüdern.} und der Mutter gab er Würze.\footnote{Jenauer: köstliche Früchte.} Da aß und trank er, sammt den Männern, die mit ihm waren, und blieben über Nacht allda. Des Morgens aber stund er auf, und sprach: Lasset mich ziehen zu meinem Herrn. Aber ihr Bruder und Mutter sprachen: Laß doch die Dirne einen Tag oder zehn bei uns bleiben, darnach sollst du ziehen. Da sprach er zu ihnen: Haltet mich nicht auf, denn der HErr hat meinen Weg gefertiget, lasset mich, daß ich zu meinem Herrn ziehe. Da sprachen sie: Lasset uns die Dirne rufen, und fragen, was sie dazu sagt. Und riefen der Rebekka, und sprachen zu ihr: Willst du mit diesem Manne ziehen? Sie antwortete: Ja, ich will mit ihm. Also ließen sie Rebekka, ihre Schwester, ziehen mit ihrer Amme, sammt Abrahams Knecht und seinen Leuten. Und sie segneten Rebekka, und sprachen zu ihr: Du bist unsere Schwester, werde zu viel tausendmal tausend, und dein Same besitze die Thore seiner Feinde. Also machte sich Rebekka auf mit ihren Dirnen, und setzte sich auf die Kameele, und zog dem Manne nach. Und der Knecht nahm Rebekka an, und zog hin. Isaak aber kam vom Brunnen des Lebendigen und Sehenden, denn er wohnete im Lande gegen den Mittag, und war ausgegangen zu trachten\footnote{In der Bibel: beten.} auf dem Felde, um den Abend. Und hub seine Augen auf, und sahe, daß Kameele daher kamen. Und Rebekka hub ihre Augen auf und sahe Isaak, da fiel sie vom Kameele, und sprach zu dem Knechte: Wer ist der Mann, der uns entgegen kommt auf dem Felde? Der Knecht sprach: Das ist mein Herr. Da nahm sie den Schleier\footnote{In der Bibel: Mantel.} und verhüllete sich. Und der Knecht erzählete Isaak alle Sache, die er ausgerichtet hatte. Da führete sie Isaak in die Hütte seiner Mutter Sara, und nahm die Rebekka, und sie ward sein Weib, und gewann sie lieb. Also ward Isaak getröstet über seiner Mutter.

1. In den nächsten zwei Capiteln haben wir von Todten, Sterben und Begraben gehört;\footnote{„gehört“ fehlt in der Erlanger.} hier hören wir wiederum vom Leben. Denn weil jene Frau gestorben ist, müssen wir wieder eine neue Braut haben. Darum ist das Capitel nicht mehr denn eine Historie von der Hochzeit, als es auf die Zeit zugegangen ist. [Es] möchte abermahl jemand sagen, welch ein unnützer Schwätzer Moses sei, daß er so viel Worte macht um der Braut willen, beschreibt es so genau, daß er auch von Kameelen allein sagt. Nun, wir wollen am ersten die Historie handeln. Derselbigen nach siehst du aber einen Trost, den da haben alle Gläubigen, wie GOtt sich ihrer Sachen so hart annimmt.

2. Isaak ist verheißen, er soll ein Vater sein eines großen Volkes; so war es gewiß, daß er müßte ehelich werden, und ein Weib haben. Da sorgt GOtt auch für. Sein Leben, Sterben, Ehre, Gut, und was er hat, steht in GOttes Hand; also auch das Weib, das er haben soll, und doch nicht weiß, woher es kommen werde. Wir Narren thun alle Dinge zuvor, ehe wir GOtt darum fragen, und treiben Narrentheidung damit; darnach, wenn wir sehen, daß der Schimpf aus ist, und sich der Ernst hebt, murren wir und sind ungeduldig. Wenn wir Christen wären, und Mann und Weib so ansähen, daß
\quellseite{198}
es Gottes Creatur wäre, und daß wir alles müßten aus seiner Hand nehmen, daß wir auch nicht ein Tröpflein Wassers schöpfen möchten, Gott gäbe es denn, so würden wir keine Täckerei daraus machen. Die heiligen Väter haben es alles angenommen als eine gute Gabe von Gott; wir aber fahren dahin, als hätten wir es gethan, und Gott nicht darum müßte, noch sich’s annähme.

3. Siehe, also versorgt Gott die Seinen, ehe Abraham, Isaak und der Knecht, dazu sie mit ihrer Freundschaft, darauf gedenken. Damit je angezeigt ist, es soll alles gehen im Glauben, ihm befohlen sein alles, was wir haben, sonderlich der Ehestand, davon Früchte in aller Welt herkommen, und darum aufgezeigt ist, daß man die Leute regiere in Gottes Erkenntniß, und die Kinder lehre Gott fürchten, und also die Sorge Gott heim gebe.

4. Aber das läßt Abraham nicht nach; wiewohl er es, und Isaak, beide, zu Gott sehen, wo das Weib herkomme, thut er doch so viel, als ihm zu thun ist. Die Sorge soll er lassen und Gott trauen. [die] Arbeit soll er nicht lassen. Das Meine soll ich dazu thun, aber wie, wann, und wo es geschehe, soll ich ihm befehlen. Darum befiehlt er Gott, was er für eine soll nehmen; doch schickt er seinen Knecht aus unter die Freundschaft, weiß nicht wohin. Als sollte er sagen, wie er auch sagt: Du wirst wohl eine finden, sorge nur nicht, was und wie, du wirst einen großen Geleitsmann haben. Gott vom Himmel hat mir geschworen und gesagt: „Dies Land will ich deinem Samen geben, der wird seinen Engel vor dir her senden, daß du meinem Sohne daselbst ein Weib nehmest.“ Als wollte er sagen: Will er meinem Sohne das Land geben, so wird er meinem Sohne ja auch ein Weib geben; so trost und pocht er auf die Zusagung. Will es aber Gott nicht schaffen, so komme wieder, spricht er, er wird wohl finden, woher es kommen soll; ja, er trost und pocht noch mehr, daß er sagt: „Er wird seinen Engel vor dir her schicken“; so fest hält er an der Verheißung Gottes. Das ist ein Stück, daß er den Knecht abfertigt, daß er sucht, so weit als ihm zu suchen ist.

5. Item, er fährt noch weiter, nimmt einen Eid von dem Knechte, daß er je kein Weib nehme von den Töchtern des Landes, da er wohnte. Das wußte Abraham, daß die Amoriter, und die im Lande Canaan wohnten, alle sollten vertilgt werden, wie droben gesagt ist im fünfzehnten Capitel, V. 16. 21. [Sie] mußten alle hinweg, sollte sein Same das Land einnehmen. Aus dem Worte folgte, daß er sich nicht sollte unter das Volk mengen, noch mit ihnen Freundschaft machen, weil Gott schon das Urtheil gefällt hatte, daß sie alle verstoßen müßten werden. So hat nun Abraham, so viel in ihm war, Arbeit und Mühe nicht gespart, doch die Sorge Gott befohlen. Nun, der Knecht ist ein frommer, heiliger Mann, in großer Andacht und Trauen. Abraham wird sein Gesinde also gezogen haben, daß ihrer viel darunter gläubig sind gewesen; aber auch viel böse.

6. Daß er aber seine Hände auf Abrahams Hüfte legte, halte ich, wird ein Gebrauch gewesen sein, wie wir schwören auf das Evangelium oder die Heiligen; also, daß sonderlich die großen Herren die Weise gehabt haben, also einen Eid zu nehmen. Etliche legen es geistlich aus, weil Christus von Abrahams Blute und Fleische kommen sollte, habe er darum die Hand unter Abrahams Hüfte gelegt, daß der Eid auf Christum deutete. Das lasse ich gehen in seinen Würden.

7. Dabei hat die Historie auch das in sich, daß ein frommer Christ möge der Welt ihr Recht auch lassen. Sonst hätte man mögen sagen: Was geht er mit dem Gutelwerke um, daß er so Großes darauf legt, und dem Knechte mitgibt zehn Kameele, Gold und Edelsteine, treibt groß Gepränge und Geschenke, [er] hätte es wohl mögen armen Leuten geben? Es ist eben darum geschrieben, daß nicht jemand denke, es sei Sünde, ob er mit der Braut herrlicher fährt, denn sonst im gemeinen Leben, daß man es nicht so enge spanne. Es ist nicht Sünde, ob man eine Braut hübscher schmückt denn sonst eine Magd, und mehr Gepränge treibt, weil es diese Väter gethan haben. Gott läßt es geschehen der Hochzeit zu Ehren, daß sie also gepriesen werde.

8. Darum muß man der Welt ihr Recht lassen, ausgenommen den Ueberfluß, daß man eine Braut ziemlich schmückt, geht hin, isset und trinkt, auch daß man schön tanzt, man muß darüber kein Gewissen machen, allein, daß man mehrt, wo es Ueberfluß\footnote{Erlanger: überaus.} ist. Gott kann wohl
\quellseite{199}
zusehen, dem ehelichen Stande zu Ehren, daß man fröhlich ist; wenn man schon zuweilen ein wenig zu viel thut, soll darum niemand die Gewissen betrüben.

9. Aber wo man nicht Freude darinne sucht, sondern Säue daraus werden, wie wir Deutschen thun, das taugt nichts, da sollte man wehren. Also ist es auch mit dem Schmucke; jetzt muß man so viel Perlen und Seide haben, gerade als sollte die Braut nicht geschmückt sein, sondern sehen lassen, wie schwer sie tragen könne. Wenn das geschmückt heißt, so könnte man wohl einen Narren auch schmücken, der könnte des Dinges viel tragen; aber da fehlt es am Regimente, das sollte darein sehen, und Maße setzen. Es ist auch der tollen Männer Schuld, so den Weibern so viel geben, als sie um sich können hängen; so ist es ein tolles Thier, um ein Weib, das nicht zu sättigen ist mit Geschmuck. Das sage ich nun darum, daß Gott nicht darum zürnt, ob man sich ziemlich schmückt und wohl lebt.

10. So läßt nun Abraham den Knecht auch etliche güldene Kleinode und Schmuck der Braut mitbringen, eines von einem halben Goldsekel, die andern zwei von zehn Goldsekel. Ein Goldsekel ist zehn Ortgülden, das ist drittehalben Gulden, daß es alles bei dreißig Gulden werth gewesen ist.

11. Also möchte man nun vom Tanzen reden, wie wir vom Schmuck sagen. Daß man davon sagt, es bringe viel Reizung zu Sünden, ist wahr, wenn es über die Maße und Zucht fährt; aber es kann auch einer wohl mit einer buhlen, die weder Schmuck noch Schöne hat; denn die Liebe ist blind, fällt so schier auf einen Kuhdreck,\footnote{Erlanger: Dreck.} als auf ein Lilienblatt. Darum, weil Tanzen auch der Welt Brauch ist, des jungen Volkes, das zur Ehe greift, so es auch züchtig, ohne schandbare Weise, Worte oder Gebärde, nur zur Freude geschieht, ist es nicht zu verdammen. Das hat Gott hiermit angezeigt, daß er es also läßt schreiben, daß nicht die hoffärtigen Heiligen so bald Sünde daraus machen; wenn man es nur nicht in Mißbrauch bringt.

12. Weiter geht es auf der Braut Seite also zu, daß man sieht, wie die Tochter sein gezogen ist. Wie freundlich empfähet sie den fremden Knecht, wie willig und dienstlich erbeut\footnote{Erlanger: beut.} sie sich, daß sie ihm dient nach aller Nothdurft, läuft hin und sagt es ihrem Vater an, und schafft ihm Herberge. So hat man die Zeit die Kinder gezogen, daß sie gedicht sind gewesen zu dienen, jedermann willfertig zu sein, sonderlich den Gästen und Pilgern, gerne [zu] herbergen, welche Tugend oft in der Schrift gelobt wird. Item, hernach, da er sie heimführt, und sie auf das Kameel setzt, sieht den Isaak von ferne, fällt sie vom Kameele, schlägt den Schleier um sich, und verhüllt sich, hat sich fein demüthig und züchtig gestellt gegen den Leuten, und eine ehrbare Gebärde geführt, wie den Jungfrauen fein ansteht. [Es] ist auch der beste und lieblichste Schmuck, wenn sie züchtig mit Gebärden sind, die Augen niederschlagen, und nicht um sich werfen; [das] ist lieblich, denn alles Gold, so man mag tragen, wie der weise Mann [Sprüchw. 31, 10.] sagt: Es ist mit keinem Gelde zu bezahlen, wer ein züchtig Weib hat; [es] ist aber ein seltsam Thier. Denn es ist ein schwach Blütelein um ein Weib, kann sich nicht selbst halten noch regieren.

13. Nun sollten wir weiter von der Hochzeit reden, wie es zugehen soll. Der Papst hat in seinen Rechten verboten die heimlichen Gelübde; doch wiederum, wenn es geschehen ist, verbindet er es, damit er so viel Ehen zerrissen, und wieder zusammen geknüpft hat, daß [es] Jammer ist. Also habe ich nun gesagt: Vater und Mutter soll des Kindes mächtig sein. Denn es ist kein größerer Gehorsam auf Erden, denn Vater- und Mutter-Gehorsam, auch keine größere Gewalt. Wo nun nicht ist Vater und Mutter, da soll sein Vetter, und was die nächste Freundschaft ist und Vormünder. Darum habe ich weiter gesagt: Eine Magd oder Knabe soll sich selbst nicht verloben; denn man bisher aus dem ehelichen Leben eine Lächerei und Schimpf gemacht hat, daß sich die jungen Leute geschämt haben, ein ehelich Gemahl zu begehren, damit hat man sie heimlich in Winkel getrieben.

14. So sollte es nun sein: Wenn sich gleich eines verlobte hinter der Eltern Willen, soll es bei den Eltern stehen, ob es gelte oder nicht. Du darfst nicht darum gen Rom ziehen; es soll's thun Vater und Mutter, wo es ihnen gefällt, daß sie es bestätigen; wo nicht, daß sie es nur frisch zerreißen, daß es nicht binde. Wenn das geschähe, so bliebe viel Irrthums und Hurerei nach, da sonst eitel Unglück und Herzeleid nachfolget.
\quellseite{200}
. das Evangelium mit seinen Gütern und Gaben anzubieten und auszutheilen.

23. Weiter, Laban, der Braut Bruder, heißt auf deutsch weiß, gleichend, das hell daher glänzt und gleist, bei dem findet man die Braut. Das ist die Synagoga und Judenthum, das ging fein gezieret und geschmückt her, allerdings blank und gleichend. Denn das Evangelium trifft nicht die gar ruchlosen Geister, sondern die sich geübt haben, wollten gerne fromm sein; da findet es gemeiniglich arme, elende Gewissen, unter denen findet es die Braut, [die] ist gefangen unter dem Gesetze, die muß es frei machen und zu Christo bringen.

24. Das bedeutet auch, daß sie heraus geht mit dem Kruge aus des Vaters Hause, Wasser zu holen. Zum ersten ist sie Hausdirne, nicht die Frau; das ist das jüdische Volk, das sich getragen hat mit dem Gesetzkrug, das ist ein Wasser, das nur die Thiere tränkt, den äußerlichen Menschen, die nichts vom Glauben wissen; trägt ihn auf dem Rücken, das ist die Bürde, die die tragen, so mit dem Gesetze umgehen.

25. Doch dient das Gesetz dazu, daß das Volk fein im Zaume gehalten werde, daß es ein fein, züchtig Leben führe, gleichwie hier Rebekka, die nun fort eine Frau soll werden. So ist das Predigtamt nichts denn ein Wasserkrug, daraus man die Kameele tränkt; die Kameele aber sind die großen, hoffärtigen Heiligen, dazu müssen die Apostel auch des Wassers trinken. So haben wir zwei Predigtämter, Gesetz und Evangelium. Das Gesetz ist Gleisnerei, da kommt das Evangelium, holt uns heraus zum Herrn, der uns zur Frau macht.

26. Dazu stimmt auch der Braut Name Rebekka, das ist, Gemästete, oder eine feiste Plunisch\footnote{So ist dies Wort in der Wittenberger und in der Jenaer geschrieben. Erlanger: Plünisch.} und dicke Proßel, das ist die Synagoga, gemästet und feist mit großen Werken, muß sich aber danach durch das Evangelium einziehen und schlank machen, ist noch fett und voll eigener Werke, damit sie wollen den Himmel pochen, wie der fette Eglon [Richt. 3, 17.].

27. Das ist die Bedeutung, dadurch vorgebildet ist das Predigtamt und Evangelium, was es ausrichtet in der Welt, wie Christus regiert, und läßt die Apostel seine Braut herzuholen, die sie nicht missen, wer sie sein solle; die nimmt er darnach zu sich; ihre Mutter und Freunde läßt er im Lande, und gibt ihnen Geschenke, der Braut aber gibt er es gar, setzt sie in alle Güter, die er hat. Die Christenheit ist zusammen gelesen aus Juden und Heiden,\footnote{„und Heiden“ fehlt in der Erlanger.} ist zu Christo gekommen gar mit einander; aber die Synagoga hat nur ein Stück und Partheke, ein geringes Kleinod davon.

28. Daß aber Abraham dem Knechte verbeut, er soll je seinem Sohne kein Weib nehmen denn von seiner Freundschaft, nicht von den Cananitern, bedeutet, daß Christus nicht verheißen ist denn den Juden. Da mußte die Braut hergeholt werden; wir sind aber hinzu gekommen, gleich wie die Dirne, so Rebekka mitgeführt hat. Aber das Haupt und Beste der Christenheit ist das jüdische Volk gewesen, das ist das Geblüt. Christus ist ihr Vetter, kommt von den Vätern her; die Apostel sind auch eigentlich zu den Juden geschickt. Da sie aber nicht wollten hören, kehrten sie sich zu den Heiden, wie die Rebekka ihre Amme und Mägde mit sich nimmt, und was sie angehört. Das ist die Historia mit der Bedeutung; nun wollen wir Abraham begraben, und seine Legende ausmachen.

\section{Das fünfundzwanzigste Capitel}

V. 1—10. Abraham nahm wieder ein Weib, die hieß Ketura, die gebar ihm Simron und Jaksan, Medan und Midian, Jesbak und Suah, Jaksan aber zeugete Seba und Dedan. Die Kinder aber von Dedan waren: Assurim, Letusim und Leumim. Die Kinder Midian waren: Epha, Epher, Hanoch, Abida und Eldaah. Diese alle sind Kinder der Ketura. Abraham aber gab alle sein Gut Isaak, aber den Kindern, die er von den Kebsweibern hatte, gab er Geschenke, und ließ sie von seinem Sohne Isaak ziehen, weil er noch lebte, gegen den Aufgang in das Morgenland.
\quellseite{201}
eine reine Jungfrau Christo zubrachte.“ Christus hat eine Braut, die will er, daß sie eine Jungfrau und rein bleibe, wie er. Aber rein wäre sie nicht, hätte er sie nicht rein gemacht; wie abermal St. Paulus Eph. 5, 25—27. sagt: „Er hat sich selber für sie gegeben, auf daß er sie heiligte, und hat sie gereiniget durch das Wasserbad im Wort; auf daß er ihm darstellete eine herrliche Gemeinde, die nicht Flecken oder Runzeln habe, sondern daß sie heilig sei und unsträflich.“ Er hat sich selber für sie gegeben, auf daß er sie heiligte, und hat sie gereiniget durch das Wasserbad im Wort; auf daß er ihm darstellete eine herrliche Gemeinde, die nicht Flecken oder Runzeln habe, sondern daß sie heilig sei und unsträflich.“ Er hat das Wort hinzu gethan, das Evangelium gepredigt, und dadurch sie gereinigt, als sie hat geglaubt. So ist die Christenheit seine Braut, die den Namen von ihm hat, daß sie heißt, wie er heißt, und was der Bräutigam hat, ist alles auch ihr.

20. Wo kommt aber die Braut her, wer hat sie geholt und zugebracht, daß\footnote{Erlanger: da.} sie zusammen kommen? Der Vater mußte einen Knecht mit Kameelen ausschicken und sie holen lassen. Die Braut sucht den Bräutigam nicht, ist auch nicht so der Gebrauch, sondern er muß das erste Wort haben, und um sie werben. Damit ist auch bedeutet, daß nicht in uns steht, anzufahen ein einig gut Werk. So sendet der himmlische Vater seinen Boten, läßt um sie werben, daß wir\footnote{Erlanger: sie.} nicht rühmen dürfen, wir sind die ersten gewesen, und haben es angefangen; er will den Ruhm behalten, daß er das erste Werk thue durch sein Wort. Die Knechte aber und Boten sind die heiligen Propheten und Apostel, sind von Gott in die Welt geschickt, daß sie sein Wort führen, und die Christenheit zu ihm bringen. Wie schickt er sie aber? Also, daß sie nicht wissen, wohin er sie schickt. [Er] nennt wohl das Land, und die Stadt, und Freundschaft; welches aber die Magd sein soll, sagt er nicht, bis daß sie hinkommen; da tritt sie ihm unter die Augen unversehens.

21. Damit ist alles gedeutet das Wesen des Evangelii, wie es geht in der Welt. Sie ziehen nur hin und lassen das Wort gehen, lassen Gott dafür sorgen, der kann sie wohl finden und unter Augen schicken, daß sie da ist, ehe man sich umsieht. Ehe er in das Haus kommt, geht die Braut heraus und begegnet ihm. Also, wenn sie nur das Amt gehen lassen, predigen das Wort, so kommt die Rebekka, die es sein soll, nimmt es auf und glaubt, wird also die Braut. Da gibt ihr nun der Knecht seines Herrn Güter. Das legt Jakobus aus 1 Cor. 4, 1.: „Dafür halte uns jedermann, nämlich für Christi Diener und Haushalter des Geheimnisses Gottes.“ Zum ersten gibt ihr der Knecht sonderliche Kleinode, wie man den Bräuten gibt, inaures oder Stirnspangen von einem halben Seckel Goldes, das ist, fast von einem ungarischen Gulden; darnach auch ein Armgeschmeide, das sie an die Hände und Arme legt, zehn Seckel Goldes, das sind vier und zwanzig Gulden; die zwei drückt er sonderlich aus. Das erste Kleinod an der Stirne ist nicht köstlich, das andere ist noch zwanzigmal so köstlich an [den] Händen. Die geistlichen Güter nun, die er mit sich bringt, und Geschenke, so er gibt der Braut, Mutter und Freunden, das sind die Geheimnisse Gottes, wie sie Paulus heißt, und kurzum, nichts denn das Evangelium; das bringt mit sich Stirnspangen und Armgeschmeide, das ist, es bringt mit sich den Glauben, und darnach die Gaben des Geistes.

22. So werden die zwei angezogen in der Offenbarung Johannis [Cap. 14, 9.], die Stirne und Hände, daß der Endchrist verboten hat, daß niemand kaufen noch verkaufen sollte, er hätte denn ein Malzeichen, Siegel, und Character der Stirne und Hände, gleichwie hier. Das sind die zwei Stücke, Glaube und des Heiligen Geistes Gaben. Die Hände bedeuten das äußerliche Leben, das ein Christenmensch führt; die Stirne aber ist das Angesicht vor Gott. Der Glaube scheint klein [zu] sein, ist aber viel edler und besser; die Liebe scheint viel größer, wie sie auch ist, darum sie sich weit ausbreitet und jedermann dient. Denn, der Glaube allein gegen Gott handelt, und geht auf Eine Person; Liebe aber geht auf alle Menschen. Summa, so deuten hier die Kleinode mannigfaltige Gaben Gottes, wie sie Paulus nennt und erzählt 1 Cor. 12, 4. und auch\footnote{„auch“ fehlt in der Jenaer.} Röm. 12, 6., so der Christenheit gegeben sind, wenn sie gläubig ist worden. Der Glaube muß aber vor sein, darnach kommt denn Gott mit allerlei seinen schönen Gaben, zieret sie mit denselben, als, Weissagen, Regieren, Weisheit, Verstand, Geister unterscheiden etc. Solches alles ist nun gegangen zu der Apostel Zeit, so da in das jüdische Land gesandt sind,
\quellseite{202}
warum er es gethan hätte, sprach er [Richt. 15, 11.]: „Wie sie mir gethan haben, so habe ich ihnen wieder gethan.“ Darum, wenn da Vernunft sollte richten, schließet sie, es sei ein unchristlich Werk. So thut GOtt immerdar, und noch, und wird es wohl immerdar thun. Wir sollten ihn lassen schließen, so wollen wir ihn meistern.

7. Also muß auch der gute Abraham in seinen letzten Tagen ein Weib nehmen, allein darum, daß er uns zu Narren mache, daß die Hure, die Vernunft, still schweige und lasse ihr Urtheil stehen. Also auch, das Simson thut aus dem Heiligen Geist, kann ein anderer nicht thun ohne den Heiligen Geist. Es ist einerlei Werk, aber der Geist scheidet es, wie auch Kains und Abels. Denn Simsons Gemüth ist so gestanden: Es lag ihm nicht daran, daß ihm das Weib genommen war; er war aber gesetzt zum Schußherrn und Richter über das Volk, dazu hatte ihn GOtt geordnet; darum war er schuldig seines Amtes halben, daß er den Feinden Leide thäte, und suchte nur Ursache, wie er konnte, daß er die Philister unterdrückte und sein Volk schützte und rettete. Derhalben nahm er nun diese Ursache an, daß sie ihm mit Frevel hatten sein Weib genommen, daß sie es auch mußten selbst sagen, sie haben es verdient. Also lerne solche Historien und Exempel recht ansehen.

8. Was nun die Kinder Keturä sind, daraus groß Volk ist geworden, fast alle in Arabia, wie man findet in [den] Propheten, haben beide Arabien innen, Reich-Arabia, und Stein-Arabia, lasse ich gehen, es gehört nicht für die Laien. Da nun die Kinder gereifet haben, ist Abraham gestorben, als er gelebt hatte hundert und fünf und siebenzig Jahre. [Er] ist über hundert Jahre ein Pilgrim und Fremdling gewesen in dem Lande, denn er eben im fünf und siebenzigsten Jahre seines Alters ausgezogen war; so lange ist er da ein Fremdling gewesen, hat kein eigen Gut gehabt. Was wird ihm in den hundert Jahren Unglücks unter Augen gestoßen sein, daß sein Herz manchen harten Puff hat müssen [aus]halten; [er] ist nimmer sicher gewesen, immer hin und her geworfen, wie ein Blatt vom Baume; zuletzt kriegt er dennoch noch ein eigen Begräbniß. Es ist kurz beschrieben und gesagt, hundert Jahre; es ist aber lange zu leben.

9. Was thut er mehr? [Er] macht eine Ordnung beim Leben, daß Isaak, der Sohn von der rechten Hauptfrau, im Hause bleiben soll, und in den Gütern und Erbe sitzen, die andern setzt er aus, und gibt ihnen ihr Theil mit.

10. Also ist er gestorben in einem geruhigen Alter; er war des Lebens voll und satt, spricht der Text. Und ist auch ein großer Preis der lieben Väter, und eine edle Gabe von GOtt, daß sie des Lebens müde und satt sind gewesen, und des Todes begehrt haben; [sie] sind gerne und fein ruhiglich\footnote{In den alten Ausgaben: rüglich.} gestorben. GOtt hat sie so wohl gemartert, daß sie wohl des Lebens genug haben gehabt, nach dem Tode verlangt, und sind froh worden.

11. So haben nun ihn die zween Söhne, Ismael und Isaak, begraben. Daraus zu merken ist, daß der Sohn Ismael nicht also vertrieben ist, als sollte er nimmer wiederkommen, weil er je hier wieder zum Begräbniß kommt. Darum ist es nicht unglaublich, doch nicht gewiß, wie etliche sagen, Keturä sei die Hagar gewesen, die habe er nach Sara Tode wieder zu sich genommen. Das ist aber gewiß, daß er den Ismael wieder\footnote{„wieder“ fehlt in der Erlanger.} zu sich genommen hat, doch also, daß er nicht durfte im Hause wohnen; so ist ja die Mutter freilich auch oft bei ihm gewesen, daß er also geblieben ist in Bekentnisse seines Vaters.

12. Weil nun Abraham todt ist, werden wir fortan vom andern Vater, Isaak, hören eine andere Legende, wie er nach des Vaters Tode gesegnet ist. Aber Abrahams Legende ist fast die größte und edelste in diesem Buche, darinne wir kein Stück sehen von dem Gaukelwerke, das wir getrieben haben mit unsern Legenden und Werken, Beten, Fasten, Wachen etc. Man findet nichts davon, was er habe gegessen oder getrunken, wie er sich in die Zeit gestellt habe. Wenn etwas davon [da] stünde, wie sollten sie sich brüsten mit ihren Werken und geistlichem Leben; nun aber nichts da steht denn Glaube, daß GOtt immer sein Wort treibt und übt mit Verfolgung, und sein Wesen immer geht in der Liebe und heiligen Kreuze, so wird die Welt zum Narren über GOttes Worte und GOttes Heiligen. Aber siehe darauf, und lerne daraus, daß kein Leben gelte, es gehe denn auch also. Solches werden wir nun weiter sehen und hören in folgenden Historien von Isaak und Jakob.
\quellseite{203}
V. 11. Und nach dem Tode Abrahams segnete Gott Isaak, seinen Sohn, und er wohnte bei dem Brunnen des Lebendigen und Sehenden.

13. Gott hat den Isaak gesegnet (sagt Moses), daß man sehe, daß die erste Geburt Ismaels, seines Bruders, vor GOtt nicht gilt, ja, er wird verworfen, und kein Segen Isaak gegeben, welcher doch vor der Welt dem Erstgebornen gebührte, zu einem Zeichen, daß des Fleisches oder leibliche Geburt nicht vor GOtt Abrahams Kinder macht, wie das St. Paulus in seinen Episteln [Röm. 9, 8. Gal. 4, 28.] erklärt. Denn sollte hier des Fleisches Ansehen gelten, so wäre es ja recht und billig, daß dieser leibliche und erstgeborne Sohn, Ismael, den Vorzug und Segen behalten hätte, und nicht ohne Ursache und Schuld, wie es die Vernunft ansieht, ihm entzogen wäre. Weil nun dieser nicht angenommen wird, wie viel weniger sollten die andern Juden, die nicht so nahe Abrahams Kinder sind, und sich doch des Fleisches rühmen, vor GOtt gelten.

12—18. Dies ist das Geschlecht Ismaels, Abrahams Sohnes, den ihm Hagar gebar, die Magd Sarä aus Egypten, und das sind die Namen der Kinder Ismael, davon ihre Geschlechter genannt sind. Der erstgeborne Sohn Ismaels, Nebajoth, Kedar, Adbeel, Mibsam, Misma, Duma, Masa, Hadar, Thema, Jetur, Naphis und Kedma. Dies sind die Kinder Ismael mit ihren Namen in ihren Höfen und Städten, zwölf Landesfürsten. Und das ist das Alter Ismaels, hundert und sieben und dreißig Jahr; und ward krank und starb, und ward gesammelt zu seinem Volk. Er wohnte aber von Hevila an bis gen Sur gegen Egypten, denn man gen Assyrien geht, und überfiel alle seine Brüder.

14. Diese Ordnung hält die Schrift immerdar, wie wir gesehen haben, von Anfange an, daß sie die, so von und nach dem Fleische geboren sind, zum ersten beschreibt, und jetzt hernach die aus dem Geiste sind gekommen, darum, daß wir alle, wie St. Paulus sagt, erstlich natürlich vom Fleische müssen geboren werden, darnach durch das Wort geistlich. [Es] ist auch darinne die heimliche Deutung, daß das Gesetz und Gesetzes Volk muß vorgehen, darnach das Evangelium. Gnade wird nicht gegeben, es sei denn zuvor Zorn durch Erkenntniß der Sünde gefühlt. Der alte Mensch muß zuvor hingerichtet werden und sterben, darnach kommt der neue und lebt.

15. Zum andern ist hier angezeigt, wie wir oben [Cap. 4, 17.] gehört, daß Kain mit seinen Kindern am ersten das Land besitzt, und allerlei Hantierung erfindet, und sich immer so fort mehrt und ausbreitet; also auch hier: dieser Ismael setzt sich nieder, denkt im Lande zu bleiben und [zu] regieren, zeugt zwölf Söhne, die macht er alle zu Landesfürsten, wird ein gewaltiger, mächtiger Herr, [es] geht ihm glücklich an nach seinem Vornehmen, überkommt eher zwölf Fürsten, denn Isaak zween Söhne. So [ver]fährt Gott, mehrt ihn so reichlich, als wollte er den Isaak nicht mehren, den er doch gesegnet hatte, wirft diesen empor im Lande, daß er stille sitzt und regiert, und läßt Isaak im Lande irregehen als einen Fremdling, welchem er durch den Vater verheißen hatte, das Land zu geben. Also treibt es Gott mit den Seinen, kartet das Spiel so seltsam, daß es scheint, als sollte nur das Widerspiel geschehen; läßt es dem Glauben also gehen, auf daß er geübt werde, rein und frei bleibe, und stark werde. Jetzt läßt er Ismael obliegen und herrschen; aber nach Isaaks Tode machte er seine Kinder so groß und stark, daß sie alles unter sich brachten.

16. Zum dritten will Moses anzeigen, was Ismael für Wesen, Regiment und Handel geführt habe, als er spricht: Und er überfiel alle seine Brüder, das ist, daß er, wie oben [Cap. 16, 12.] von ihm gesagt ist, ein wilder Mensch worden ist, der niemand gehorsam noch unterthan hat wollen sein, und sich nicht allein wider Isaak, sondern auch an die andern Brüder legt; wie zuvor auch im sechsten Capitel [V. 8.] von den Nephilim oder Tyrannen, und im zehnten\footnote{Wittenberger und Jenaer: elften.} von Nimrod geschrieben ist. Also hat er mit dem Schwerdte gelebet, damit die Leute überältigt, unter sich gebracht und gewonnen. Hat auch ein groß Land unter sich gehabt, fast das ganze Arabien; daher die Juden noch auf den heutigen Tag ein Theil des Landes, so man nennt Arabiam Petraeam, das ist, Stein-Arabien, nennen Kedar, nach dem einen Sohne Ismaels.

17. Siehe, so sind alle Kinder vom Fleisch, das ist ihr Wesen und Handel, haben groß Glück,
\quellseite{204}
Reichthum und Gewalt, sonderlich die Werktheiligen, die den Namen haben, daß sie Gottes Volk heißen. Dagegen sind die rechten Gottes Kinder im Geiste und Glauben, müssen verachtet und für nichts angesehen sein, dazu von jenen überwältigt und verfolgt werden. Ist nun die Deutung und Figur, daß dies eigentlich der Werktheiligen Art ist, die im Gesetze leben, und Frömmigkeit in Werken suchen, daß sie sich immerdar wider die zu Felde legen, und\footnote{„und“ fehlt in der Erlanger.} verfolgen, so da rechtschaffen im Glauben gehen; wie es St. Paulus ausspricht Gal. 4, 29. Das ist Ismael. Folgt nun die Historie vom andern Sohne:

V. 19—23. Dies ist das Geschlecht Isaaks, Abrahams Sohnes. Abraham zeugete Isaak. Isaak aber war vierzig Jahr alt, da er Rebekka zum Weibe nahm, die Tochter Bethuel, des Syrers, von Mesopotamia, Labans, des Syrers, Schwester. Isaak aber bat den Herrn für sein Weib, denn sie war unfruchtbar. Und der Herr ließ sich erbitten, und Rebekka, sein Weib, ward schwanger. Und die Kinder stießen sich mit einander in ihrem Leibe. Da sprach sie: Da mir es also sollte gehen, warum bin ich schwanger worden? Und sie ging hin, den Herrn zu fragen. Und der Herr sprach zu ihr: Zwei Völker sind in deinem Leibe, und zweierlei Leute werden sich scheiden aus deinem Leibe, und ein Volk wird dem andern überlegen sein, und der Größere wird dem Kleinern dienen.

18. Dieser Patriarch hat nicht eine lange Legende, wie Abraham und hernach Jakob. Es gehen aber von Stunde wieder zween Brüder an. Da hebt sich viel Fests\footnote{Da hebt sich viel Fests — da wird großes Wesen gemacht. (Dies s. v. Fest. Vergl. Cap. 26, § 1.)} mit dem Weibe Isaaks, daß [es] vor der Vernunft anzusehen ist, als treibe Gott eitel Narrentheidung. Darum siehe dich vor, daß du nicht narrst mit der Schrift; denn Gott hat es wohl mehr gethan, daß er mit einem geringen Werke alle Welt zu Narren macht. Das erste, das hier steht, ist, daß Isaak versucht wird, der die Verheißung hatte, daß sein Same sollte werden wie Sand am Meere: so stellt sich Gott so närrisch dazu, als werde nichts daraus. Isaak nimmt ein Weib im vierzigsten Jahre, geht dahin fünf, zehn, fünfzehn, bis in das zwanzigste Jahr. Wie viel hat er Hochzeiten gesehen, daß sich’s um ihn täglich gemehrt hatte, denen doch Gott nichts verheißen hatte; sieht dazu seinen Bruder Ismael, daß er so viel Kinder kriegt. Warum nimmt er nicht auch noch ein Weib, wie Abraham that [Cap. 16, 1. ff.], weil er sahe, daß seine Rebekka unfruchtbar war?

19. Das erste, das ihn gestärkt hat, ist gewesen Gottes Wort, daran er fest ist blieben hangen, und Gott hinegestellt [hat], wann und wie er seine Zusagung wollte erfüllen. Zum andern hat er auch das Exempel seines Vaters gehabt, daß Abraham von der unfruchtbaren Sara in seinem Alter, als er sie bei dreißig Jahren oder länger gehabt hat, dennoch zuletzt noch ein Kind kriegte, und nun gelernt hatte, wie Gott seine Verheißung erfüllte wider und über Menschen Vernunft und Verstand. [Er] hat vielleicht auch gesehen, daß die andern Kinder Abrahams von den andern Weibern übel gerathen sind, und ihm nicht wohl mit ihnen ist gegangen.

20. Als nun Isaak eine lange Zeit geharrt hatte, und sich so viele Jahre verzog, hat er zu Gott gebeten, ohne Zweifel nicht ein Mal, sondern oft und viel, daß sein Weib möchte fruchtbar werden; da wird er erhört, daß er zwei Kinder für eines kriegt. Da siehst du, was das Gebet der frommen Gläubigen vermag, daß es muß erhört werden, und reichlich erlangen, darum man bittet. So sehen wir in dem allen nichts, denn seine Exempel des Glaubens, wie zuvor in dem Patriarchen Abraham.

21. Da gibt uns auch der Text von schwangern Weibern zu predigen. Als die Rebekka schwanger wird, widerfährt ihr eine sonderliche Noth, daß sie muß leiden, daß sich die Kinder in ihrem Leibe mit einander stoßen. Das hat ihr freilich nicht wohlgethan, daß sie sich hat müssen erwägen ihres Lebens, und gefürchtet, sie müßte mit den Kindern bleiben. Darum sie auch klagt und schrieet: „Da mir es also sollte gehen, warum bin ich schwanger worden?“ Solche Noth und wunderlich Wesen führt sie so hart, daß sie hingeht, den Herrn zu fragen; sonst, wo sie nicht in Furcht des Todes gewesen wäre, wäre sie nicht gegangen zu fragen. Wo ist sie denn hingegangen zum Herrn? Sie ist ja gen Himmel nicht gefahren.

22. Die Juden sagen, sie sei zu dem Altvater Sem, dem Sohne Noah, gegangen, und habe
\quellseite{205}
denjenigen gefragt, welcher nach der Rechnung zu der Zeit noch gelebet hat, und hernach bei fünfzig Jahren. [Es] ist auch wohl glaublich. Denn er ist der älteste Vater gewesen, hat lange Zeit gelebt, wohl erfahren und geübt in geistlichen Sachen, und voll Heiliges Geistes, daß er hat wissen in solchen Sachen zu rathen. Auch sind noch neben ihm etliche mehr gewesen desselbigen Geschlechtes, auch heilige Väter. Zu solchen ist sie nun gegangen, ohne Zweifel nicht ohne Befehl, sondern mit Rath ihres Mannes, und [hat] gefragt, was doch das werden möchte, oder was sie daraus gedenken sollte.

23. Da kriegt sie solche Antwort: Laß dich es nicht wundern, daß dir es also geht, denn „du hast zwei Völker in deinem Leibe“ \&c. „und der Größere wird dem Kleinen dienen“. Das scheint, als sei es ein zauberisch Wort, als unter den Heiden solches viel war; aber das ist ein sonderlich Werk Gottes, daß er sagt, der Größte soll dem Kleinsten dienen. Das ist es, daß wir so oft gesagt haben, und hernach viel mehr sehen werden, ist auch beschlossen im Himmel, daß [es] immerdar also gehen muß, wie mit den zwei Kindern, daß eines muß oben liegen, das andere sich unter die Füße treten lassen. In geistlichen Sachen, sage ich, nicht in äußerlichem Wesen und Gütern; also, daß in dem Volk, welches sich rühmt Gottes, seines Wortes und Gesetzes, muß [es] sich also theilen in die zween Haufen, und welche die Kleinsten und Geringsten sind, müssen hervor gezogen und empor gehoben werden, die Größten aber und Mächtigsten werden verworfen. Das ist so viel geredet: Gottes Volk gewinnt, wenn es gleich unten liegt, und der andere Haufe, wenn er am höchsten und festesten sitzt, so fällt er und hat verloren. Wie ist denn das möglich, daß der gewinnen soll, der da unten liegt? Vor Gottes Augen muß man es ansehen; da heißt es also: Vor der Welt verloren, vor Gott gewonnen; und wiederum.

24. Also ist dieser Spruch und Figur erfüllt zu der Apostel Zeit, so lange das Evangelium gepredigt ist, als sie von Juden und Heiden verfolgt und vertrieben wurden; da kann Vernunft und Welt nicht anders sagen: Bei den Juden gewonnen, bei den Aposteln verloren. Was geschah aber? In dem, daß die Juden gewinnen, liegen sie unter, und jene, die verlieren, behalten den Sieg. Wie so? Denn die Juden werden eben durch der Apostel Wort, um Christi Blutes willen, niedergelegt, zerstöret und verheeret, also, daß sie noch verachtet sind, und stinken in aller Welt; die Apostel aber und Märtyrer sind in höchsten Ehren vor Gott, und gepriesen über alle Welt, daß nie kein König noch Kaiser solche ewige Ehre und Ruhm auf Erden erlangt hat.

25. Also wird es zu unsern Zeiten auch gehen, so wirs werth sind. Der Papist mit seinen Bischöfen und Tyrannen werden feindlich rühmen und gewonnen schreien, wenn sie uns, die wir das Evangelium recht haben, genug verfolgt, gemartert, verbrannt und umgebracht haben. Aber, wenn man meint, es sei aus mit uns, so wird Gott offenbar machen, wer gewonnen habe, und sie gestürzt werden und zu Scheitern gehen, wir aber ewiglich bleiben. Summa, das ist Gottes Art und Werk, daß er immerdar treibt: Wer unterliegt, der liegt oben. In weltlichen Sachen geht es anders zu, da spricht man wohl: Gott helfe dem Stärksten; hier aber hilft Gott dem Schwächsten.

26. Rebekka, der zwei Kinder Mutter, bedeutet die christliche Kirche oder heilige Schrift. Denn also nennt Gott die Schrift oder sein Wort seinen\footnote{In den Ausgaben: seiner.} Mutterleib und seine Brüste, daraus seine Christen geboren werden, wie ein Kind aus Mutterleibe geboren wird und darinne ernährt; wie Jesaias [Cap. 46, 3.] spricht: „Höre mich, du Haus Jakob, und alle Uebrigen von Israel, die ihr getragen werdet in meinem Bauche, und von meiner Mutter.“ Darinne hebt sich nun der Hader, daß sich die zwei mit einander stoßen; das muß sie fühlen und leiden, daß sie meint, sie müsse mit den Kindern bleiben, das ist, daß sich die zwei Völker wider einander auflegen. Und hebt sich der Hader um der Schrift, welche die rechten Christen für sich haben; die Werktheiligen aber wollen sie mit Gewalt auf ihr Ding reißen, und verfolgen jene darüber, daß es scheint, und die Vernunft sagen muß, sie müssen zu Trümmern gehen.

27. Darum schließest sie also: es sei nicht fein, daß solche Zwietracht in der Welt sei; man sollte doch ein gemein Concilium machen, daß man der Sache eins würde, und ein jeglicher wüßte, weß er sich halten sollte. Wie denn jetzt meinen die Allerbesten und Klügsten in der Welt. Ja, harre
\quellseite{206}
so lange, bis das geschieht, und fahre dieweil zum Teufel. Der Papst hat es auch gethan, wollte solche Zwietracht aufheben, und ein einig Regiment machen, auf daß der Hader gestillt würde, daß sich die Kinder nicht mehr schlügen in der Mutter Bauche. Was hat er aber ausgerichtet? Nichts denn Teufel und Tod.

28. Gott aber sagt also dazu: Es wird nichts daraus, [es] wäre auch nicht recht noch gut, daß es stille sei; ich will es also haben, daß von den Kindern zwei Völker kommen sollen, die ewiglich müssen wider einander sein; darum sagt St. Paulus [1 Cor. 1, 19].: „Es müssen Spaltungen sein.“ Aber der stärkste und größte Haufe sind die Jesuiten, die mit Werken umgehen und ohne Glauben sind; aber Gott leidet endlich nicht, daß sie überlegen sind, sondern kehrt es um, und macht, daß Jakob den Esau unter die Füße tritt.

29. Also geht es jetzt zu unsern Zeiten auch. Nachdem das Evangelium offenbar und am Tage ist, müssen wir uns wider das Papstthum legen, und sagen, daß [es] eitel Teufelslehre ist, wie es Paulus [1 Tim. 4, 1.] heißt. Wir können nicht schweigen, oder müßten Gott zum Feinde haben; so können sie es nicht leiden, haben Gewalt und das Schwert, und wollen uns todt haben. Darum ist es eben so viel, wenn man Frieden will haben in der Christenheit, als hätte man keine Christenheit.

30. Also schließe nun: Wo das Evangelium angeht, da hebt sich der Hader, daß man in die Schrift kommt, und sich muß schlagen auf allen Seiten mit den Rotten; wo das nicht ist, so ist es ein gewisses Zeichen, daß da keine Christen sind. Weil Rebekka gute Tage hat, trägt sie keine Kinder.

31. Solche Zwietracht aber richtet Gott nur darum an, daß der kleine Haufe empor komme und obliege, daß man sehe, wie alle Gewalt und Stärke muß gewonnen geben, und nichts vermag wider Gottes Wort und die Seinen, wie gering und schwach sie sind.

32. Nun, diesen Spruch: „Der Größte soll dem Kleinen dienen“, hat St. Paulus angezogen und wohl ausgelegt Röm. 9, 10—13.: „Da Rebekka von Isaac schwanger war (spricht er), ehe die Kinder geboren waren, und weder Gutes noch Böses gethan hatten, auf daß der Vorsatz Gottes bestünde nach der Wahl, ward zu ihr gesagt, nicht aus Verdienst der Werke, sondern aus Gnaden: Der Größte soll dienstbar sein dem Kleinen; wie geschrieben stehet: Jakob habe ich geliebet, Esau habe ich gehasset.“ [Er] will also schließen, daß sich Jakob nichts könne rühmen, denn deß: Was ich habe, das ist mir beschert und geschenkt von Gott, ehe denn ich jung ward in Mutterleibe. Also, daß dieser Text mit aller Gewalt niederlegt den freien Willen und alle Verdienste. Denn, was könnten sie thun, oder womit sollten sie es verdienen, daß dieser obliegen, oder jener unterliegen soll?

33. Also haben wir das göttliche Urtheil von den zwei Brüdern, daß, was da hoch ist, vor Gott nichts sei, und müssen herunter geworfen werden; und was niedrig ist, vor ihm hoch ist und obliegen müsse, auf daß da gewehret werde beide, aller Vermessenheit und Verzweiflung, und die Armen Trost und Trost haben, aber die Großen und Gewaltigen sich fürchten und demüthigen.

Da nun die Zeit kam, daß sie gebären sollte, siehe, da waren Zwillinge in ihrem Leibe. Der erste, der heraus kam, war röthlich, ganz rauh, wie ein Fell, und sie nannten ihn Esau. Danach darnach kam heraus sein Bruder, der hielt mit seiner Hand die Ferse des Esau, und hießen ihn Jakob. Sechzig Jahr alt war Isaac, da sie geboren wurden. Und da nun die Knaben groß wurden, ward Esau ein Jäger und ein Ackermann, Jakob aber ein einfältiger Mann, und wohnte in den Hütten. Und Isaac hatte Esau lieb, darum, daß er von seinem Weidwerk aß; Rebekka aber hatte Jakob lieb.

34. Die Schrift will das Wunderwerk anzeigen, daß die Kinder in Mutterleibe getrieben haben, und darnach, wie jeglicher gerne der erste heraus gewesen wäre. Der erste, sagt Moses, ist röthlich gewesen, nicht, daß er ganz roth war, sondern sein schwarzbräunlicht, hart und mager; wie man auch liest von David [1 Sam. 16, 12.], daß er ein Bräunling war, mit hübschen Augen \&c. \unsicher{unleserlich} Esau aber heißt auf deutsch, thätig oder schätzig und mager, als wir auf unsere Weise sagen, Thäterich,\footnote{So die Erlanger. Wittenberger: Ackerreich; Jenair: Ackerich. Wir haben der ersteren Lesart den Vorzug gegeben, weil Luther (Wald, alte Ausg., Bd. XIV, 1292, § 8) dem Namen „Thätig“ die Auslegung gibt: „reich in Gott, heilig, göttlich“.} darum, daß er vor der Welt sollte ein Held werden.
\quellseite{207}

35. Jakob aber hat den Namen von der Ferse oder Fußsohle, als der einen mit den Fersen oder Füßen tritt, darum, daß er seinen Bruder hier bei der Ferse hält, und hernach im 27. Capitel, V. 36., Esau selbst sagt: Er heißt wohl Jakob, denn er hat mich nun zweimal untergetreten. Dieser, spricht der Text, war ein einfältiger Mann, und wohnete in den Hütten.

36. Einfältig, auf hebräisch [תם] Tham, heißt so viel, als wir sagen, ein unschuldiger, frommer Mann, ohne Arg, der niemand Schaden noch Leid thut. Dieser hat gewohnt in den Hütten, das ist, er ist daheim bei der Mutter geblieben, und [hat] sich nicht hervor gebrochen. Esau aber wird ein Ackermann und ein Jäger, schickt sich zur Nahrung, als dächte er ewig zu leben, wird rüstig und reisig zu Sachen.

37. Da siehe nun abermal Gottes Gericht, also, daß auch Isaak betrogen wird, und sein Datum setzt auf den Sohn Esau, und hat ihn zum liebsten, darum, sagt der Text, daß er von seinem Weidwerke aß. [Er] hat also gedacht: Rebekka ist eine Närrin, hält von dem Jakob viel; aber dieser Sohn, meint er, sollte es thun; Jakob aber mußte der Aschenbrödel sein, ohne daß ihn die Rebekka lieb hatte. Daher kannst du schließen, daß dieser Esau gar ein köstlicher Mann gewesen ist, der daher gegangen ist in höchstem Schein und seinem ehrbaren Wesen, daß es dem Vater ganz wohl gefallen hat, daß er schlecht meinte, er wäre der rechte Sohn, auf den der Segen und Verheißung lautete. [Es] betrog ihn auch, daß er der Erstgeborne war, wie der Text sagt; von Jakob aber hält er nichts, weil er so schlecht und einfältig daher geht.

38. Aber Gott läßt es eine Weile geschehen, daß der gute Vater also irrt, endlich aber kehret er es um, und macht es allezeit anders, denn wir es vorschlagen; denn er will schlecht, daß man ihn fürchte. Wer das nicht thut, sondern läßt ihm seinen Sinn und Dünkel wohlgefallen, so stürzt er es gewißlich um. Die Rebekka hat also gedacht: Ach daß Gott meinen Sohn auch annähme, daß er möchte gesegnet werden! denn sie hat ja den Spruch für sich gehabt: „Der Größte soll dem Kleinen dienen“; wiewohl es viel anders vor den Augen schien. Jakob hat auch wohl halb gezweifelt, weil er sah, wie der Vater so viel von seinem Bruder hielt. Also sind sie auf dieser Seite in der Furcht gestanden, Isaak aber und Esau meinten, sie wären der Sache gewiß; darum fehlten sie auch. Wie sich aber Esau in das Wort geschickt hat, das von den zweien Brüdern gesagt war, weiß ich nicht, daß er also irrte und fehlte, so doch die Worte hell und klar waren: Der Größte soll dienen, der Kleine soll Herr sein. [Es] ist aber damit angezeigt, wie Gott seine Heiligen zuweilen läßt weiblich irren, auf daß sich niemand auf Heilige verlasse, noch auf Concilia, und alles, was da groß scheint, sondern nichts sehen noch hören wolle, denn was Gottes Wort sagt.

Da kam Esau vom Felde und war müde, und sprach zu Jakob: Laß mich kosten das rothe Gericht, denn ich bin müde. Daher heißt er Edom. Aber Jakob sprach: Verkaufe mir heute deine Erstgeburt. Esau antwortete: Siehe, ich muß doch sterben, was soll mir denn die Erstgeburt? Jakob sprach: So schwöre mir heute; und er schwur ihm, und verkaufte also Jakob seine Erstgeburt etc.

39. Das Wort [אדם] Edom heißt auf deutsch „roth“, davon kriegt Esau hier den Namen. Droben, als er geboren ward, sagt der Text, wie er am Leibe röthlich war, oder ein Bräunling. Aber davon gibt ihn die Schrift nicht den Namen, sondern um der rothen Suppe willen; damit sie hat wollen die heimliche Deutung anzeigen, wie wir hören werden.

40. Dieser Text zeigt nun an, wie Esau ein stolzer Heiliger gewesen ist, daß er die erste Geburt so geringe achtet, und um die rothe Suppe gibt; hat also gedacht: Ha, ich bin doch der Erste, daß mir alle Herrschaft und Recht der ersten Geburt wohl bleibt. Denn es war zu der Zeit und immerdar bei den Juden ein köstlich Kleinod um die erste Geburt, also, daß sie mit sich brachte beide, das Priesteramt und Regiment, dazu zwei Theile des Erbgutes. Item, daß der Erstgeborne sitzen bliebe in des Vaters liegendem Gute, und der Stamm auf ihn bliebe, daß die Linie des Geschlechts nach ihm geführt ward.

41. Aber das war hier viel größer, daß Christus dem Abraham und seinem Samen verheißen war, darum war die Ehre der ersten Geburt noch höher. Aber weil sie oft auf solch Recht stolzierten und trotzten, wendete es Gott auch um, und verwarft sie, daß die Jüngsten dazu erwählt wurden; als wir auch gesehen haben in Adams und Abrahams Kindern, und hernach von Jakobs und Josephs Söhnen; denn er will
\quellseite{208}
nichts ansehen noch leiden, was nicht in der Furcht geht.

42. Man mag hier Jakob entschuldigen, daß er dem Bruder das rothe Gericht nicht wollte geben, ohne um die erste Geburt, daß er es habe gethan aus Eingeben des Heiligen Geistes; denn ohne Zweifel seine Mutter ihm immerdar die Worte vorgesagt hat: Der Größte soll dem Kleinsten dienen.

43. Daß aber die Sophisten sagen, Esau habe gesündiget, daß er geistliche Güter verkauft, und heißen es Simoniam, lasse ich fahren, es sind ihre Träume. Die Schrift straft das, daß er die erste Geburt verachtete und in Wind schlug, welches doch eine feine Gabe von GOtt war, und gibt sie um den rothen Brei. Er sollte Leib und Leben darüber gelassen haben, ehe er solche Ehre und Freiheit, von Gott gegeben, von sich hätte lassen kommen; weil er es aber so hochmüthig verachtet, verachtet ihn Gott wiederum, und beraubt ihn alles Rechtes und auch des Segens.

44. Darum hat Jakob nicht gesündiget, daß er sie zu sich nimmt, wie die Kinder Israel [2 Mos. 12, 35.] nicht gesündiget haben, als sie von den Egyptern silberne und güldene Geräthe und Kleider nahmen, wie sie Gott geheißen hatte. Denn die Egypter waren Gottes, mit allem, was sie hatten. Also ist auch des Esau erste Geburt Gottes gewesen, daß er sie ihm mochte nehmen und dem Bruder geben. Wie sie nun Jakob überkommen hat, es heiße gestohlen oder gekauft, so gehört sie ihm. Also lehrt dies Exempel, wie Gott die Verwegenheit und Hochmuth straft, und Demuth erhöhet. Nun wollen wir auch sehen, was hierinne bedeutet ist.

45. Wir haben gehört, wie diese zween Brüder zwei Völker bedeuten, die sich Gottes rühmen, und äußerlich nicht geschieden sind von einander; denn sie bleiben beide in einem Hause; sondern inwendig nach dem Glauben des Herzens; das sind die beiden, Juden und Christen. Die Juden waren Esau, thätig und hurtig, die sich frisch übten in Werken des Gesetzes. Da nun das Evangelium aufging, kam der Esau zum ersten hervor, der war fein röthlich; das ist, die Werktheiligen leuchten schön daher in ihren feinen Werken, daß man sie allein für Gottes Kinder hält; und müssen eben röthlicht oder braun sein, das ist, ein heftig, heiß, zornig Volk, wie braune Leute gemeiniglich sind. Das ist es auch, daß sie haarig und rauh sind, wie ein Fell, das ist, einer halsstarrigen und unschlachtigen Art, nicht säuberlich noch freundlich. [Es] ist kein unleidlicher, unsanfter und bitterer Volk, denn solche Werktheiligen; darum sie auch Christus [Matth. 12, 34.] Ottergezichte heißt; [sie] können keine Geduld noch Barmherzigkeit haben mit denen, so aus Schwachheit straucheln; urtheilen jedermann, schurren und purren, sobald man ihnen nicht recht thut.

46. Die rechten Christen aber sind sanft und gelinde, können die Schwachen wohl tragen und bessern. Diese müssen nun, wie gesagt, die Erstgebornen sein, daß sie größere Ehre, Namen und Gewalt vor der Welt haben. Aber Jakob wollte auch gerne am ersten hervor kommen; denn die Christen sehen gerne, daß jedermann das Evangelium erkennt; es wird aber nichts daraus. Esau muß vor der Welt vorgehen; jedoch hält Jakob\footnote{Erlanger: er.} mit der Hand die Ferse Esau, das ist, die Werktheiligen drücken die frommen Christen unter, außen vor der Welt; aber vor Gott kehrt sich's wieder um, daß sie jene unter die Füße treten, wie wir hören werden.

47. Weiter habe ich oben [§ 26] gesagt, daß Rebekka bedeutet die heilige Schrift, Natur aber, der Vater, bedeutet Gott. Wenn wir nun gegen einander halten die Wohltthat, den fleischlichen Juden von Gott widerfahren, und auch den rechten Gläubigen, so scheint es, daß Gott die Juden lieber habe gehabt denn die Apostel; [er] stellt sich gegen ihnen als ein Vater, diese aber heißt er gehen, als jedermann's Knechte, und stellt sich, als kenne er sie nicht. Den Juden hat er viel Landes und Gutes verheißen, die Apostel läßt er Kummer und Noth leiden; jedoch hat Jakob das beste Theil für sich, daß sich die Mutter zu ihm hält. Also hat dies einfältige, fromme Häuflein die Schrift und Gottes Wort für sich, damit sie sich trösten, und daran genügen lassen, fragen nicht viel nach der Welt.

48. Was bedeutet es aber, daß Esau die Erstgeburt dahin gibt um die rothe Suppe? Was in der Schrift von Essen und Trinken steht, soll man auf die Lehre deuten. Das ist nun die, so man getrieben hat im ganzen alten Testamente, von zeitlichen Gütern, wie sie sagen im 144. Psalm, V. 12.: „Daß unsere Söhne aufwachsen, wie die Pflanzen, in ihrer Jugend, und
\quellseite{209}
unsere Töchter, wie die ausgehauenen Erker, gleichwie die Paläste; daß unsere Winkel voll seien, die heraus geben Vorrath, bis wieder zu Vorrath; daß unsere Schafe tragen tausend und hunderttausend draußen;“ solches hat man uns eben auch gelehrt, daß ist, daß bräune Süpplein vorgetragen, danach ihn lüstet. Wenn es wohl zugeht, daß wir keinen Mangel haben, sondern Keller und Boden voll, so spricht man: Hier wohnt Gott; wie sie daselbst, B. 15., sprechen: „Wohl dem Volke, dem es also gehet.“ Der rauhe Esau bedeutet nun dieselbigen Prediger, die vom Glauben nichts wissen, und predigen, was man gerne hört und zu guten Tagen dient.

49. Daß er nun vom Acker kommt und müde worden ist, bedeutet, daß solche Predigt den Puff nicht leidet. Also sagt Paulus Gal. 5, 11.: „Ich aber, so ich die Beschneidung noch predige, warum leide ich denn Verfolgung? So hätte das Aergerniß des Kreuzes aufgehört.“ Wer von Werken predigt, der wird nicht verfolgt. Denn, was die Vernunft begreifen kann, da hat sie Lust zu; was sie aber nicht fassen kann, das will sie nicht leiden. Darum wollen sie nicht daran, vom Glauben zu predigen, wie er allein rechtfertig macht, sondern predigen also, daß auch der Bauch genug habe; daß rothe Süpplein schmeckt ihnen zu wohl, da lassen sie nicht von. Jakob aber läßt das rothe Gericht fahren, gönnt es dem Esau wohl, überkommt aber ein Höheres und Besseres, nämlich die Erstgeburt.

50. Also ging es zu der Apostel Zeiten. Die Juden behielten das rothe Mus, blieben auf Werken und zeitlichem Gute; sie aber ließen äußerlich Ding fahren, und ergriffen die Erstgeburt durch den Glauben, und wurden Könige und Herren der ewigen Güter. Daher heißen die Juden und Werktheiligen noch heutiges Tages wohl Edom, daß sie gerne hören Menschenlehre und Werke. Und wie Esau die Erstgeburt verachtete, und sprach: „Was ist sie mir nütze, ich muß doch sterben“, also thun diese auch, und sprechen: Ach, sollten wir unser Ding und alles zeitliche Gut lassen fahren, das wäre uns zu schwer, wir haben doch nicht mehr davon, weil wir leben, und verlieren doch darüber beide, zeitliche und ewige Güter.

\section{Das sechsundzwanzigste Capitel}

V. 1—33. Es kam aber eine Theurung in das Land über die vorige, so zu Abrahams Zeiten war, und Isaak zog zu Abimelech, der Philister König, gen Gerar. Da erschien ihm der Herr und sprach: Zieh nicht hinab in Egypten, sondern bleibe in dem Lande, das ich dir sage; sei ein Fremdling in diesem Lande, und ich will mit dir sein, und dich segnen; denn dir und deinem Samen will ich alles dies Land geben, und will meinen Eid bestätigen, den ich deinem Vater Abraham geschworen habe. Und will deinen Samen mehren wie die Sterne am Himmel, und will deinem Samen alles dies Land geben, und durch deinen Samen sollen alle Völker gesegnet werden, darum, daß Abraham meiner Stimme gehorsam gewesen ist, und hat gehalten meine Sitten, meine Gebote, meine Aufsätze und meine Gesetze. Also wohnte Isaak zu Gerar. Und wenn die Leute am selben Orte fragten von seinem Weibe, so sprach er: Sie ist meine Schwester, denn er fürchtete sich zu sagen: Sie ist mein Weib, sie möchten mich erwürgen um Rebekka willen, denn sie war schön von Angesicht. Als er nun eine Zeitlang da war, sahe Abimelech, der Philister König, durch das Fenster, und ward gewahr, daß Isaak scherzte mit seinem Weibe Rebekka. Da rief Abimelech dem Isaak und sprach: Siehe, es ist dein Weib, wie hast du denn gesagt, sie ist meine Schwester? Isaak antwortete ihm: Ich gedachte, ich möchte vielleicht sterben müssen um ihretwillen. Abimelech sprach: Warum hast du denn uns das gethan? Es wäre leicht geschehen, daß jemand vom Volke sich zu deinem Weibe gelegt hätte, und hättest also eine Schuld auf uns gebracht. Da gebot Abimelech allem Volke und sprach: Wer diesen Mann oder sein Weib antastet, der soll des Todes sterben. Und Isaak säete in dem Lande, und fand desselben
\quellseite{210}
Jahres hundertfältig.\footnote{In den Ausgaben: „hundert Schäffel“. Doch nachher, wo dieser Text wiederholt wird: „hundertfältig“; nur die Erlanger hat auch da: „Schäffel“.} Denn der HErr segnete ihn, und er ward ein großer Mann, ging und nahm zu, bis er fast groß ward, daß er viel Gutes hatte an Schaafen und Rindern, und ein groß Gesinde. Darum richteten ihn die Philister, und verstopften alle Brunnen, die seines Vaters Knechte gegraben hatten zur Zeit Abrahams, seines Vaters, und fülleten sie mit Erde, daß auch Abimelech zu ihm sprach: Zeuch von uns, denn du bist uns zu mächtig worden. Da zog Isaak von dannen, und schlug sein Gezelt auf im Thale Gerar, und wohnete allda. Und als er sich nun gesetzt hatte, ließ er die Wasserbrunnen graben, die sie zu Abrahams Zeiten, seines Vaters, gegraben hatten, welche die Philister verstopft hatten nach Abrahams Tode, und nannte sie mit dem Namen, da sie sein Vater mit genannt hatte. Auch gruben Isaaks Knechte im Thale, und fanden daselbst einen Brunnen lebendiges Wassers. Aber die Hirten von Gerar zankten mit den Hirten Isaaks, und sprachen: Das Wasser ist unser; da hieß er den Brunnen Esek,\footnote{Handglosse: Esek heißt Hohn, wenn man jemand Gewalt und Unrecht thut.} darum, daß sie ihn da verhöhnet hatten. Da gruben sie einen andern Brunnen, da zankten sie auch über, darum hieß er ihn Sitna.\footnote{Handglosse: Sitna heißt Widerstand, daher der Teufel Satan heißt, ein Widerwärtiger.} Da machte er sich von dannen, und grub einen andern Brunnen, da zankten sie sich nicht über, darum hieß er ihn Rehoboth,\footnote{Handglosse: Rehoboth heißt Raum oder Breite, daß nicht enge ist.} und sprach: Nun hat uns der HErr Raum gemacht, und uns wachsen lassen im Lande. Darnach zog er von dannen gen Bersaba. Und der HErr erschien ihm in derselben Nacht, und sprach: Ich bin deines Vaters Abraham GOtt, fürchte dich nicht, denn ich bin mit dir, und will dich segnen, und deinen Samen mehren um meines Knechtes Abraham willen. Da bauete er einen Altar daselbst, und rief den Namen des HErrn an, und richtete daselbst seine Hütten auf, und seine Knechte gruben daselbst einen Brunnen. Und Abimelech ging zu ihm von Gerar, und Ahusath, sein Freund, und Phichol, sein Feldhauptmann. Aber Isaak sprach zu ihnen: Warum kommt ihr zu mir? habt ihr mich doch, und habet mich von euch getrieben. Sie sprachen: Wir sehen mit sehenden Augen, daß der HErr mit dir ist, darum sprachen wir: Es soll ein Eid zwischen uns und dir sein, und wollen einen Bund mit dir machen, daß du uns keinen Schaden thust; gleichwie wir dich nicht angetastet haben, und wie wir dir nichts, denn alles Gutes gethan haben, und dich mit Frieden ziehen lassen. Du aber bist nun der Gesegnete des HErrn. Da machte er ihnen ein Mahl, und sie aßen und tranken. Und des Morgens frühe stunden sie auf, und schwur einer dem andern. Und Isaak ließ sie gehen, und zogen von ihm mit Frieden. Deselbigen Tages kamen Isaaks Knechte, und sagten ihm an von dem Brunnen, den sie gegraben hatten, und sprachen zu ihm: Wir haben Wasser gefunden. Und er nannte ihn Saba,\footnote{Handglosse: Saba heißt ein Eid oder Schwur, Beer aber heißt ein Brunnen.} daher heißt die Stadt Bersaba, bis auf den heutigen Tag.

1. Da hast du wohl angezeigt, wie Moses ein unnützer Wächer ist, daß er von unnützen Sachen so viel Fests\footnote{Fests = Aufhebens.} macht. Was geht es uns doch an, daß GOtt der Allmächtige solche Dinge also läßt auslaufen in der Welt, daß jedermann muß davon lesen? Aber alle Vernunft muß dazu sagen, es sei eitel Narrentheidung; haben wir sonst nichts zu schaffen noch zu lesen, denn wie Isaak Brunnen gräbt, und mit dem Weibe scherzt? Ist das so eine köstliche Historie, da die Macht an liegt? Da er von heiligen Vätern schreiben wollte, konnte er nicht von guten Werken und Exempeln schreiben, als beten, fasten, fasten etc., und solches Narrenwerk daheim lassen, welches sonderlich den geistlichen Leuten verdrossen zu lesen ist? denn sie sollen nicht von unnützen weltlichen Sachen studiren. Er hat zwar damit Ursach genug gegeben dazu, daß man die Bibel unter die Bank gestoßen hat, und nicht werth geachtet, daß unsere hohen Doctores darinne studiren sollten. Warum schreibet er auch so närrisch Ding?

2. Das erste, so wir oft gehört haben, sehen wir hier auch: Isaak mußte der größten Erzväter einer werden, an dem viel gelegen war vor GOtt. Warum, ob wir es nicht verstanden, was Moses hiermit meint, sollen wir nicht mit der tollen Vernunft zufahren, und schließen, es sei Narrheit. Denn (wie oft gesagt) GOtt hat
\quellseite{211}
Lust dazu, daß er uns zu Narren mache; legt uns eben solche närrischen Werke vor, darum, daß er die großen Heiligen niederschlage, daß man sehe, wie es ihm alles allein am Glauben liegt, und nicht leiden kann, daß man irgend ein Werk rühme, sondern seine bloße Gnade. Wenn es nun seinen Ruhm und Preis hat, so liegt nicht daran, es sei ein Werk wie gering es wolle. Und gilt das allernärrischste Werk so viel, als das allerhöchste und köstlichste. [Er] läßt solche geringen Werke preisen und predigen durch alle Welt, da er der großen Werkheiligen keines sehen noch wissen will.

3. Darum ist das das höchste Stück in diesem Capitel, daß Gott mit dem Manne zweimal redet. Wirke nun hin und her, lebe sonst oder so, wenn Gottes Wort klingt, so thue die Augen, Ohren und Herz auf; wenn die Majestät redet, so hat es Kraft und Macht, weil er nun mit dem Isaac so zu reden pflegt.

4. Die erste ist fast lang, welches eitel treffliche, starke Verheißungen sind, wie zuvor Abraham gegeben, daß sie auch mit großer Gewalt, Schwur und Eide bestätigt sind. [Sie] geht uns aber vor den Ohren über, und weil es uns nicht trifft, schlagen wir es in Wind. Aber wenn wir recht aufheben, welche Worte es sind, und was sie hinter sich haben, würden wir uns davor entsetzen, und groß genug achten.

5. Weil wir nun hören (sage ich), daß Gott hier zweimal redet mit so viel Worten, sollst du gewiß dafür halten, daß ihm an den geringen Werken, Scherzen, Brunnengraben, Essen und Trinken, viel mehr gelegen ist, denn aller Klöster, Mönche und Nonnen, und viel höher zu rühmen denn alle ihr Gebet, Fasten und heiliges Leben. Was ist die Ursache? Keine andere, denn daß hier Gottes Wort steht, daran sich Isaac gehalten hat, unsere Geistlichen aber ohne Wort und Glauben fahren.

6. Mir ist auch also gewesen. Ehe ich ein Narr ward, und noch ein gelehrter Doctor wollte sein, war ich dem Buche allein darum feind, daß Moses so viel Geschwätz machte, und erzählte, wie Gott mit den Vätern redet, als ein Schuster mit dem andern; darum hielt ich von diesem Exempel nicht mehr, denn als wenn man von mir eine Legende machte, und schriebe, wie ich über Feld zöge, oder daheim säße und schliefe.

7. Also schlägt es alle Vernunft in [den] Wind, und sieht nicht darauf, was da für Wunder geschieht, daß Gott selbst redet; das sieht sie nur: O das ist ein loses Werk. Aber wiederum, der hat so lange gefastet, eitel trocken Brod gegessen und Wasser getrunken, sein Lebenslang hären Hemde getragen. St. Hieronymus lag so lange in der Wüste, und hatte einen Stein unter dem Kopfe für ein Kissen, und schlug seine Brust mit Steinen. O das sind große, heilige Werke! Moses aber weiß nichts zu schreiben, denn wie sich Isaac mit der Frau schleppet und mit ihr scherzt und narrt.

8. Darum hat Gott immer damit zu thun, daß er der tollen Hure, der Vernunft, wehre, die die Güte der Werke nach der Größe und Länge messen will, schlägt ihr zuwider und [zum] Trotz solche Werke vor, die sie für nichts hält. Daher schließe nun, daß nichts liegt an vielen, großen, langen Werken, sondern allein an Gottes Wort. Wenn das daran hängt, so halte es für keinen Scherz noch Gaukelfwerk; so bald das klingt und lautet, so ist es eitel köstlich Ding. Also, daß Gott nicht will gerechtteilt haben nach Werken, sondern nach seinem Willen. Die Vernunft spricht wohl: Das Werk ist groß, darum wird es Gott ansehen; denn was groß ist, das soll man auch hoch heben. Verharren ist es nütz und noth, solche Exempel zu beschreiben, daß nicht die Vernunft Gott ermuntere, und ihm sagen wolle, was er thun oder nicht thun sollte. Solches müssen wir oft sagen, weil es so oft angezeigt wird, daß wir sehen, wie die ganze Schrift darauf bringt; wiewohl es dennoch nichts geholfen hat.

9. So merkte nun, daß der Isaac muß ein außerwählter Kern sein, weil ihm die Ehre widerfährt, daß Gott selbst mit ihm redet; denn es ist der größte\footnote{Wittenberger und Erlanger: große.} Zorn, den er erzeigt, wenn er stille schweigt und redet nicht; und wiederum, die größte\footnote{Erlanger: große.} Gnade, wenn er sein Gebot und Willen hören läßt. Aber das ist eine überschwengliche Gnade, wenn er so freundlich und väterlich redet. Darum siehe auf die Worte, und halte sie gegen einander: „Ich will mit dir sein (spricht er), und will dich segnen.“ Wenn uns Gott also zuspräche, als er auch thut durch das Evangelium, wenn wir auch sollten gehen in ein Feuer, Hölle und Tod, was wollten wir lieber hören, denn daß er sagt: Gehe frisch hinan,
\quellseite{212}
ich will bei dir sein? Welch einen Trotz und Trost würden wir da haben, daß das Herz möchte schwärmen vor Freuden, und uns nicht fürchten vor tausend Töden?\footnote{Wittenberger und Jenaer: Todte.} Denn wenn der Gott, der alle Dinge in Händen hat, bei mir ist, und mich nicht verlassen will, was will mir denn schaden? Darum, wenn alle Teufel Ein Teufel wären, gehe ich dennoch frisch hinan.

10. Solchen Trost hat niemand gefühlt, denn Isaak, denn er ist da in einem fremden Lande, und in der theuren Zeit. Wie mag nun dem zu Muthe sein, der ein Fremdling ist, dem jedermann feind ist, und nicht die Brodrinde gönnt? Wo soll er hin? Alle Lande sind ihm zu, [er] hat Weib, Kind und Gesinde mit dem Viehe. Wo ist Haus, Korn, Heu, Stroh, Stall und allerlei Nahrung? [Er] ist noch dazu unter den Feinden; wo will er hinaus? Solches alles muß ihm dennoch wehe gethan haben. Da kommt nun Gott, und läßt ihn nicht, und spricht: Es hat keine Noth, die Theurung soll dich nicht umbringen. Zeuch nicht in Egypten, welches ein voll Land ist, sondern bleibe hier im Lande, ich will dich dennoch ernähren, daß du genug hast; als es auch geschah, daß er hundert Scheffel kriegte, wo ein anderer kaum Einen; er muß genug haben, sollte das ganze Land verhungern.

11. Darum hat er immerdar also müssen denken: Hier ist Weib, Kind und Gesinde, aber nichts dazu, damit ich sie könnte ernähren. Wohlan, Gott hat gesagt, er wolle bei mir sein, und mich segnen; das weiß ich, daß es wird wahr sein. [Er] hätte wohl mögen sagen: Wie lange will es währen? Wird doch nichts daraus! Aber er hält am Worte und glaubt nicht allein der Zusagung von zeitlicher Nahrung in der theuren Zeit, sondern auch von dem Samen, dadurch alle Welt sollte selig werden; darum ist sein Herz fröhlich gewesen, und [hat] gedacht: Ehe ich sollte Noth leiden, müßte eher ein Engel vom Himmel kommen und mich speisen.

12. Das ist nun der rechte Kern der Schrift, welches keine Vernunft sieht noch wahrnimmt, daß solche Gewalt in Gottes Worte verfasset ist, und solcher Glaube. Dieweil predigen sie in des Teufels Namen von Werken, damit man soll gen Himmel kommen, sehen nicht, wie in dieser Historie durch und durch eitel treffliche Exempel des Glaubens sind, die man sollte predigen, und also sagen: Siehe an Isaaks Glauben, nicht die Werke. Es ist ein schlechtes Werk, daß er aus einem Lande in das andere zeucht; aber er geht dahin auf (G)Gottes Wort, das ist er gewiß; du aber gehst hin und thust so viel Werke, und hast kein Gottes Wort.

13. Darum ist sein Werk mächtig, groß, und doch nur ein Hauswerk, kein Fasten, Beten, noch Kirchen und Messe stiften, und Erben halten, die unsere Narren für köstlich halten, deren doch Gott keines ansieht, ja, verwirft und verdammt sie zur Hölle darum, daß alles ohne Glauben ist. Dieses aber alles gilt so viel vor ihm, daß er es läßt schreiben und predigen in alle Welt. Nenne mir nun jetzt ein Werk, das diesem gleich sei. Es ist ein geringes Ding; solltest du es aber thun mit allen Pfaffen und Mönchen zusammen, würdest du es nicht wissen anzugeben, sondern darüber verzweifeln und zum Teufel fahren. Ach! daß Gott erbarme, wie haben wir tollen Narren gelehrt und gehandelt.

14. Das sei das erste Stück, wie Gottes Wort Isaak gepredigt wird, und er sich sein so heftig annimmt, und zum Exempel geschrieben; er hat es beweiset, gepredigt und getrieben, doch daneben auch viel Püffe erlitten. Denn Gott hätte solche Worte vergebens mit ihm nicht geredet, wo es nicht große Noth wäre gewesen, ihn zu stärken und trösten. Darum ist das Bischofsamt, zu predigen [und] regieren, auf ihm gelegen, [er] wird aber manchmal hier und da gezappelt haben; es ist auch noch Fleisch und Blut da gewesen, daß er wohl eines solchen Trostes bedurft hat, damit er auch weiter andere tröstet und stärkt. Es müssen auch fromme Leute gewesen sein, die ihm gefolgt haben und angehangen; [es] steht darauf, daß etliche auch wohl von ihm gelaufen sind. Also ist es ihm gegangen, und uns zum Vorbilde geschrieben, die wir die Verheißung haben, ein jeglicher für sich selbst. Wir sollen genug haben, wir leben oder sterben.

15. Ja, sagst du aber, wenn mir es Gott so zugesagt hätte, wie dem Isaak, so wäre gut glauben. Antwort: Wer weiß, wie er es ihm gesagt hat; es steht wohl darauf, daß er es gesagt habe durch einen ältesten Patriarchen, als Sem oder Eber. Aber wenn es gleich ein Engel vom Himmel geredet hat, ist es je so gewiß, wenn\footnote{In den Ausgaben: „als wenn“. Das „als“ scheint uns zu viel zu sein.} er es uns durch Menschen in der Schrift,
\quellseite{213}
oder durch Priester sagt. Die Zunge ist nicht Menschen, sondern Gottes Zunge; also auch die Predigt. Darum, wenn Gott selbst da stünde, und ein Kind in die Hand nähme, und taufete es, wäre es nichts\footnote{Erlanger: nicht.} gewisser, als wenn er es den Christen befiehlt, und durch uns Menschen thut und redet. Es ist ja sein, und nicht unser Wort, es fehlt nur daran, daß wir den Glauben nicht haben. Isaak wird auch viel Leute gehabt haben, die nichts von ihm gehalten, und ihn\footnote{„ihn“ fehlt in der Erlanger.} gar höhnlich verspottet haben; darum muß er sich immer an das Wort gehalten haben.\footnote{„haben“ fehlt in der Wittenberger und in der Erlanger.} So will er, daß auch wir allein auf das Wort sehen, er rede es, wie und wo er wolle. Das ist das Hauptstück dieses Capitels; wollen es nun nach dem Texte durchlaufen.

16. Das erste Stück, daß zu des frommen, heiligen Vaters Isaak Zeiten eine Theurung im Lande ist, gibt uns eine Frage: wie es zugeht, daß Gott gemeiniglich Theurung in das Land schickt, sonderlich wenn sein Wort geht? wie wir droben [Cap. 12, 10.] in Abrahams Historien auch gehört haben, und hernach unter Jakob und Joseph [Cap. 42, 1. ff.] und oft unter den Kindern Israel. Als, unter den Propheten Elia [1 Kön. 17, 1. 18, 2.] und Elisa [2 Kön. 6, 25.] war eine merkliche theure Zeit, und so fortan. Auch zu Christi Zeit, und der Apostel, ließ man von großer Theurung unter dem Kaiser Claudio [Apost. 11, 28.]. Jetzt hebt es sich auch hin und wieder in Landen, nachdem das Evangelium wieder aufgekommen ist, daß jedermann anhebt zu klagen, mehr denn zuvor jemals, daß es überall mangelt an Geld und Nahrung, gleichwie zu des Propheten Jeremiä Zeiten geschah, wie er [Cap. 44, 17. 18.] schreibt: Da traten zusammen Weib und Mann, und sagten: Ist es nicht eine Klage: dieweil wir dieneten und opferten der Königin des Himmels, hatten wir Brods genug, und ging uns wohl; weil wir aber aufgehöret haben, ihr zu dienen, müssen wir Noth leiden, und sind durch das Schwert und Hunger verzehret? Darum fuhren sie auch zu, und trieben es wie zuvor.

17. Also geht es nun: So bald man Gott einen Dienst soll thun, so fehlt es überall, und will jedermann verhungern. [Die] Ursache habe ich gesagt, und ist wohl möglich, wo das Evangelium im Schwange bleibt, daß uns alles zu enge und schmal werden wird. Eine Ursache ist, daß, wo das Evangelium angeht, da geht auch Sünde an, daß man es verachtet, verdammt, verfolgt, und lästert. Zuvor konnte man geben zweihundert Gulden, da man jetzt nicht Einen gibt armen Leuten; da bauete und stiftete man alle Städte voll Kirchen und Klöster; jetzt können wir alle nichts zuwege bringen.

18. Das ist eine Ursache, daß Gott die Welt wiederum plagt, und läßt die Nahrung zu kurz werden, weil man das Evangelium verachtet. (Es wird auch bei der Plage nicht bleiben, sondern [er] wird uns baß heimsuchen, wenn das Evangelium seinen Lauf gethan, und Raum genug gegeben ist zur Buße. Die Sünde will Gott nicht ungestraft lassen, daß man so schimpflich und verächtlich handelt mit seinem Worte. Je länger er Geduld trägt, je ärger wir es machen, und je unsinniger sie werden; so muß er denn drein schlagen.

19. Die andere Ursache gebe ich dem Teufel, doch durch Gottes Willen. Denn weil er ein Fürst der Welt ist, wie er sich rühmt gegen Christo, im Matthäo [Cap. 4, 8. 9.], so gebührt ihm auch, daß er dem Dinge widerstehe, so nicht für sein Reich ist. Daß man nun nicht so viel gibt, weil man das Evangelium predigt, als zuvor, ist nicht Wunder. Ein Narr wäre er, wenn er es gestattete. Weil wir nun wider ihn fechten, muß er uns also halten, daß wir Hungers sterben müssen, so viel an ihm liegt.

20. Darum ist es nicht ein böses Zeichen, ob er sich mehrt und sträubt, daß man jetzt nicht Einen guten Prediger ernähren kann, da man zuvor zweihundert Mönche füllte. Denn vorhin dienten sie ihm; darum mußte er seine Knechte auch versorgen; darum hat er alle Klöster und Stifte so wohl versehen, und allen genug gegeben; nun er aber sieht, daß man ein Loch in sein Regiment will reißen, wehrt er auf allen Seiten.

21. Ich habe einmal ein Exempel gehöret, wenn es nicht lügerlich lautete, reimte es sich wohl hierher: Wie einmal ein Hausvater war,
\quellseite{214}
ein ruchloser, roher Christ, der sich voll und fluchte ohn Unterlaß im Hause, lehrte es seine Kinder auch; da hatte er alles voll, Keller, Boden, Haus und Hof, bis einmal ein frommer Mann in sein Haus kam, und ihn bat, daß er solches abstelle, so lange, daß er sich ließ überreden, und aufhörte. Da fing es so bald an zu fehlen in allen Winkeln. Der Teufel that auch recht, daß er ihm nimmer wollte dienen und lohnen, weil er ihm nicht mehr diente.

22. Also liest man auch von St. Ambrosio, ob es wahr sei, weiß ich nicht, doch ist es der Wahrheit nicht unähnlich: daß er einst in eines reichen Mannes Haus kam, da sahe er, daß alles Dinges die Fülle war, und der Mann bekannte, daß [es] ihm alle seine Lebtage nach allem seinem Willen gegangen wäre. Da erschrak der fromme Mann, und sprach zu denen, die bei ihm waren: Hier ist hohe Zeit, daß wir davon gehen, denn da regiert der Teufel gar. Und als er heraus kam, fiel das Haus ein, und die Erde verschlang alles, was da war.

23. Also soll es zugehen: Der Teufel ist der Welt Herr, darum wird er Gottes Dienern das Futter nicht geben. Sollst du es aber haben, so muß [es] Gott geben, und dich ernähren. Darum schickt es Gott also, daß die, so Christen sind, müssen in der theuren Zeit leben. Wenn sich nun solch Schreien und Klagen hebt, weil man das Evangelium predigt, da mußt du klug sein, und wissen recht zu antworten. Zuvor hat der Teufel alles genug gegeben, da mir seine Knechte waren; nun aber das Evangelium aufgekommen ist, zieht er die Hand zurück. So straft Gott auch die Welt, wie [es] recht ist; aber stehe fest, und glaube an Gott, so wirst du mitten in der Theurung genug haben. Und mich dünkt, daß der Spruch im 37. Psalm, V. 25., hieher gehe und sehe: „Ich bin jung gewesen und alt worden; aber ich habe noch nie gesehen einen Gerechten Hunger leiden, oder seine Kinder nach Brod gehen.“

24. Gott läßt es geschehen, daß die Frommen in der theuren\footnote{Erlanger: theurer.} Zeit leben müssen, auf daß er beweise, wie sein Wort wahrhaftig ist, und er seine Kinder ernähre aus dem Glauben. Und [es] ist ein Zeichen, daß das Evangelium recht sei, Furcht bringe, und nicht vergebens gepredigt werde. Wo es nicht wäre, so dächte ich, der Teufel fühlte es nicht. Weil er es aber fühlt, so sieht man, daß es Kraft habe in den Leuten, und sie zurück zieht. Aber lasse ihn ziehen, es wird dennoch Saat bleiben. Wer nicht glaubt, der sterbe Hungers; wer aber glaubt, der wird genug haben, denn Gott lügt nicht, sollte es auch Korn vom Himmel regnen. Es thut’s auch Gott nur darum, daß er die Ungläubigen strafe, und die da glauben versuche und stärke. Das ist ein Stück dieses Capitels. Folgt das andere:

V. 6—8. Also wohnte Isaac zu Gerar. Und wenn die Leute am selben Orte fragten von seinem Weibe, so sprach er: Sie ist meine Schwester. Denn er fürchtete sich zu sagen, sie ist mein Weib, sie möchten mich erwürgen um Rebekka willen; denn sie war schön von Angesicht. Als er nun eine Zeitlang da war, sahe Abimelech, der Philister König, durch das Fenster, und ward gewahr, daß Isaac scherzte mit seinem Weibe Rebekka.

25. Es ist leicht zu merken, warum der gute Vater sich also gefürchtet habe. Zu der Zeit war das Evangelium noch nicht ausgegangen, daß zu hoffen wäre gewesen, daß etliche fromme Leute im Lande sollten sein; es ist allein ein Regiment nach der Natur und Vernunft hin gewesen. Wo aber Gottes Wort und Gnade nicht lebt bei den Leuten, darf man sich nichts Gutes versehen; da kann Vernunft und Natur nicht fromm bleiben; sind alle Lügner und Schälke; oder, müssen sie fromm sein, so thun sie es aus Zwang. Derselben kannst du denken, daß dies Land und Königreich ein weltlich Wesen mußte sein, daß jedermann gethan hat, was er wollte, und sonderlich der König hat ihm mit gutem Fuge mögen das Weib nehmen.

26. Und ist abereins,\footnote{. abereins = abermals.} daß Moses Gottes Thorheit vorschlägt, daß er solch närrisch Werk her schreibt, wie Isaac sagte, sie sei seine Schwester, und der König habe zum Fenster ausgesehen, daß er mit ihr scherzte; nicht, daß er etwas Unzüchtiges gesehen habe, sondern daß er aus äußerlichen Geberden gemerkt habe, daß es nicht seine Schwester wäre, sondern müßte sein Weib sein.

27. Mußte er denn so eben das schreiben? Wie ich zuvor gesagt habe, sage ich noch. Wenn du in der Schrift siehst solche thörliche Gottes
\quellseite{215}
Merke, daß du wissest, es gelte, die Welt zu schänden. Darum muß der geringe Heilige mit geringen Werken betreten, und vor der Welt verspottet sein; sie sieht aber nicht, wie solches Werk gegangen ist im Glauben, in Gottes Wort und Wohlgefallen. So fährt sie zu, und mißt die Werke nach ihrem eigenen Gutdünken, so doch alle Schrift sagt, daß alle Werke ihr Gewicht haben von Gottes Wohlgefallen, und nichts liegt an der Menge und Größe der Werke. Also muß dies Werk auch köstlich sein und gepredigt werden, wie sehr es die Vernunft und heuchlerische Heilige für Narrheit halten. Dagegen müssen sie auch sehen, daß Gott nach ihrer Heiligkeit nichts fragt, läßt kein Wort davon schreiben, und ihnen zuwider mit solchem Narrenwerke umgeht.

28. Aber doch ist es Isaak nicht große Freude gewesen, daß er mußte in dem fremden Lande sein, und ist alle Tage mit dem Weibe in Fahr gestanden seines Lebens, und ihrer Ehre. Darum hat Gott nicht gewollt, daß er lange in Fährlichkeit wäre; wie hier selbst der König spricht: „Du hättest eine große Sünde auf uns mögen bringen.“ Darum hat es Gott so geschickt, daß es offenbar würde, auf daß ihm, auch Rebekka, nichts Leides widerführe. Da blickt hervor die Sorge und der Fleiß, so Gott trägt für die, so ihm trauen. Er hat das Weib müssen wagen, und in die Schanze setzen, und Gott heimgeben, daß er sie verwahrete, wo sie ihm würde genommen. [Dies] ist ein groß Exempel des Glaubens, und dabei ein großer Trost, daß sie Gott bewahrt und rein behält. Solches hat nun Moses nicht umsonst so oft geschrieben, wie wir es auch droben [Cap. 12, 13. 20, 2.] von Abraham zweimal gehört haben. Isaak ist auch ohne Zweifel wohl versucht worden; denn es ist gar schwer, so bloß alles Gott heimzustellen, daß er hat mögen denken: Du heißest mich in das Land ziehen, und steckst mich in allerlei Fährlichkeit meines Leibes, und meines Weibes; es ist ihm aber nütze, seinen Glauben zu stärken. Denn so fährt Gott immerdar, auf daß er beweise, wie stark der Glaube sei, und immer stehe wider Schande, Tod und alles Unglück.

29. Hier fragt man nun: ob Isaak gelogen habe, daß er spricht, sie ist meine Schwester? Das lasse ich fahren. Ist es Sünde, so sei es Sünde; wir haben nicht im Sinne, die lieben Heiligen zu entschuldigen, daß sie nie gesündigt haben. Gottes Gnade wollen wir in ihnen hoch preisen, aber ihr Wesen wollen wir nicht so hoch heben. Doch wäre es leichtlich zu entschuldigen, es sei keine Lüge, denn sie war seine Schwester nach dem Geiste.

30. Es ist aber nicht Wunder, ob ein Christen-mensch strauchelt. Wenn wir in solcher Fahr ständen, würden wir zumalen auch nicht stracks zu sagen. Es bleibt dennoch Gottes Gnade, die ihn erhält und schützt, ob er schon fällt. Das ist nun die Gnade, daß Gott durch Abimelech, den König, ausrichtet, daß er ihm Ruhe und Frieden gibt auch unter den Leuten, daß er sich setzt, und im Lande nährt. Das sei die erste Aufrichtung, in diesem Capitel beschrieben. Folgt weiter:

V. 12—16. Und Isaak säete in dem Lande, und fand desselben Jahres hundertfältig; denn der Herr segnete ihn. Und er ward ein großer Mann, ging und nahm zu, bis er fast groß ward, daß er viel Gutes hatte an kleinem und großem Vieh, und ein großes Gesinde. Darum neideten ihn die Philister, und verstopften alle Brunnen, die seines Vaters Knechte gegraben hatten, zur Zeit Abrahams, seines Vaters, und fülleten sie mit Erde, daß auch Abimelech zu ihm sprach: Zeuch von uns, denn du bist uns zu mächtig worden.

31. Gott hat den Patriarchen hin und her geworfen, wie einen Ball, auch wohl zwischen die Sporn gefaßt, und würde gemacht in seinem ganzen Leben, daß man je sehe, wie seine Exempel des Glaubens hier vorgebildet sind. So ist nun hier angezeigt das Glück, wie es den Frommen geht. Und siehe, wie fein beschreibt es Moses, wie die Schrift schreiben soll: gibt Isaak nichts, das er erworben habe, daß ja alles Gut, auch zeitlich, müsse von Gottes Segen kommen; denn ihrer immerdar sind viel gewesen, die auch gearbeitet, und doch nichts erworben haben. Nun hat ihm Gott eine Zeit Ruhe gegeben, aber es geht bald wieder Unglück an.

32. Also geht unser Leben. Eitel Frieden können wir nicht haben, darum muß es genenget sein, daß jetzt die Sonne scheint, jetzt Wind und Regen kommt. So muß Isaak nun wieder Verfolgung und Neid leiden, [sie] werden ihm so feind, daß sie ihm des Wassers nicht gönnen, und alle seine Brunnen verstopfen. Das Land
\quellseite{216}
hat eine sonderliche Plage Wassers halben, daß es nicht überall wohl zu finden ist. Darum haben sie die Brunnen köstlich und werth gehalten, daher sich viel Hader erhoben hat; denn es ist ein bergig, heißes und trocken Land, hat nicht viel Bäche und Wasserströme.

B. 17. 18. Da zog Isaak von dannen, und schlug sein Gezelt auf im Thale Gerar, und wohnete allda, und ließ die Wasserbrunnen wieder aufgraben, die sie zu Abrahams Zeiten, seines Vaters, gegraben hatten, welche die Philister verstopfet hatten nach Abrahams Tode.

33. Da beschreibt Moses vier Wallfahrten nach einander des heiligen Vaters. Die erste zu dem König, wie gehört ist. Die andere gen Gerar. Darnach ist er aber vertrieben, bis ins vierte Mal, daß er immer ein armer Pilgrim sein mußte. [Es] wäre nicht wohl möglich gewesen, ohne hohen Glauben zu bestehen, mit Weib, Kind, Gesinde immer auf[zu]brechen, und irre zu ziehen, und keine gewisse Stätte, Haus, Stall, Acker oder Wiesen zu haben, sondern alles mit großer Kost und Mühe mit sich [zu] schleppen, und im fremden Lande alles kaufen und mieten. Der Glaube hat ihn müssen stärken, trösten und erhalten, sonst wäre er bald unleidlich worden; auf daß man sehe, wie die lieben Väter mehr gelitten haben, denn wir bedenken.

34. Deß hat er sich getröstet und gehalten, daß GOtt ihm zusagte: „Ich will bei dir sein.“ Es ist niemand bei ihm denn GOtt; [er] stellt sich aber also, als wäre er auch nicht bei ihm. Also haben wir das größte Stück der Legende dieses Patriarchen in diesem Capitel, eitel theuer und edel Ding, gefällt GOtt so wohl, deß\footnote{Wittenberger und Erlanger: daß.} war er sicher, denn er hatte GOttes Wort. Vor der Welt war es ein jämmerliches, elendes Leben; aber GOtt hat es angesehen in Gnaden. Es muß also gehen, daß die Welt das Kreuz davor macht, und sich scheuet; sie will wissen, wo sie sitzen und bleiben soll; weiß sie es nicht, so hebt sie nichts an. Dieser aber hangt allein an dem Worte. Wo er hinkommt, ist es alles wider ihn, deß muß er sich vorhin erwägen. Siehe, das ist die einfältige Legende, darinne kein köstlich gleichend Werk steht, sondern eitel Kreuz, Verfolgung und Armuth; aber in großem Glauben.

35. So ist nun das Erste hier, daß er gen Gerar gezogen ist, und sich da setzen will, und macht da zween Brunnen, die sein eigen sollten sein, wie sie seines Vaters gewesen waren; und bauete noch einen Brunnen dazu, den hieß er [עֵשֶׂק] Esek, das ist, Hohnbrunn, den haben sie lassen fahren; wiewohl es schwer ist. Denn die Hirten von Gerar sprachen: Er ist unser, wir sind Herren, so bist du ein Gast; nehmen mit Gewalt und rauben, hätten wohl mit ihm genug daran gehabt, treiben ihn aber hinneben mit allem, das er hat. Das muß er leiden und weichen, schilt und rächt nicht, thut nicht mehr, denn daß er den Brunnen Esek nennt.

36. Da er hinweg kommt, bauet er aber einen; den muß er aber lassen gehen, und sich davon machen. Darum heißt er ihn [שִׂטְנָה] Sitna, das ist, Widerstand, davon das Wort Satan kommt [Hiob 1, 6. 7.], damit Christus und die Schrift [Matth. 4, 10.] den Teufel nennen, das ist, ein Widersacher. Als auch Paulus den Papst nennt zum Thessalonichern [2. Ep. 2, 4.] ἀντικείμενος, adversarius, der nur Widerspiel thut, feiert nicht, und thut anders nichts, denn was Christo wider ist. Den Namen gibt Christus dem Teufel nicht umsonst; er weiß wohl, wie er nicht ruht, greift uns auf allen Seiten an. So heißt er Petrum auch im Matthäo [Cap. 16, 23.]: „Hebe dich hinter mich, Satan.“ Also nennt hier Isaak den Brunnen, darum, daß sie sich immer wider ihn setzten, und nicht erlauben noch gönnen wollten, daß er seines Vaters Brunnen ausgrabe.

37. Da weicht er abermal, bis er einmal Frieden überkommt, und gräbt einen Brunnen, den heißt er [רְחֹבוֹת] Rehoboth, das ist, Raum und Breite; als sollte er sagen: GOtt sei gelobt, daß wir einmal Raum haben. Welches Wort genugsam anzeigt, wie er zuvor gedrängt ist worden, daß man ihn hat wollen weder sehen noch hören. Aber da bleibt er auch nicht lange an dem Orte, und reist wieder gen Bersaba, da er hatte mit seinem Vater lange gewohnt. Was nun Moses schreibt, wie er mit dem Könige sich vereiniget, und einen Bund macht, ist aber ein Trost, daß GOtt die Gläubigen nicht läßt, wenn sie nur fest halten; schafft ihnen Unglück genug, läßt sie aber nicht ohne Trost. Das sei von der Historie gesagt. Nun sollten wir auch die Mysteria und heimliche Deutung handeln.

38. Wir haben gehört, daß Isaak Christi
\quellseite{217}
Figur ist, also wird er noch eine Weile in dem Capitel bleiben. Nun steht Christi Reich also, daß es ist ein Reich des Lebens, und so starkes Lebens, daß es mitten im Tode lebt; und so starker Gnade, daß sie in der Sünde überhand behält, und dem Teufel mitten im Rachen regiert. Denn es ist ein geistlich Reich, darum scheint es nicht, sondern eben scheint das Widerspiel. Wenn man die Christen ansieht, als, die heiligen Märtyrer, sieht man nichts denn eitel Tod; wie der 44. Psalm, V. 23., sagt: „Wir werden um deinetwillen täglich erwürget, und sind geachtet, wie die Schlachtschafe.“ Wo die Christenheit ist, da muß es Blut kosten, oder [es] sind nicht rechte Christen. Es sind nicht Weideschafe, sondern Schlachtschafe, immer eines nach dem andern hin. So ist das Ansehen des christlichen Lebens nichts denn Schwachheit, Tod und Sünde; noch regiert er es im Leben, geistlich, welches niemand sieht, sondern allein der Glaube faßt.

39. Deß Bild trägt nun Isaac, wie er zuvor Christi Figur war mit seinem Opfer, verurtheilt und hingegeben zum Tode, und doch lebend geblieben. Also ist er auch hier Christi Figur in seinem Reiche, wie er hier fährt in dem fremden, unbekannten Lande. Wenn man sein Leben ansieht, steht es alle Stunden in Fahr des Todes, dazu seines Weibes: noch wird er erhalten, nur zum Zeichen, daß ein christlich Leben auch soll also gehen, und in aller Schande und Fahr stehen. Die Christen müssen den Titel führen, und in der Fahr stecken, daß [es] Ketzerei sei, und eine Hure heiße; die andern wollen rechten Glauben haben, und GOttes eigen sein.

40. Was ist es aber, daß zuvor eine Theurung in das Land kommt? Das ist es: Wenn das Evangelium recht angeht, muß sich zuvor ein Hunger und Kummer heben im Gewissen. Das Evangelium kommt [zu] niemand, der da Fülle und gute Tage hat, sondern allein [zu] den geängsteten Gewissen, die in großem Hunger sind, und eine ledige Seele haben, die gerne solche tröstliche Predigt hört.

41. Daß nun hernach steht vom Brunnengraben, und daß [es] Isaac eine Weile glückselig geht, ist alles der Lauf der Christenheit; eine Weile hat sie Ruhe, daß sie zunimmt, und bessert sich; als, zu der Apostel Zeit ging sie bald auf, und wuchs; aber bald erhebt sich Hader, Zank und Ketzerei.

42. Der Brunnen aber, den sie zustopften, ist nichts denn der Frevel, den sie thun durch Menschenlehre, über die heilige Schrift, welche die Christenheit eine Zeitlang führte und rein hatte; aber bald kam der Teufel, und führte Menschenstand hinein, daß man ihrer nicht mehr genießen konnte, legten sie aus nach ihrem Kopfe. Das heißt eben, Erde in [den] Brunnen getragen, und damit gestopft, daß er je anzeige die, so die Schrift verderben mit irdischem Verstande, und sie lenken nach ihrem fleischlichen Wahne; wie die Ketzer Arius und Pelagius und andere, und hernach der Pabst, der sich rühmt, wie er allein der Schrift Meister, und sie Macht auszulegen habe; welcher nichts Anderes getrieben hat im ganzen geistlichen Rechte, und durch alle hohen Schulen, denn eitel Erde und Mist eingeführt, und die Schrift so verderbt und verschlemmt, daß [es] nicht möglich ist, aus ihren Büchern etwas zu verstehen aus der Schrift, sondern nur verhindert, und zurück daraus studiren.

43. Also wird es uns noch weiter gehen, wie es schon geht, durch manche Rottengeister, und bleiben bis an [den] jüngsten Tag. Wir haben den Brunnen auch gegraben und aufgethan; so kommen immer andere, die ihn wieder mit Erde zustopfen. Deß mögen wir uns kühnlich versehen: Wenn wir uns aus der Schrift tränken und stärken wollen, so will man es uns wehren und stärken. Darum muß sie auch den Namen haben: Hohn und Widerstand, das ist, daß wir darüber gehöhnt werden, und Widerstand leiden müssen; wo nicht, so haben wir die Schrift nicht recht, bis so lange, daß uns GOtt Raum gibt, daß wir bei rechtem Verstande der Schrift bleiben, und Ketzerei überwinden.

44. So haben wir die Historie mit der heimlichen Deutung, darin wir sehen, wie es alles zu thun ist um das Evangelium, und Reich Christi allein zu predigen. [Es] muß aber Widerstand leiden, und verfolgt werden. Nun folgt eine andere Legende des Patriarchen Jakob.
\quellseite{218}

\section{Das siebenundzwanzigste Capitel}

Cap. 26, 34. 35. Da Esau vierzig Jahr alt war, nahm er Weiber, Judith, die Tochter Beri, des Hethiters, und Basemath, die Tochter Elon, des Hethiters, die waren beide ungehorsam Isaak und Rebekka.

Cap. 27, 1—35. Und es begab sich, da Isaak war alt worden, daß seine Augen dunkel worden zu sehen, rief er Esau, seinem größern Sohne, und sprach zu ihm: Mein Sohn. Er aber antwortete ihm: Hier bin ich. Und er sprach: Siehe, ich bin alt worden, und weiß nicht, wann ich sterben soll. So nimm nun deinen Zeug, Köcher und Bogen, und gehe auf das Feld, und jage mir ein Wildpret, und mache mir ein Essen, wie ich es gerne habe, und bringe mir es herein, daß ich esse, daß dich meine Seele segne, ehe ich sterbe. Rebekka aber hörte solche Worte, die Isaak zu seinem Sohne sagte. Und Esau ging hin auf das Feld, daß er ein Wildpret jagte, und heimbrächte. Da sprach Rebekka zu Jakob, ihrem Sohne: Siehe, ich habe gehört deinen Vater reden mit Esau, deinem Bruder, und sagen: Bringe mir ein Wildpret, und mache mir ein Essen, daß ich esse, und dich segne vor dem HErrn, ehe ich sterbe. So höre nun, mein Sohn, meine Stimme, was ich dir heiße. Gehe hin zu der Heerde, und hole mir zween gute Böcklein, daß ich deinem Vater ein Essen davon mache, wie er es gerne hat. Das sollst du deinem Vater hinein tragen, daß er esse, auf daß er dich segne vor seinem Tode. Jakob aber sprach zu seiner Mutter Rebekka: Siehe, mein Bruder Esau ist rauh, und ich glatt; so möchte vielleicht mein Vater mich betasten, und würde vor ihm geachtet, als [ob] ich ihn verleitet,\footnote{Dazu hat die Jenaer die Randglosse: „betrigen wollt“, was auch unsre Bibel bietet.} und brächte über mich einen Fluch, und nicht einen Segen. Da sprach seine Mutter zu ihm: Der Fluch sei auf mir, mein Sohn, gehorche nur meiner Stimme, gehe und hole mir. Da ging er hin und holte es, und brachte es seiner Mutter. Da machte seine Mutter ein Essen, wie sein Vater gerne hatte. Und nahm Esaus, ihres größern Sohnes, köstliche Kleider, die sie bei sich im Hause hatte, und zog sie Jakob an, ihrem kleinern Sohne. Aber die Felle von den Böcklein that sie ihm um seine Hände, und wo er glatt war am Halse, und gab also das Essen mit Brod, wie sie es gemacht hatte, in Jakobs Hand, ihres Sohnes. Und er brachte es hinein zu seinem Vater, und sprach: Mein Vater. Er antwortete: Hier bin ich. Wer bist du, mein Sohn? Jakob sprach: Ich bin Esau, dein erstgeborner Sohn, ich habe gethan, wie du mir gesagt hast. Stehe auf, setze dich, und iß von meinem Wildpret, auf daß mich deine Seele segne. Isaak aber sprach zu seinem Sohne: Mein Sohn, wie hast du so bald funden? Er antwortete: Der HErr, dein GOtt, bescherte mir es. Da sprach Isaak zu Jakob: Tritt herzu, mein Sohn, daß ich dich betaste, ob du seiest mein Sohn Esau, oder nicht. Also trat Jakob zu seinem Vater Isaak, und da er ihn betastet hatte, sprach er: Die Stimme ist Jakobs Stimme, aber die Hände sind Esaus Hände. Und er kannte ihn nicht, denn seine Hände waren rauh, wie Esaus, seines Bruders Hände, und segnete ihn. Und sprach zu ihm: Bist du mein Sohn Esau? Er antwortete: Ja, ich bin es. Da sprach er: So bringe mir her, mein Sohn, daß ich esse von deinem Wildpret, daß dich meine Seele segne. Da brachte er es ihm, und er aß, und trug ihm auch Wein hinein, und er trank. Und Isaak, sein Vater, sprach zu ihm: Komm her und küsse mich, mein Sohn. Er trat hinzu und küssete ihn. Da roch er den Geruch seiner Kleider, und segnete ihn, und sprach: Siehe, der Geruch meines Sohnes ist wie ein Geruch des Feldes, das der HErr gesegnet hat. GOtt gebe dir vom Thau des Himmels, und von der Fettigkeit der Erde, und Korn und Weins die Fülle. Völker müssen dir dienen, und Leute müssen dir zu Füße fallen. Sei ein Herr über deine Brüder, und deiner Mutter Kinder müssen dir zu Füße fallen. Verflucht sei, wer dich verfluchet; gesegnet sei, wer dich segnet. Als nun Isaak vollendet hatte den Segen über Jakob, und Jakob kaum hinaus gegangen war von seinem Vater Isaak, da kam Esau, sein Bruder, von seiner Jagd, und machte auch ein Essen, und trug es hinein zu seinem Vater, und sprach zu ihm: Stehe auf, mein
\quellseite{219}
Vater, und ich von dem Wildpret deines Sohnes, daß mich deine Seele segne. Da antwortete ihm Isaak, sein Vater: Wer bist du? Er sprach: Ich bin Esau, dein erstgeborner Sohn. Da entsetzte sich Isaak über die Maße sehr, und sprach: Wer? Wo ist denn der Jäger, der mir gebracht hat, und ich habe von allem gegessen, ehe du kamst, und habe ihn gesegnet? Er wird auch gesegnet bleiben. Als Esau diese Rede seines Vaters hörte, schrie er laut, und ward über die Maße betrübt, und sprach zu seinem Vater: Segne mich auch, mein Vater. Er aber sprach: Dein Bruder ist gekommen mit List und hat deinen Segen hinweg. Da sprach er: Er heißt wohl Jakob, denn er hat mich nun zweimal untergetreten.\footnote{\unsicher{Handbalge}: Esau, heißt eine Fußsohle, daher kommt Jakob oder Jacob, ein Untertreter, oder der mit Füßen tritt, und bedeutet alle Gläubigen, die durch das Evangelium die Welt und das Fleisch, und den Teufel mit Sünde und Tod unter sich treten.} Meine erste Geburt hat er dahin, und siehe, nun nimmt er auch meinen Segen; und sprach: Hast du mir nicht einen Segen vorbehalten? Isaak antwortete, und sprach zu ihm: Ich habe ihn zum Herrn über dich gesetzt, und alle seine Brüder habe ich ihm zu Knechten gemacht, mit Korn und Wein habe ich ihn versehen; was soll ich doch dir nun thun, mein Sohn? Esau sprach zu seinem Vater: Hast du denn nur Einen Segen, mein Vater? segne mich auch, mein Vater. Und hub auf seine Stimme und weinete. Da antwortete Isaak, sein Vater, und sprach zu ihm: Siehe da, du wirst eine fette Wohnung haben auf Erden, und vom Thaue des Himmels von oben her. Deines Schwertes wirst du dich nähren, und deinem Bruder dienen. Und es wird geschehen, daß du sein Joch ablegest, und von deinem Halse reißest. Und Esau war Jakob gram, um des Segens willen, damit ihn sein Vater gesegnet hatte, und sprach in seinem Herzen: Es wird die Zeit bald kommen, daß mein Vater Leide tragen muß, denn ich will meinen Bruder Jakob erwürgen. Da ward Rebekka angesagt dies Wort ihres größern Sohnes, und schickte hin, und ließ Jakob, ihren kleinern Sohn, rufen, und sprach zu ihm: Siehe, dein Bruder Esau dräuet dir, daß er dich erwürgen will. Und nun höre meine Stimme, mein Sohn, mache dich auf, und fleuch zu meinem Bruder Laban in Haran, und bleibe eine Weile bei ihm, bis sich der Grimm deines Bruders wende, und bis sich sein Zorn wider dich von dir wende, und vergesse, was du an ihm gethan hast; so will ich darnach schicken, und dich von dannen holen lassen. Warum sollte ich euer beider beraubt werden auf Einen Tag? Und Rebekka sprach zu Isaak: Mich verdrießt zu leben vor den Töchtern Heth; wo Jakob ein Weib nimmt von den Töchtern Heth, die da sind wie die Töchter dieses Landes, was soll mir das Leben?

1. In diesem Capitel ist das Vornehmste der Segen, den Isaak seinem Sohne Jakob gegeben hat. Damit geht an der dritte Patriarch, von welchem Christus kommen sollte, davon wir viel hören werden, wie ihn Gott geführt hat. Proben [Cap. 25, 31. ff.] haben wir gehört, wie die zween Brüder mit einander kaufschlagten\footnote{kaufschlagen = einen Handel abschließen.} um die erste Geburt; und Esau nahm das rothe Gemüße, und gab dafür sein Recht, das er hatte, damit er auch vor den Leuten die rechte Erstgeburt verlor, wie es ihm Gott zuvor genommen hatte, als er sprach zu Rebekka: Der Größere soll dem Kleinern dienen. Darum beschreibt nun Moses, wie es gekommen ist, daß der Segen von Esau fällt auf Jakob.

2. So ist nun in dem Capitel nichts beschrieben nach der Historie, denn wie Gott wahrhaftig ist, und wunderbarlich, daß er hält, was er redet [Ps. 33, 4.], und daselbige Halten so seltsam und wunderlich\footnote{\unsicher{Tenaar}: wunderbarlich.} ausrichtet, daß es weit übertrifft alle Vernunft. Geredet hat er, daß der Kleinste sollte Herr werden, und der Größeste dienen. Das müßte wahr werden, hat sich aber verzogen eine gute Weile. Denn der Esau war nun schon ein ehelicher Mann von vierzig Jahren, und zeugte bald Kinder, bis Jakob fast bei siebenzig Jahren war, wie wir hören werden, ehe er ein Weib nahm, und mußte lassen den Bruder Recht haben und herrschen, und ist ein Aschenbrödel; doch hat er immer an dem Worte gehalten. Dieweil fährt Esau fort, ist der liebe Sohn, hat das Recht gar inne, und wiewohl er es droben [Cap. 25, 33.] verkauft hatte, dachte er also: Verkaufst hin, verkaufst her, dennoch bleibe ich Herr im Hause, dazu behalte ich den Segen. Darum nimmt er zwei Weiber, und setzt sich ein, und hält Haus.

3. Also zeigt uns die Schrift an Gottes Wunderwerk, daß die Verheißung so lange aufgezogen
\quellseite{220}
ist, ehe denn er den Segen überkommt, der ihm gebührt, und läßt neben ihm den Esau aufwachsen, und in das Regiment greifen, der doch den Segen verloren hatte. Und zeigt merklich an den großen Trotz Esaus, daß er der Sache so gewiß ist, und nimmt zwei Weiber, wider Vater und Mutter Willen, die ihnen beide ungehorsam sind gewesen, zwo stolze Metzen, wie sie pflegen zu thun, wenn sie die Schlüssel kriegen. [Sie] haben also gedacht: Wir sind Töchter des Landes, Esau soll froh werden, daß wir in sein Haus und Geschlecht kommen; denn er ist fremd und ein Gast im Lande.

4. Das hat nun Isaac und Rebekka verdrossen. Aber was sollen sie thun? Sie leiden es. Isaac läßt es gehen, und bleibt gleichwohl auf dem Wahn, Esau soll der rechte Sohn sein, dem der Segen gebührt. Was meinst du, daß auch Jakob die Zeit hat müssen leiden, weil es die Eltern mußten leiden, und gehalten haben, daß es einen Schein hatte, als sollte es alles widersinnisch gehen, was Gott geredet hat. Solch Bild muß Jakob tragen, daß man lerne Gott erkennen in seinen Werken. Also mußte er hören und fühlen: Du solltest der Größeste werden, mußt aber der Kleinste sein; Vater und Mutter müssen nicht recht haben, und der Weiber Trotz und Muthwillen leiden; sollte denn Jakob Recht haben? Wie oft hat er gedacht: Soll das der Segen sein? Meine Eltern müssen die Weiber regieren lassen, was sollte ich denn dawider schaffen? Ist das gehalten, was Gott redet?

5. Aber so geht es in allen Sachen, die Gott treibt. Er that es seinem eigenen Sohne auch. Als er ihn wollte zum höchsten Könige machen Himmels und Erden, fährt er zu, und läßt ihn an das Kreuz hängen, und also zurichten, als nie keinen Menschen. Also hat er den heiligen Märtyrern auch mitgespielt. Darum läßt er den Titel von sich ausrufen, daß er heiße Wunderbarlich; wie der 4. Psalm, 2. 4., sagt: „Erkennet doch, daß Gott seine Heiligen wunderlich führet“; darum sagt er, V. 5.: „Zürnet ihr, so sündiget nicht“; item, V. 6.: „Opfert rechte Opfer, und hoffet auf den Herrn“, das ist, haltet stille, und lasset es gehen; wenn er sich noch so wunderlich stellt, muß es dennoch geschehen; es ist seine Weise nicht anders. Also wird es uns gehen, wenn wir sterben sollen; da soll er sich wohl stellen, als sei er uns feind, und sei nichts denn Teufel um uns.

und sei nichts denn Teufel um uns. Wer da nicht seine Weise weiß und kennt, der steht übel. Vernunft vermag es nicht, der Glaube muß aber also denken: Gott hat es geredet, das ist wahr, er stelle sich nun, wie er wolle.

6. Das ist nun der Anfang, daß Jakob mußte glauben, er sollte der Größte werden, und dennoch Knecht im Hause bleiben,\footnote{Erlanger: bleibe.} und sehen, daß sie ihren Hochmuth trieben zu Trotz Vater und Mutter. Da hat Esau gedacht: Nun habe ich es in der Hand; ist sicher und gewiß. So läßt ihn Gott aufsteifen und prangen; aber ehe er sich umsieht, so stürzt er ihn; als er am allersichersten ist, liegt er danieder, und Jakob bringt den Segen hinweg. Das heißt gespielt nach Gottes Weisheit. Es sind einfältige, närrische Historien, aber man könnte Gott nicht also abmalen mit seiner Farbe, als er da abgemalt ist. Isaac lässet er betrogen werden und hebt es durch das Weib an; aber so närrisch, daß Jakob selbst daran verzagt.

7. Wer gibt ihr aber den Trotz, daß sie also herfährt, und spricht: „Der Fluch gehe über mich, mein Sohn.“ Da muß das Ding ausgerichtet werden, daß Jakob Herr wird durch des Weibes Glauben, nicht durch den heiligen Mann Isaac, und gewinnt so einen großen Geist, daß sie es wagen darf, ob Isaac gleich flucht. Der Teufel sollte es wagen, daß mir ein solcher heiliger Mann fluchte. Ich wollte lieber, daß mich alle Päpste und die Welt verdammte und verfluchte. Sie wußte, daß er der Mann war, dem Gott so viel zugesagt hatte, und auf ihn ein Auge hätte, und mit ihm redete; noch darf sie so trotzen. Das wäre nicht möglich einem geringen Geiste zu thun; es muß ein sonderlicher, großer Glaube in dem Weibe gewesen sein, der sich tröstlich auf das Wort verlassen hat: „Der Kleine soll Herr sein, der Größere soll dienen.“ Daher schließt sie: Soll das wahr sein, so wird Gott meinem Sohne nicht lassen fluchen; und ob er gleich fluchte, so müßte der Fluch auch nicht recht sein. Welches gerade so viel ist, als wenn ich spräche: Ich glaube an Christum, und wenn er selbst käme, und fluchte mir, wollte ich mich nicht daran kehren. Denn sie hat so wenig gezweifelt, daß der Mann fromm wäre, und Gottes Wort hätte, als ich von Christo. Darum hat sie gedacht
\quellseite{221}
: Ist es ein Fluch, so muß [es] entweder ein Schein sein, oder GOtt wird es wenden. Wie ich müßte thun am Todtbette, wenn er spräche, er wollte meiner nicht, daß er es entweder nicht mit Ernst meinte, oder GOtt würde es nicht leiden.

8. Darum ist GOtt ein wunderbarer Herr in seinen Werken. Da muß alle Vernunft zu Boden sinken, und der Glaube nach nichts fragen, denn nach GOttes Worte. Er wende und mache es, wie er wolle, so muß es wahr bleiben. Das sind die trefflichen Exempel, wie GOtt fährt mit seinen lieben Heiligen. Das sei nun von den zwei Stücken gesagt, wie GOtt wahrhaftig ist; aber also, daß er es wunderlich hinaus führt, auf daß wir getrost sind, ob er verzeucht, und [es] nicht hinaus geht, wie wir denken. Also muß Jakob zum Segen kommen, sollte auch die ganze Welt dawider sein, kriegt ihn aber so wunderlich, daß [es] niemand konnte ermessen. So wunderlich wird auch Rebekka geführt; sie kann dem Vater den Segen nicht nehmen, thut ihm kein Leid, läßt ihm nach, wie er es machen will, noch gibt GOtt einen Geist, daß sie es findet, und den Segen hinweg stiehlt, ohne sein Wissen und Willen, also, daß er gleichwohl auf Jakob bleibt, auf daß man eigentlich sehe, daß es GOttes Werk ist.

9. Nun sollen wir dies Capitel auch geistlich handeln. Es sind viel guter Sprüche und Stücke darinne. [Wir] wollen das vor uns nehmen, das auch in die Historie gehört, wie zwei Völker sollten kommen von den zwei Kindern, nämlich Israel und Edomäi, davon viel in den Propheten geschrieben ist. Da ist in dem Capitel beschrieben und angezeigt fast der Gang, den da gehabt haben die zwei Völker; wie jetzt möchte sein zwischen Deutschen und Walen. [Sie] haben sich wunderlich geschlagen; Edom hat sich lange feindlich gewehrt, doch zuletzt herunter lassen treten. Esau ward zeitlich ein großer Herr, nahm zeitlich das Land Edom ein. Und weil er so regierte, und mußte in seinem Lande, zeucht Jakob in Egypten, und blieb ein armer Fremdling, daß sein Volk hernach groß Herzeleid tragen mußte, und mehr denn dreihundert Jahre gehen, ehe es dahin kam, daß es Esau unter sich brächte, da Jakob, Moses, Aaron etc. alle todt waren, bis an David und Salomo, die brachten es erst unter sich. Da siehst du abermal, wie GOtt seine Verheißung vollführt, daß es scheint, als habe er die Seinen verworfen, und die er verwirft, bald empor hebt.

10. Das hat müssen bedeuten, daß der Text erstlich schreibt, wie Jsaak dunkel siehet, das ist, GOtt thut die Augen zu, und regieret, als kenne er Jakob nicht, daß er daselbig Volk so lange läßt gehen, und hiebei Esau empor hebt. Denn also steht sein Sinn: Esau ist mein Sohn, dem will ich den Segen geben. Darum muß Jsaak GOttes Figur sein, der thut, als sähe er nicht; das ist, wenn GOtt das Regiment führt mit seinen Heiligen, stellt er sich gleich, als wollte er der Frommen nicht wissen, hat allerdings die Geberde, als wollte er Esau segnen, und alles geben; noch kriegt es endlich der Jakob. Darum, ob GOtt wohl nicht blind ist, oder dunkel sieht, stellen sich doch seine Werke also; aber alles uns zu Trost. Laß ihn nur dunkel sehen, und thun, als sähe er uns nicht, er wird es dennoch recht treffen.

11. Also ist es endlich genommen, daß durch das Volk hernach Edom herunter geworfen ist, wiewohl sie es gar ungerne gethan haben; sind einander feind gewesen, sonderlich, da sie sahen, daß GOtt den Israel segnete. Da huben sie an zu schreien, wie hier der Text sagt, wie Esau so jämmerlich heult und klagt. Aber GOtt sagt also dazu: Ich kann ihm nicht thun: „Er hat den Segen dahin, und wird ihn auch behalten.“ Doch sagt er, du sollst auch ein Regiment haben, und dich deines Schwertes nähren, und es wird eine Zeit kommen, daß du sein Joch von deinem Halse reißest, das ist, du wirst ihm nicht ewig unterworfen sein. Davon lies die Historie im andern und vierten Buche der Könige [2 Sam. 8, 14. 2 Kön. 8, 20.], wie Edom abgesatelt ist und dem jüdischen Volke nicht mehr unterthänig wollte sein, und ist ein eigen Königreich worden.

12. Also werden wir in der Bibel alle Historien dahin sehen gehen, wie GOtt wahrhaftig, und doch sich stellt als unwahrhaftig, daß man klug sei, und ihn recht lerne kennen, wie er sein Thun ausführt über alle Sinne und Vernunft. Da wie sehr noth wäre es, daß man die Worte recht verstünde! Was tolle Geister werden noch vom Evangelio predigen, aber sich nichts darein wissen zu richten, wie GOttes Wahrheit zugeht! Er ist so klug, daß er dich so irre macht in dem Erfüllen, daß du nicht wissen kannst, woran du seiest. Darum mag man wohl mit Furcht und Zittern bitten, und nicht mit ihm trotzen. Gewiß
\quellseite{222}
ist er; will aber auch gefürchtet sein, daß du nicht stolzirest, sondern bittest, daß du könnest fest halten. Es sind ihrer wohl, die da meinen, sie haben GOtt gefangen; aber er ist zu flug, will den Hochmuth nicht leiden, seiner Züchtigung soll man sich allein trösten, und ihn läßig walten, wie er es machen wolle. Summa: Es wird es niemand lernen, es sei denn eine Erfahrung da. Als, wenn wir in Todesnöthen oder Fahr kommen, da lernt sich’s recht, sonst bleibt es nur bei den Worten. Wenn es zum Treffen kommen soll, so vergißt sich’s wohl, denn da sieht man es gar viel anders. Das ist noch nach der Historie hin geredet. Nun sollen wir es auch auf Christum deuten; wenn wir den nicht haben, so haben wir es nicht recht getroffen.

13. Der rechte Segen, davon wir oben [Cap. 22, 18.] gehört haben, trifft das Evangelium an, als wir so oft gesagt haben, was der Segen sei, Abraham verheißen in seinem Samen. Denn dieser Segen zwischen den zweien Brüdern, daß Jakob ein groß Volk und reich sollte werden, ist auch ein Segen; jenes aber ist der rechte Hauptsegen, daß wir gesegnet werden durch Christum, wider den Fluch, der über Adam gegangen ist, uns von demselben zu entledigen, also, daß dadurch verkündigt und gebracht wird eitel Gnade und Leben.

14. Hier müssen uns nun die zween Brüder vortragen, wie es zugeht, daß der Segen so seltsam und wunderlich gefällt,\footnote{gefallen = wohin kommen, fallen.} daß es kein Mensch glaubt. Da lasse ich aber bleiben: Isaak den Vater, und die zween Brüder zwei Völker; diese, die da wandeln in einem schönen, scheinenden Leben; die andern, die Aichenbrödel, die da im Glauben wandeln. So ist nun Esau der Mann, der da regiert, wie die Juden, so das Gesetz predigten, GOttes Volk,\footnote{So die Jenaer. Erlanger: „wie die Juden das Gesetz predigten, das Gottes-Volk.“ Wittenberger: „das sind die Juden, das Volk des Gesetzes.“} die in einem ehrbaren Leben gingen, und nicht anders meinten, denn sie säßen GOtt in dem Schooße, und war doch nichts mehr, denn allein äußerlicher Schmuck, und körperliche Ehre.

15. Darum sagt die Schrift wohl, daß Rebekka bei sich hatte des Esau köstliche Kleider, die sie Jakob anzog. Die Juden haben den Vortheil gehabt, daß sie die Schrift GOttes hatten, wie Paulus Röm. 3, 2. sagt. Die edlen, köstlichen Kleider, da die Wahrheit GOttes in liegt, die hatten sie im Hause bei sich, zogen sie aber nicht an, sondern Jakob that sie an. Das hat sich begeben, da Christus kam. Vor dem ging die Synagoge und das Gesetz in ihrem Schwange; aber als er kam, muß ein ander Volk daher, dem zog sie die Kleider an. Das war der Jakob, das ist der rechte Christus, der in der Schrift verwickelt ist.

16. Also bekleidet die rechte Rebekka, die Christenheit, ihren Sohn mit den Kleidern, das ist, mit den Sprüchen der Schrift, die von der Gnade und Christo sagen, nicht von den Werken. Daraus wird nun gegeben der rechte Segen, und wird genommen dem, der mit Werken umgehet; das ist, eben zu der Zeit, da sie auf das höchste das Gesetz trieben, fahren Christus und die Apostel zu, und gewinnen die Schrift, nehmen den Segen hinweg von GOtt, und er stellt sich doch also dazu, als sähe er es nicht, daß es scheint vor uns, als gehe es ohngefahr zu.

17. Dabei sagt die Schrift, wie Jsaak gesagt hat: „Die Stimme ist Jakobs Stimme, aber die Hände sind Esaus Hände.“ Denn weil Esau rauh war, dachte die Mutter: Wie thun wir, daß wir den Sohn auch rauh machen, daß [es] der Vater nicht wisse? Und „schlachtete zwei Böcklein, und that ihm die Felle um die Hände“, daß Jsaak nicht anders meinte, es wären Esaus Hände, und hörte doch Jakobs Stimme. Also hat ihn GOtt in dem Stücke der Sinne beraubt, daß es also mußte geschehen.

18. Was ist nun das? Die Stimme bedeutet nichts anders denn die Predigt; die Hände aber bedeuten das Leben und Werk, damit man betrogen wird. Einer predigt nicht wie der andere. Esau ist der Jäger, der treibt das Gesetz, und will die Leute mit Gewalt fromm machen; Jakob aber treibt nicht, sondern lockt mit der Gnade. Das sind zweierlei Stimmen, daß man eine wohl vor der\footnote{Erlanger: für die andern.} andern hört, doch führen sie beide gleiche Werke. Jakob hält auch das Gesetz, und thut dasjenige, das Esau lehrt, und lehrt doch anders. Nach dem äußerlichen Wandel ist es alles gleich; da heißt es also: Ist einer fromm, so sind sie es alle. Das Evangelium aber fragt nicht, wie die Werke sind.

19. Darum ist es ohn Gefahr,\footnote{Gefahr = Betrug, Hinterlist. (Dietz.)} daß die Christen halten alles, was andere Leute halten,
\quellseite{223}
ohne allein, daß die Stimme muß anders sein; da muß man scheiden das Gesetz und Evangelium. In das Gewissen soll man kein Gesetz treiben,\footnote{So die Jenaer. Wittenberger: „Dem Gewissen soll man kein Gesetz predigen.“ Erlanger: „Ins Gewissen soll man kein Gesetz predigen, in die Hände soll man kein Evangelium thun.“} den Händen soll man kein Evangelium predigen. Dem Gewissen soll man also predigen, daß man allein durch den Glauben an Christum müsse fromm und selig werden, und was man für Werke thue, daß man sie ja nicht thue, dem Gewissen zu helfen, sondern äußerlich dem alten Adam soll man auflegen alles, was das Gesetz fordert, also, daß das Gesetz außen, das Evangelium inwendig bleibe; so ist es recht, und bleibt die Stimme Jakobs Stimme, aber die Hände Esaus Hände. Das Gesetz thut nichts, denn zwingt und treibt immer, und hält die Leute mit Gewalt innen; aber das ist der Fehl daran, wenn man es will dem Gewissen predigen, das ist Esaus Stimme, die hört Gott nicht. Also redet Paulus davon, 1 Cor. 9, 27., und spricht: „Ich zähme meinen Leib, und betäube ihn, daß ich nicht den andern predige, und selbst verwerflich werde“; also, daß man alle Gesetze thue, und doch sage, es hilft dem Herzen und Gewissen nichts.

20. Wiederum wäre es auch nicht fein, daß man es umkehrte, und das Gesetz in die Gewissen triebe, und die Freiheit herausen ließe. Es muß Jakobs Stimme und Esaus Hände sein, nicht Jakobs Hände und Esaus Stimme. Wenn es eitel Jakobs Stimme und auch Jakobs Hände wären, taugte [es] auch nicht; wie die sind, die da sagen: Sollen die guten Werke nicht helfen, so wollen wir gute Gesellen sein, und nichts thun. Siehe, also ist in den zwei Brüdern abgemalt der Unterschied zwischen dem zweierlei Volk. Jene verlassen sich auf ihr Recht, meinen, sie sind das liebe Kind, und verlieren den Segen, und fahren zum Teufel; diese trösten sich allein Gottes Wortes, ergreifen den Segen, und kommen gen Himmel.

21. Danach hebt sich die Verfolgung, daß Esau dem Jakob so feind und gram wird, daß er ihn denkt zu erwürgen, daß ist auch noch der Werkheiligen und Esauiten Art, die nicht können leiden, daß ihr Thun und Werke nichts sollen gelten, heben an zu toben, und werden der rechten Christen tödliche Feinde. Darum ist es eine große Erkenntniß, wenn man das Evangelium recht erkennen soll, und recht damit fahren. Darum will Paulus [2 Tim. 2, 15.] haben, daß man das Wort Gottes recht schneide, nämlich, in die zwei Stücke, wie ich gesagt habe, daß man das Gewissen frei halte, und den Leib belade mit Gesetzen; und nicht wiederum, beide beschwert, oder beide frei gelassen. Die Seele soll leben im Geiste, aber der Leib soll sterben um der Sünde willen, sagt St. Paulus Röm. 8, 13. Das sei die geistliche Deutung dieses Capitels.

22. Daß nun jemand möchte bewegen, wie Jakob so untreulich mit seinem Bruder gehandelt habe, als Isaac selbst sagt: „Dein Bruder ist mit List gekommen, und hat dir den Segen hinterrückt“; item, daß auch Isaac den Segen wohl hätte mögen widerrufen, und auch närrisch genug ist, daß er also sollte betrogen werden, so er doch die Stimme Jakobs kennt, daß es sicher unglaublich lautet. Zum ersten habe ich zuvor mehr gesagt, daß wir nicht des Sinnes sind, die heiligen Väter gar rein zu machen, und können leiden, daß sie auch zuweilen gestrauchelt haben, und dann am meisten, wenn sie gegangen sind in sonderlichen Werken. Als hier: Der Jakob hat den rechten Segen für sich, und Gott wollte es also haben, daß er ihn überkommen sollte; so waren nun Jakob und Rebekka der Sache gewiß, daß sie wohl thaten. Ob sie aber Unrecht gethan haben, daß sie mit Hinterlist mit dem Vater umgingen, und Jakob ihm vorlog, und sprach: „Ich bin dein Sohn Esau“, lasse ich geschehen. Wenn das Hauptwerk gut ist, kann Gott wohl durch die Finger sehen, ob man ein wenig zu viel thut.

23. Also lesen wir ein Exempel von dem Könige David im andern Buche Samuelis [Cap. 16, 2. 3.], als er vor\footnote{So die Jenaer. Wittenberger und Erlanger: von.} seinem Sohne Absalom floh, und in der Wüste war, fromm und demüthig, kommt zu ihm ein Bube, der hieß Siba,\footnote{In der Bibel: Ziba.} vermacht und verleugnet seinen Herrn Mephiboseth, des Königes Saul Sohn, daß David zufährt unverhört, und verheißt ihm, in alle Güter seines Herrn zu setzen. Das war auch zu viel von dem heiligen Manne, daß er sich läßt überreden durch eine giftige Zunge, und mit Unrecht und Gewalt fährt.
\quellseite{224}

24. Das läßt GOtt geschehen, daß seine trefflichen Heiligen auch bei uns hernieden bleiben, und wir sehen, daß sie auch Fleisch und Blut sind; sonst würde folgen, daß man verzweifelte. Zum andern, daß man die Heiligen nicht\footnote{„nicht“ fehlt in der Wittenberger und in der Erlanger.} zu weit von uns scheidete, wie man doch gethan hat. Wenn sie haben können machen, daß sie keine Sünde hätten gehabt, so sollte es köstlich Ding sein. Damit wir sind trostlos worden, und haben verzweifeln müssen.

25. Darum wollte ich viel lieber, so viel als ich könnte, sofern es [es] nicht wider die Schrift wäre, die Heiligen beisammen bleiben. Denn es ist GOttes Gnade geschändet, und Christo zu nahe geredet, daß man sie so weit will von uns scheiden, als Himmel und Erde. Ich soll den nächsten, geringsten Christen nicht geringer achten denn St. Petrum und alle Heiligen im Himmel. Es hat keiner mehr Gnade; mehr Geschenke und Kleinod haben sie wohl; ich bin eben so theuer erkauft als sie, so sind sie eben in dem Fleische und Blute gesteckt, darinne ich stecke. Daraus können wir Hoffnung und Trost gewinnen; sonst machen wir Abgötter aus ihnen, und verachten die Gnade GOttes.

26. Ja, sagen sie, man muß ja die Heiligen ehren, und uns gering achten, und mit Demuth sich gegen ihnen stellen? Antwort: Hierinne gilt nicht Demuth, wenn es GOttes Gnade trifft, sondern Hochmuth. Auf dein Ding und Wesen sollst du demüthig sein, denn du bist nichts denn ein Madensack; aber auf Christi Güte kannst du nicht stolz genug sein, und mußt sagen: Wenn ich noch zehnmal so unfläthig wäre, habe ich dennoch das Blut, das mich rein und heilig macht, und Christo so viel gekostet hat, mich zu erlösen, als St. Petrum. Sie sind eben so tief gewesen als wir; so sind wir eben so hoch als sie, also daß keiner mehr hat, denn ich, ohne daß sie wohl stärker sind im Glauben gewesen; doch ist die Kost und Hauptgut nicht größer. So lerne erkennen, was wir an GOtt haben; denn er läßt uns solches allein zu Trost schreiben, und wir Narren machen es uns nur zum Schrecken. Sie, die heiligen Väter, dürfen es nicht, denn sie sind lange todt; wir dürfen es aber, daß wir daraus lernen GOtt erkennen; und eben darum müssen wir Sünder sein, daß ihr Leben unser Trost sei. Narr ich, so wollte es GOtt, [ich] stehe aber wieder auf. Abraham, Isaak, Jakob, David, Petrus und andere viel mehr haben auch genarret. Damit können wir Sünde und Teufel widerstehen.

27. Item, hier lernen wir auch, was Segen und Fluchen heißt. Denn wir haben gehört, wie der Segen also geht über Jakob, daß er soll Korn und Weins die Fülle haben, und Herr sein über seine Brüder, und über viel Land und Leute herrschen. Da siehst du noch kein geistlich Gut, ist noch alles leiblich. Darum heißt segnen eigentlich, etwas Gutes wünschen, und fluchen, etwas Uebels wünschen; wie er denn spricht: „GOtt gebe dir“ etc. Nun ist das ein Wunsch, der allein zeitlich Gut betrifft, Reichthum und herrlich Regiment und große Herrschaft, nicht allein über die Brüder, sondern auch über ander Volk. Aber wie lange hat sich’s hernach verzogen! Doch ist er sicher, daß der Segen wahr ist. Isaak hat es wohl gefühlt, daß der Geist redete, daß es wahr müßte werden; darum hat er es auch bestätigt, und nicht widerrufen. In solchen Sachen haben sie nicht gefahren nach Menschenkünsten, wie Petrus [2. Ep. 1, 21.] sagt: „Es ist noch nie keine Weissagung aus menschlichem Willen hervor gebracht, sondern die heiligen Menschen GOttes haben geredet, getrieben von dem Heiligen Geist“; sonst hätten sie es nicht geredet; GOtt hätte es auch nicht beschrieben lassen.

28. Solches ist nun erfüllt eine lange Zeit hernach, wie es die Historien ausweisen; aber nichtsdestoweniger ist es nicht vergessen, daß ihm auch viel Widerstand würde begegnen. Denn was aus GOtt kommt, kann unangefochten nicht bleiben, es muß durch das Feuer, wie es auch ist. Wenn dich GOtt hieße eine Maus fangen, oder einen Strohhalm aufheben, das doch ein närrisch Ding wäre, und die Welt würde es gewahr, daß er es geheißen hätte, würde sie dich’s mit Frieden nicht lassen thun. Wie gering er ein Ding heißt, so hängt sich der Teufel daran, will GOttes Wort und Werk nicht leiden, und richtet alles Unglück an.

29. Also auch hier hat GOtt dem Volke das köstliche Land zugesagt zu geben. Weil es nun GOtt gibt, hat der Teufel sich dawider aufgelehnt, und so viel Verfolgen und Fluchen erweckt; wie man liest in den Büchern der Könige, daß sich daran hingen die Assyrer, Egypter, Edomiter, Moabiter, Egypter, und andere Völker mehr, daß kein Volk auf Erden war, dem man
\quellseite{225}
also zugesetzt hätte von allen Orten; wie David selbst schreibt im 83. Psalm, V. 3—9.: „Siehe, deine Feinde toben, und die dich hassen, heben den Kopf auf. Sie trachten listiglich wider dein Volk, und rathschlagen wider deine Verbogenen. Sie sprechen: Kommt her, lasset uns sie zunichte machen, daß sie kein Volk seien, daß des Namens Israel nicht mehr gedacht werde; denn sie haben sich berathen in ihrem Herzen, und machen mit einander einen Bund wider dich, die Hütten der Edomiter [und Ismaeliter, der Moabiter und Hagariter, der Gebaliter, Ammoniter]\footnote{Die eingeklammerten Worte fehlen in den Ausgaben, wie wir vermuthen, aus Versehen.} und Amalekiter, der Philister sammt denen von Tyro. Assur hat sich auch zu ihnen geschlagen, und sind worden ein Arm der Kinder Lot.“

30. Da zählt der Prophet alle den großen Haufen, der sich wider sie erhob; aber sie hatten einen Schutzherrn, der sie wohl handhaben konnte; darum schlugen sie auch getrost um sich, daß er oft mit Einem Manne oder Engel ein ganz Heer hinweg schlug. Daß man in der Historie durch und durch sieht, wie der Segen wunderbarerlich erfüllt ist, daß mit dem Volke alle Welt zu schaffen hatte, und Gott sie oft ließ unterliegen, daß die Feinde sie meinten gar zu vertilgen; wie in dem genannten 83. Psalm, V. 5., steht: „Sie sprechen: Kommt her, lasset uns sie zunichte machen, daß sie kein Volk seien, daß des Namens Israel nicht mehr gedacht werde.“ Ja, sagt Gott dazu, ich habe einen Segen darüber gesprochen, daß sie wohl bleiben werden; wie man liest im vierten Buche Mosis [Cap. 23, 20.], wie Bileam dem Volke Israel sollte fluchen; aber da er das Maul aufthat, konnte er nichts denn segnen.

31. Also muß alles verfolgt, geschändet und gelästert werden, was Gottes Volk ist. Darum spricht auch Isaak: „Verflucht sei, wer dich verfluchet“, als wollte er sagen: Rüste dich nur darauf, du wirst leiden müssen, daß man dich verfolgt als kein Volk auf Erden. Aber den Trost sollst du dagegen haben, daß, wer dich verflucht, soll wieder verflucht sein, und soll ihm\footnote{„soll ihm“ fehlt in der Erlanger.} nimmer wohl gehen.

32. Das ist nun der leibliche Segen, bedeutet aber den rechten geistlichen Segen; darum ist es viel ein anderer Segen, denn der über Abraham gesprochen [Cap. 12, 3.]: „In deinem Samen sollen gesegnet werden alle Geschlechter auf Erden.“ Denn diesen Segen redet Gott durch einen Mann, und auf vergängliche Dinge; jenen aber redet er selbst, nicht von Wein oder Korn, sondern schlecht von der Person, über die der Segen soll gehen; darum muß er geistlich sein. Aus welchem weiter folgt: Wenn er über alle Heiden gehen, und gleich sein soll, so muß Abraham gleich sein allen Heiden; nicht wie zu Jakob gesagt wird, daß er Herr soll sein über seine Brüder und viel ander Volk. Weil nun dieses alles in diesem Segen gleich ist, muß es ein solches Ding sein, das sie alle begreifen können; welches ja nicht leiblich sein mag. Denn die Welt kann nicht bestehen also, daß alle gleich werden; es muß eines regieren, das andere unterthan sein; und der Segen will doch jedermann gleich haben, und über alle gehen.

33. Das ist das Evangelium, ein guter Wunsch, daß uns viel Gutes widerfahren soll. Denn also spricht das Evangelium über jedermann: Wir sollen die Sünden vergeben sein, Christus mit allen seinen Gütern und ewig Leben soll dein sein. In dem hohen Segen hat einer nicht mehr noch weniger, denn der andere, [sie] sind alle Brüder und gleiche Erben.

34. Mehr weiß ich von dem Capitel nicht viel zu sagen, denn daß Esau über den Bruder klagt und spricht: „Er heißt billig Jakob, denn er hat mich zweimal untergetreten.“ Das Wörtlein Eseb, habe ich gesagt, heißt auf deutsch eine Fußsohle; davon ist er genannt, darum, daß er des Bruders Fußsohle in der Hand hielt, da sie geboren wurden [1 Mos. 25, 26.]. Hier aber deutet es Esau ein wenig anders, will also sagen: Ja, er heißt wohl Fußsohle, denn er hat mich zweimal unter die Füße getreten. Was aber das Füßetreten bedeutet, haben wir vor gehört, nämlich, daß die Christen mit Füßen treten Sünde, Tod, Teufel und alles, was da geboren wird aus dem Fleisch.
\quellseite{226}

\section{Das achtundzwanzigste Capitel}

V. 1—9. Da rief Isaac seinen Sohn Jakob, und segnet ihn, und gebot ihm, und sprach zu ihm: Nimm nicht ein Weib von den Töchtern Canaan, sondern mache dich auf und zeuch zu Mesopotamien zu Bethuels, deiner Mutter Vaters, Haus, und nimm dir ein Weib daselbst von den Töchtern Labans, deiner Mutter Bruder. Aber der allmächtige GOtt segne dich, und mache dich fruchtbar, und mehre dich, daß du werdest ein Haufen Völker, und gebe dir den Segen Abrahams, dir und deinem Samen mit dir, daß du besitzest das Land, da du Fremdling innen bist, das GOtt Abraham gegeben hat. Also fertigte\footnote{fertigen = abfertigen, entsenden.} Isaac den Jakob, daß er in Mesopotamien zöge zu Laban, Bethuels Sohne, in\footnote{In den Ausgaben: von.} Syrien, dem Bruder Rebekka, seiner und Esaus Mutter. Als nun Esau sahe, daß Isaac Jakob gesegnet hatte, und abgefertigt in Mesopotamien, daß er daselbst ein Weib nähme, und daß, indem er ihn segnete, ihm gebot und sprach: Du sollst nicht ein Weib nehmen von den Töchtern Canaan, und daß Jakob seinem Vater und seiner Mutter gehorchte, und in Mesopotamien zog; sahe auch, daß Isaac, sein Vater, nicht gerne sahe die Töchter Canaan: ging er hin zu Ismael, und nahm über die Weiber, die er zuvor hatte, Mahalath, die Tochter Ismaels, des Sohnes Abrahams, die Schwester Nebajoths, zum Weibe.

1. Das wird die Lehre sein von Isaac und Rebekka, ohne daß er noch sterben wird, im fünf und dreißigsten Capitel, V. 28. 29. Zum ersten müssen wir es abfertigen mit dem Esau. Da siehst du, daß die Schrift den Jacob preist, wie er Vater und Mutter gehorcht hat.

2. Esau aber hatte zwei Weiber, die waren wunderlich und widerwertig, daß ihnen beiden übel gefiel, daß die Mutter auch klagte, und sprach [Cap. 27, 46.]: „Mich verdrießt zu leben vor den Töchtern Heth.“ So fährt er nun zu, als er sieht, daß Jakob nun in Gnaden ist, und hat den Segen hin, und zeucht zu seinem Vetter, der Mutter Bruder, allda ein Weib zu nehmen; stellt sich auch freundlich, und befreyundet sich auch nahe mit seinem Vetter Ismael, und nimmt desselben Tochter. Das ist nun im andern Gliede gefreiet unter Geschwister-Kindern, welches die Schrift nicht verboten hat. Solches hat er gethan, die Eltern zu versöhnen, und gefallen zu thun. Darauf folgt nun ein feiner Text von der Leiter, die Jakob im Traume sah, und hernach von dem Gelübde, das er that.

V. 10—15. Aber Jakob zog aus von Bersaba und reisete gen Haran, und kam an einen Ort, da blieb er über Nacht; denn die Sonne war untergegangen. Und er nahm einen Stein des Orts, und legte ihn zu seinen Häupten, und legte sich an demselben Ort schlafen. Und ihm träumte, und siehe, eine Leiter stand auf Erden, die rührete mit der Spitze an den Himmel. Und siehe, die Engel GOttes stiegen daran auf und nieder, und der HErr stand darauf, und sprach: Ich bin der HErr, Abrahams, deines Vaters, GOtt, und Isaacs GOtt, das Land, da du auf liegest, will ich dir und deinem Samen geben. Und dein Same soll werden, wie der Staub auf Erden, und du sollst ausgebreitet werden gegen dem Abend, Morgen, Mitternacht und Mittag, und durch dich und deinen Samen sollen alle Geschlechter auf Erden gesegnet werden. Und siehe, ich bin mit dir, und will dich behüten, wo du hinzeuchst, und will dich wieder her bringen in dies Land; denn ich will dich nicht verlassen, bis daß ich thue alles, was ich dir geredet habe.

3. Erstlich soll man die Historie immer ansehen nach dem Worte Gottes, denn man wird keine Legende mehr finden, da GOtt so viel mit Leuten geredet habe. Derhalben übertrifft sie alle Legenden. [Sie] ist auch darum geschrieben, daß man sähe, was rechte ausbündige Historien sind; nicht, die voll großer Wunder und köstlicher Werke sind, sondern nur, wo viel GOttes Wort innen ist. Vernunft sieht allein nach seltsamen, wunderlichen Geschichten und Thaten; da haben wir viel von gehalten; aber so klug sind wir nicht gewesen, daß wir fragten, ob auch GOttes Wort darinne wäre. Was ist es, daß einer alle die höchsten und größten Werke gethan hätte, und hätte kein GOttes Wort?
\quellseite{227}

4. Darum halte ich von diesen Historien, die haben Kraft, Saft und Mark, gehen immer in dem Worte, führen keine selbstersonnenen Werke, daß ihnen alle anderen das Wasser nicht reichen. Sollte es nicht ein Herz über alle Dinge erfreuen, wenn Gott also redet: „Ich will bei dir sein, und dich behüten, wo du hin zeuchst“ 2c.? Es sind eitel herzliche, lebendige Worte, welche anzeigen, daß, was er gethan,\footnote{Erlanger: hat.} hat Gott wohlgefallen, und hat eitel rechtschaffene, güldene Werke gethan. Nach solchen sollten wir auch trachten, daß wir kein Werk thäten, denn aus Gottes Wort, in welchem Stande es auch sei, daß mir gewiß sind, daß [es] Gott geboten hat und haben wolle.

5. Das ist nun das Hauptstück der Legende, daß Gott sein Wort so reichlich schlägt über das ganze Leben, daß der Mann muß heilig sein, von der Scheitel bis auf den Fuß. Ob er schon zuweilen strauchelt, soll es ihm nicht schaden, noch darum verdammt sein, so wenig als das Wort kann verdammt werden, welches niemand genugsam preisen kann.

6. So sagt nun Moses, wie Jakob ausgezogen ist von Bersaba, welches liegt in einem Winkel, hart an Egypten. [Es] ist der äußerste Winkel des gelobten Landes, und ist\footnote{In den Ausgaben: hat.} er gereiset gen Haram, zwischen Morgen und Mitternacht, und unterwegs auf der Straße ist er gekommen an einen Ort, der hat geheißen Luz, welches darnach ist worden Bethel, da Salomo seinen Tempel bauete. Denn die Städte hat Gott sonderlich geehrt, wie wir droben [Cap. 22, 9.] gehört haben, daß Isaak auf demselben Berge geopfert ist. Da ist nun Jakob über Nacht geblieben und [hat] geschlafen. Denn da er hingekommen ist, war die Sonne untergegangen, sagt Moses, also, daß es fast eine Tagesreise gewesen ist von seinem Auszuge. Da hat er sich niedergelegt, und einen Traum und Gesicht gesehen, wie der Text erzählt, und an ihm selbst klar getrieben ist. Was es aber bedeutet, werden wir hernach hören, wollen jetzt bei der Historie bleiben.

7. Das Wort, das der Herr da redet, und ihn segnet, gibt wohl zu verstehen, wie es um Jakob gestanden ist. Denn Gott ist nicht ein unnützer Wächter, daß er sein Wort in Wind werfe, und rede, wo es nicht noth ist. Darum weist sich's aus, daß Jakob in großem Leid und Angst gewesen sei, hat sich müssen verstohlen hinweg machen, weil ihm sein Bruder den Tod gebräuet hatte, also, daß er in Fahr des Lebens gestanden ist. Denn sein Bruder war auch befreundet im Lande, dadurch er gehen mußte. Wie bald wäre es geschehen, daß sie ihn erwürgt hätten! Weil er nun in solcher Fahr\footnote{Wittenberger und Erlanger: solche Fahr.} steht, und sich nicht zu schützen weiß, ist Gott da, und sagt ihm zu: weil ihm niemand helfe, so wolle er ihm helfen und schützen. [Er] schickt es auch eben darum also, daß er ohne Hülfe und Beistand muß dahin gehen, auf daß er es selbst ohne Mittel thue. Als sollte er sagen: Sei du getrost, und lasse deinen Bruder mit der Freundschaft auf dich stellen; ich will noch wohl eine Leiter herab bauen, und bei dir sein, wenn du da liegst und schläfst, und gar nichts dafür sorgst; ich habe mehr Engel im Himmel, denn Menschen auf Erden. Wenn die bei uns sind, wer will uns [etwas] thun?

8. Und hieher wird der Spruch geschlossen sein im 91. Psalm, V. 11. 12.: „Er hat seinen Engeln befohlen über dir, daß sie dich behüten auf allen deinen Wegen, daß sie dich auf den Händen tragen, daß du deinen Fuß nicht an einen Stein stoßest.“ Denn die lieben Väter haben aus der Historie viel gelernt. Also liest man auch [2 Kön. 6, 14—20.] von dem Propheten Helisäus, wie die Syrer eine Stadt umgaben, darinne er war, und wollten ihn holen. Aber als Helisäus früh aufstand mit seinem Knaben, und sahen, daß die Stadt belegt war, erschrak der Knabe, und sprach: „Awe, wie wollen wir nun thun?“ [Es] antwortete Helisäus: „Fürchte dich nicht, denn derer ist mehr, die bei uns sind, denn derer, die bei ihnen sind.“ Da wurden dem Knaben die Augen geöffnet, und [er] sahe, daß der Berg voll feuriger Rosse und Wagen war, um Helisäum den Propheten her. Darnach machte er das Volk blind, daß sie sich ließen von der Stadt führen, bis mitten in die Hauptstadt Samariam, und öffnete ihnen die Augen wieder; da sahen sie, wie sie mitten in der Stadt unter den Feinden waren.

9. Solches ist alles darum geschrieben, auf daß wir gewiß sein, daß Gott die Seinen nicht lassen kann; wenn auch alle Welt wider uns ist,
\quellseite{228}
so ist er da und hilft. Derhalben, wenn wir im Glauben sind, so stehen so viel Engel auf uns; und [er] stellt sich doch so närrisch dazu, daß uns dünkt, wir sind verloren, und müssen mit Schanden sterben, so stehen die Engel, und stehen uns zu, die uns wohl in einem Augenblicke könnten aushelfen. Gott sieht auch eine Zeitlang zu.

10. Warum läßt er uns denn nicht so bald helfen? Darum, daß seine Gnade, Werk und Wort desto baß erkannt werden, und beweise, daß er im Tode helfen kann, welches wir sonst nicht wüßten, wo wir es nicht erfahren. Wer nun ein Christ ist, der hat gut sterben; denn also: Wenn die Seele ausfährt, so steht es voll, voll Engel umher, die sie Gott in den Schooß hinein tragen. Aber solch tröstlich Ding widerfährt niemand, ehe denn er sieht, wie er von allen Creaturen verlassen ist, und niemand helfen kann, denn Gott allein.

11. Gewißlich ist es wahr, wenn wir es nur könnten glauben. Also, da dies arme Kind in höchsten Nöthen liegt, und seines Lebens nicht sicher ist, muß er so großen Trost und Hülfe sehen. Gott gab ihm solchen Trost nicht, da er bei Vater und Mutter war, und sich nichts besorgen durfte; sondern mußte vor von ihnen kommen, und den zehnten Vers aus dem 27. Psalm singen: „Mein Vater und Mutter haben mich verlassen, aber der Herr hat mich aufgenommen“; daß er sagen könnte: Ich will lieber Gott im Schooße sitzen denn Vater und Mutter. Darum muß er von ihnen kommen unter die Feinde; da sie sitzen und regieren, da schützt ihn Gott auf das höchste durch seine Engel.

B. 16—19. Da nun Jacob von seinem Schlafe aufwachte, sprach er: Gewißlich ist der Herr an diesem Orte, und ich wußte es nicht. Und fürchtete sich, und sprach: Wie schrecklich ist diese Stätte! Es ist hier nichts anders, denn ein Haus Gottes, und ein Thor zum Himmel. Und Jacob stund des Morgens frühe auf, und nahm den Stein, den er zu seinen Häupten gelegt hatte, und richtete ihn auf, und goß Oel oben darauf, und hieß die Stätte Bethel, vorher hieß sonst die Stadt Lus.

12. Hier müssen wir einmal von Kirchen predigen. Denn unsere Papisten haben diesen Text für sich gezogen auf ihre Kirchen, weil Jacob die Stätte nennt ein Gotteshaus. Es war aber noch keines, sondern er richtete einen Stein zum Male auf, und begob ihn mit Oele; item, thut ein Gelübde dazu, daß da sollte ein Gotteshaus werden, und will von allem, das er kriegte, den Zehnten dazu geben. Da siehe zum ersten, daß Gottes Haus nicht heißt ein köstlich groß Gebäu, wie wir haben. Denn so redet auch David [Ps. 5, 8.]: „Ich will in dein Haus gehen auf deine große Güte, und anbeten gegen deinem heiligen Tempel“; und war doch noch kein Tempel gebauet, sondern [sie] hatten allein den Tabernakel, oder die Hütte Mosis.

13. Also sagt hier Jacob: Der Stein solle ein Gotteshaus werden. Warum nennt er es denn also? Nicht darum, daß vonnöthen wäre, ein groß Gebäude dazu aufzurichten, sondern [hat] allein ein Zeichen und Mal gestiftet, daß Gott allda wohne; will also sagen: Hier soll Gottes Wohnung sein; darum nennt er es hebräisch BethEl. Nun fragt Gott nicht darnach, ob es groß oder nicht gebauet, gewölbt und gemehlt sei; ja, er wohnt da, und bauet ihm doch sein Lebetage kein Haus. Was gehört denn dazu, daß Gott allda wohnt? Nichts mehr, denn daß Gott da sei mit seinem Worte; wo das geht, da wohnt er gewißlich, und wiederum, wo das Wort nicht ist, da wohnt er nicht, man baue ihm ein Haus, so groß man wolle.

14. Das lerne und merke, zu antworten denen, die von Gotteshäusern und Kirchen rühmen, und den Spruch dahin dehnen. Es heißt nicht Gottes Haus um deines Dienstes oder Stiftens willen; um deines Werkes willen will er nicht da wohnen. Denn dein Werk soll so frei sein, daß es an keinen Ort gebunden ist, es wäre denn, daß dir Gott etwas Sonderliches befohlen hätte. Aber darum soll es sein Haus heißen, wenn er kommt und uns dient, und läßt sein Wort da ausgehen.

15. Das ist nun aus dem Texte klar, als Jakob selbst spricht, da er im Gesichte Gottes Wort gehört hatte: „Hier ist nichts anders, denn Gottes Haus, und ein Thor zum Himmel.“ Ist es schon allda, warum will er es denn da aufrichten? Darum ist es da, daß Gott daselbst wohnt mit den Engeln, und sein Wort hören läßt. Deß nimmt ein Gleichniß. Das heißt Meister Hansens Haus; darum, daß der Knecht darinne arbeitet? Nein, sondern daß er darinne wohnt und Wirth ist, daß es alles muß gehen, wie er will; nicht davon, daß der Knecht darinne dient. Also auch hier, wo man hinkommt, da
\quellseite{229}
Gott regiert und schafft durch sein Wort, da ist sein Haus, daß du sagen magst: Hier ist eigentlich Gottes Wohnung, eine Leiter und Thor zum Himmel. Denn da steht uns der Himmel offen, daher wir das Wort haben, welches uns das Leben und den Himmel gibt.

16. Demnach kannst du urtheilen, und schließen, daß unsere Klöster und Stiftkirchen nicht Gottes Häuser sind, denn da ist es alles also gestiftet, daß man kein Gottes Wort darinne hat. Denn wo Gott wohnt, da schweiget er nicht stille, und wo er redet, da wohnt er auch. Darum soll man kein Haus Gottes nennen, man sei denn gewiß, daß er da rede. Denn ehe Jakob Willen hatte, oder je gedachte, etwas da zu machen, hört er Gott reden und predigen, und um des Wortes willen sagt er: „Hier wohnet Gott.“ Derhalben müssen unsere Kirchen und Stifte eitel Teufelshäuser sein, weil sie ohne Gottes Wort, aus menschlichem Gutdünken aufgerichtet sind, und allein von dem Werke und Dienste, den wir gestiftet haben, Gottes Häuser heißen.

17. O was hat der Text den Juden auch Irrthum gegeben, gerade wie uns, daß sie frei dahin schließen: Hier hat Jakob gelegen und Gott gesehen, [das] ist eine heilige Stätte, darum wollen wir da eine Kirche bauen, und Gottesdienst aufrichten; wie der König Jerobeam that [1 Kön. 12, 31.]. Dawider schrieen alle Propheten und sagten: Was hebt ihr an? Redet doch Gott nicht da, und ihr richtet es ohne und wider seinen Befehl auf. Da legten sich die falschen Propheten dagegen und verdammten sie, und richteten so viel zu, daß sie mußten um das Leben kommen. Was hat sie nun geirrt?\footnote{geirrt = irre gemacht.} Eben das, daß sie den Text nicht recht angesehen haben, und das Maul aufgesperrt gegen der Stätte, und meinten, es müßte Gott da wohnen, darum, daß sie ihm darinne dienten, und nicht nach dem Worte frageten, ob er da redete.

18. Du darfst nicht denken auf herrliche Gebäude, Dach, Gewölbe, Altar, Singen und Klingen. Wo Gottes Wort klingt, es sei im Walde oder Wasser, oder wo es ist, da ist ein BethEl, daß man darf sagen: Hier wohnt Gott. Denn, wie gehört, ehe je ein Stein an dem Orte gelegt war, hieß Jakob den freien Platz ein Haus Gottes. Woher wußte er das? Daher, daß er da an dem Orte Gott hat hören reden. Darum soll man ihm kein Haus bauen, man müsse denn zuvor, daß er da wohne. Dann wohnt er aber da, wenn er sein Wort läßt predigen, in uns wirkt, und durch den Glauben erkannt wird.

V. 20—22. Und Jakob that ein Gelübde und sprach: So Gott wird mit mir sein, und mich behüten auf dem Wege, den ich reise, und Brod zu essen geben und Kleider anzuziehen, und mich mit Frieden wieder heim zu meinem Vater bringen, so soll der Herr mein Gott sein,\footnote{Randglosse: Mein Gott sein. Nicht, daß er vorhin nicht sein Gott gewesen sei, sondern er gelobt einen Gottesdienst aufzurichten, da man predigen und beten sollte, da will er den Zehnten zu geben den Predigern, wie Abraham Melchisedek den Zehnten gab.} und dieser Stein, den ich aufgerichtet habe, soll ein Gotteshaus werden, und alles, was du mir gibst, deß will ich dir den Zehnten geben.

\footnote{Hier haben die Wittenberger und die Jenaer Ausgabe die Ueberschrift: „Von den Gelübden.“} 19. Ich habe bisher genug gepredigt und geschrieben, wie alle unsere papistischen Klostergelübde wider Gottes Wort und verdammlich sind. Wenn nun die Papisten auftreten und diesen Text anzeigen wider uns, und darauf drängen, daß da ein Exempel steht des heiligen Vaters, der Gott ein Gelübde gethan hat, welches von Gott gelobt ist, und hat es auch gehalten, und du bist ein Mönch gewesen, hast so viel gelobt und hältst nichts: was willst du darauf antworten? Willst du sagen: Ich habe es gehört oder gelesen, es gelte nichts, und sei nicht christlich; so wirst du nicht bestehen. Du mußt den Teufel wissen zu schlagen, das Schwert zu nehmen, und dich wehren. Darum lasset uns sehen, was dieser Text schließet oder nicht, wie er bringt, und nicht bringt.

20. Zum ersten: Die Gelübde, die jetzt zu Boten geben,\footnote{Walch und die Erlanger: „zu Zeiten geben“. „Zu Boten geben“ heißt: haften, einwürgen. Vgl. Walch, St. Louiser Ausg., Bd. VIII, 49, § 109.} haben gemeiniglich drei Fehler, welche alle wider diesen Text sind. Das erste ist, daß die Gelübde, so man thun will in geistlichen Sachen, sollen an keinen Heiligen gerathen, sondern an Gott allein. Das fasse du wohl. Also spricht hier Jakob: „Der Herr soll mein Gott sein, und dieser Stein soll ein Gotteshaus werden“; nicht, daß er gelobe Abraham, Isaac, Noah, oder irgend einem Heiligen.
